% Options for packages loaded elsewhere
\PassOptionsToPackage{unicode}{hyperref}
\PassOptionsToPackage{hyphens}{url}
%
\documentclass[
]{article}
\usepackage{amsmath,amssymb}
\usepackage{iftex}
\ifPDFTeX
  \usepackage[T1]{fontenc}
  \usepackage[utf8]{inputenc}
  \usepackage{textcomp} % provide euro and other symbols
\else % if luatex or xetex
  \usepackage{unicode-math} % this also loads fontspec
  \defaultfontfeatures{Scale=MatchLowercase}
  \defaultfontfeatures[\rmfamily]{Ligatures=TeX,Scale=1}
\fi
\usepackage{lmodern}
\ifPDFTeX\else
  % xetex/luatex font selection
\fi
% Use upquote if available, for straight quotes in verbatim environments
\IfFileExists{upquote.sty}{\usepackage{upquote}}{}
\IfFileExists{microtype.sty}{% use microtype if available
  \usepackage[]{microtype}
  \UseMicrotypeSet[protrusion]{basicmath} % disable protrusion for tt fonts
}{}
\makeatletter
\@ifundefined{KOMAClassName}{% if non-KOMA class
  \IfFileExists{parskip.sty}{%
    \usepackage{parskip}
  }{% else
    \setlength{\parindent}{0pt}
    \setlength{\parskip}{6pt plus 2pt minus 1pt}}
}{% if KOMA class
  \KOMAoptions{parskip=half}}
\makeatother
\usepackage{xcolor}
\usepackage{color}
\usepackage{fancyvrb}
\newcommand{\VerbBar}{|}
\newcommand{\VERB}{\Verb[commandchars=\\\{\}]}
\DefineVerbatimEnvironment{Highlighting}{Verbatim}{commandchars=\\\{\}}
% Add ',fontsize=\small' for more characters per line
\newenvironment{Shaded}{}{}
\newcommand{\AlertTok}[1]{\textcolor[rgb]{1.00,0.00,0.00}{\textbf{#1}}}
\newcommand{\AnnotationTok}[1]{\textcolor[rgb]{0.38,0.63,0.69}{\textbf{\textit{#1}}}}
\newcommand{\AttributeTok}[1]{\textcolor[rgb]{0.49,0.56,0.16}{#1}}
\newcommand{\BaseNTok}[1]{\textcolor[rgb]{0.25,0.63,0.44}{#1}}
\newcommand{\BuiltInTok}[1]{\textcolor[rgb]{0.00,0.50,0.00}{#1}}
\newcommand{\CharTok}[1]{\textcolor[rgb]{0.25,0.44,0.63}{#1}}
\newcommand{\CommentTok}[1]{\textcolor[rgb]{0.38,0.63,0.69}{\textit{#1}}}
\newcommand{\CommentVarTok}[1]{\textcolor[rgb]{0.38,0.63,0.69}{\textbf{\textit{#1}}}}
\newcommand{\ConstantTok}[1]{\textcolor[rgb]{0.53,0.00,0.00}{#1}}
\newcommand{\ControlFlowTok}[1]{\textcolor[rgb]{0.00,0.44,0.13}{\textbf{#1}}}
\newcommand{\DataTypeTok}[1]{\textcolor[rgb]{0.56,0.13,0.00}{#1}}
\newcommand{\DecValTok}[1]{\textcolor[rgb]{0.25,0.63,0.44}{#1}}
\newcommand{\DocumentationTok}[1]{\textcolor[rgb]{0.73,0.13,0.13}{\textit{#1}}}
\newcommand{\ErrorTok}[1]{\textcolor[rgb]{1.00,0.00,0.00}{\textbf{#1}}}
\newcommand{\ExtensionTok}[1]{#1}
\newcommand{\FloatTok}[1]{\textcolor[rgb]{0.25,0.63,0.44}{#1}}
\newcommand{\FunctionTok}[1]{\textcolor[rgb]{0.02,0.16,0.49}{#1}}
\newcommand{\ImportTok}[1]{\textcolor[rgb]{0.00,0.50,0.00}{\textbf{#1}}}
\newcommand{\InformationTok}[1]{\textcolor[rgb]{0.38,0.63,0.69}{\textbf{\textit{#1}}}}
\newcommand{\KeywordTok}[1]{\textcolor[rgb]{0.00,0.44,0.13}{\textbf{#1}}}
\newcommand{\NormalTok}[1]{#1}
\newcommand{\OperatorTok}[1]{\textcolor[rgb]{0.40,0.40,0.40}{#1}}
\newcommand{\OtherTok}[1]{\textcolor[rgb]{0.00,0.44,0.13}{#1}}
\newcommand{\PreprocessorTok}[1]{\textcolor[rgb]{0.74,0.48,0.00}{#1}}
\newcommand{\RegionMarkerTok}[1]{#1}
\newcommand{\SpecialCharTok}[1]{\textcolor[rgb]{0.25,0.44,0.63}{#1}}
\newcommand{\SpecialStringTok}[1]{\textcolor[rgb]{0.73,0.40,0.53}{#1}}
\newcommand{\StringTok}[1]{\textcolor[rgb]{0.25,0.44,0.63}{#1}}
\newcommand{\VariableTok}[1]{\textcolor[rgb]{0.10,0.09,0.49}{#1}}
\newcommand{\VerbatimStringTok}[1]{\textcolor[rgb]{0.25,0.44,0.63}{#1}}
\newcommand{\WarningTok}[1]{\textcolor[rgb]{0.38,0.63,0.69}{\textbf{\textit{#1}}}}
\usepackage{longtable,booktabs,array}
\usepackage{calc} % for calculating minipage widths
% Correct order of tables after \paragraph or \subparagraph
\usepackage{etoolbox}
\makeatletter
\patchcmd\longtable{\par}{\if@noskipsec\mbox{}\fi\par}{}{}
\makeatother
% Allow footnotes in longtable head/foot
\IfFileExists{footnotehyper.sty}{\usepackage{footnotehyper}}{\usepackage{footnote}}
\makesavenoteenv{longtable}
\setlength{\emergencystretch}{3em} % prevent overfull lines
\providecommand{\tightlist}{%
  \setlength{\itemsep}{0pt}\setlength{\parskip}{0pt}}
\setcounter{secnumdepth}{-\maxdimen} % remove section numbering
\newlength{\cslhangindent}
\setlength{\cslhangindent}{1.5em}
\newlength{\csllabelwidth}
\setlength{\csllabelwidth}{3em}
\newlength{\cslentryspacingunit} % times entry-spacing
\setlength{\cslentryspacingunit}{\parskip}
\newenvironment{CSLReferences}[2] % #1 hanging-ident, #2 entry spacing
 {% don't indent paragraphs
  \setlength{\parindent}{0pt}
  % turn on hanging indent if param 1 is 1
  \ifodd #1
  \let\oldpar\par
  \def\par{\hangindent=\cslhangindent\oldpar}
  \fi
  % set entry spacing
  \setlength{\parskip}{#2\cslentryspacingunit}
 }%
 {}
\usepackage{calc}
\newcommand{\CSLBlock}[1]{#1\hfill\break}
\newcommand{\CSLLeftMargin}[1]{\parbox[t]{\csllabelwidth}{#1}}
\newcommand{\CSLRightInline}[1]{\parbox[t]{\linewidth - \csllabelwidth}{#1}\break}
\newcommand{\CSLIndent}[1]{\hspace{\cslhangindent}#1}
\ifLuaTeX
  \usepackage{selnolig}  % disable illegal ligatures
\fi
\IfFileExists{bookmark.sty}{\usepackage{bookmark}}{\usepackage{hyperref}}
\IfFileExists{xurl.sty}{\usepackage{xurl}}{} % add URL line breaks if available
\urlstyle{same}
\hypersetup{
  pdftitle={synapticTrack Consolidated Project Document},
  pdfauthor={Chong Shik Park},
  hidelinks,
  pdfcreator={LaTeX via pandoc}}

\title{synapticTrack Consolidated Project Document}
\author{Chong Shik Park}
\date{2026-08-18}

\begin{document}
\maketitle

\hypertarget{abstract}{%
\section{Abstract}\label{abstract}}

This document consolidates the Markdown documentation in the
\texttt{synapticTrack} repository into a single paper-oriented source.
The project author is Chong Shik Park, Department of Accelerator Science
and Center for Accelerator Research, Korea University, Sejong, 30019
Republic of Korea. It is organized to support later manuscript writing:
project motivation, literature context, physics models, simulation code
structure, simulation procedures, analysis methods, analysis procedures,
installation, and operating instructions.

\hypertarget{literature-review-and-technical-context}{%
\section{Literature Review and Technical
Context}\label{literature-review-and-technical-context}}

\texttt{synapticTrack} sits at the intersection of charged-particle beam
transport, surrogate modeling, differentiable inference, and
reinforcement learning for accelerator controls. Linear transverse beam
descriptions follow the standard Courant-Snyder/Twiss formalism for
second moments and emittance
(\protect\hyperlink{ref-courant1958theory}{Courant and Snyder 1958};
\protect\hyperlink{ref-wiedemann2015particle}{Wiedemann 2015}).
Production sampling uses Sobol low-discrepancy sequences to cover
high-dimensional control and beam-state spaces more evenly than
independent pseudo-random draws
(\protect\hyperlink{ref-sobol1967distribution}{Sobol' 1967}). The
surrogate and inference layers are implemented with PyTorch, which
provides autograd and tensor execution across CPU/GPU devices
(\protect\hyperlink{ref-paszke2019pytorch}{Paszke et al. 2019}). The
project also uses Python scientific infrastructure for arrays, tabular
data, optimization, and machine-learning evaluation
(\protect\hyperlink{ref-harris2020array}{Harris et al. 2020};
\protect\hyperlink{ref-virtanen2020scipy}{Virtanen et al. 2020};
\protect\hyperlink{ref-pandas2020}{team 2020};
\protect\hyperlink{ref-pedregosa2011scikit}{Pedregosa et al. 2011}).
Planned orbit-correction environments use the Gymnasium API and
reinforcement-learning baselines from Stable-Baselines3
(\protect\hyperlink{ref-towers2024gymnasium}{Towers et al. 2024};
\protect\hyperlink{ref-raffin2021stable}{Raffin et al. 2021}). External
high-fidelity comparison is anchored by accelerator tracking codes such
as TRACK (\protect\hyperlink{ref-ostroumov2006track}{Ostroumov and Aseev
2006}).

The practical objective is to build a staged path from deterministic
LEBT dataset generation through multi-output neural surrogates,
physics-regularized covariance outputs, beam-state inference, MEBT
surrogate modeling, and eventually RL orbit correction and
higher-fidelity reconstruction. Current maintained outputs include RAON
LEBT and MEBT notebook workspaces, staged no-offset Sobol production and
training notebooks, versioned training artifact manifests, read-only
HDF5 compatibility views, publication-ready metric tables, LINAC 2026
paper assets, and a generated static documentation site without
\texttt{polyfill.io} dependencies.

\hypertarget{project-overview}{%
\section{Project Overview}\label{project-overview}}

Source Markdown: \texttt{README.md}

\hypertarget{synaptictrack}{%
\subsection{synapticTrack}\label{synaptictrack}}

\textbf{synapticTrack} is a modular, extensible Python framework for
accelerator optimization using \textbf{machine learning} techniques.

\hypertarget{purpose}{%
\paragraph{Purpose}\label{purpose}}

The framework is currently under active development to support
\textbf{orbit correction and optimization} in \textbf{heavy-ion linear
accelerators}, specifically focused on the \textbf{LEBT (Low Energy Beam
Transport)} section.

\hypertarget{features-planned-in-progress}{%
\paragraph{Features (Planned \& In
Progress)}\label{features-planned-in-progress}}

\begin{itemize}
\tightlist
\item
  ML-based orbit correction using simulation and diagnostic data
\item
  Interoperability with multiple beam dynamics codes:

  \begin{itemize}
  \tightlist
  \item
    \href{https://www.phy.anl.gov/atlas/TRACK/}{TRACK (Fortran)}
  \item
    \href{https://github.com/MSU-Beam-Dynamics/JuTrack.jl}{JuTrack
    (Julia)}
  \end{itemize}
\item
  Support for:

  \begin{itemize}
  \tightlist
  \item
    Multi-ion species (different A/Q)
  \item
    Variable charge states and currents
  \item
    Simulation-to-reality generalization
  \end{itemize}
\end{itemize}

\hypertarget{repository-structure}{%
\paragraph{Repository Structure}\label{repository-structure}}

The repository now contains the package code, RAON lattice/notebook
workspaces, example workflows, paper assets, and generated model-output
areas. Generated run products are intentionally kept separate from
reusable source modules.

\begin{Shaded}
\begin{Highlighting}[]
\NormalTok{synapticTrack/}
\NormalTok{+{-}{-} .github/workflows/          \# GitHub Actions CI definitions}
\NormalTok{+{-}{-} RAON/                       \# RAON{-}specific lattice files and notebooks}
\NormalTok{|   +{-}{-} LEBT\_lattice/           \# LEBT TRACK lattice/input files}
\NormalTok{|   +{-}{-} LEBT\_notebooks/         \# LEBT calibration, orbit{-}correction, and extended{-}model notebooks}
\NormalTok{|   +{-}{-} MEBT\_lattice/           \# MEBT TRACK lattice/input files}
\NormalTok{|   +{-}{-} MEBT\_notebooks/         \# MEBT lattice creation/loading and segmented TRACK run notebooks}
\NormalTok{+{-}{-} TRACK/                      \# Local TRACK{-}related runtime assets, when present}
\NormalTok{+{-}{-} bin/                        \# Command{-}line helper scripts}
\NormalTok{+{-}{-} configs/                    \# Dataset, hardware, LEBT, local, and ML configuration files}
\NormalTok{+{-}{-} data/                       \# Processed/shared data products}
\NormalTok{+{-}{-} docs/                       \# Sphinx docs, task notes, paper material, and conference sources}
\NormalTok{|   +{-}{-} codex\_tasks/            \# Implementation task specifications and audit notes}
\NormalTok{|   +{-}{-} linac2026/              \# LINAC 2026 JACoW paper source, bibliography, and figures}
\NormalTok{|   +{-}{-} paper/                  \# Paper/report helper utilities}
\NormalTok{|   +{-}{-} site/                   \# Documentation/site helper utilities}
\NormalTok{+{-}{-} examples/                   \# Worked examples and legacy/reference notebooks}
\NormalTok{|   +{-}{-} analysis/               \# Scanner and Twiss{-}analysis examples}
\NormalTok{|   +{-}{-} beam/                   \# Beam object examples}
\NormalTok{|   +{-}{-} jutrack/                \# JuTrack integration examples}
\NormalTok{|   +{-}{-} lebt\_orbit\_corrections/ \# LEBT calibration, surrogate, and RL orbit{-}correction examples}
\NormalTok{|   +{-}{-} visualization/          \# Plotting examples}
\NormalTok{+{-}{-} models/                     \# Generated training outputs, summaries, comparisons, and metrics}
\NormalTok{+{-}{-} notebooks/                  \# Main staged LEBT production and model{-}training notebooks}
\NormalTok{+{-}{-} run.production/             \# Large staged production simulation outputs}
\NormalTok{+{-}{-} scripts/                    \# Repository maintenance and notebook hygiene scripts}
\NormalTok{+{-}{-} src/synapticTrack/          \# Installable Python package}
\NormalTok{|   +{-}{-} analysis/               \# Scanner and diagnostic analysis helpers}
\NormalTok{|   +{-}{-} beam/                   \# Beam, Gaussian, scanner, DataFrame, and Twiss utilities}
\NormalTok{|   +{-}{-} cli\_commands/           \# CLI preflight, TRACK checks, and validation commands}
\NormalTok{|   +{-}{-} data/                   \# Dataset schema, input/target views, batch I/O, and audits}
\NormalTok{|   +{-}{-} hardware/               \# Steering calibration and actuator abstractions}
\NormalTok{|   +{-}{-} inference/              \# Beam{-}state inference interfaces, losses, objectives, and optimizers}
\NormalTok{|   +{-}{-} io/                     \# TRACK, JuTrack, FLAME, OPAL, scanner, and beam{-}data adapters}
\NormalTok{|   +{-}{-} lattice/                \# TRACK lattice elements, conversion helpers, and JSON schema validation}
\NormalTok{|   +{-}{-} ml/                     \# Surrogate models, training, metrics, physics losses, and publication tables}
\NormalTok{|   +{-}{-} sampling/               \# Latin{-}hypercube and Sobol sampling utilities}
\NormalTok{|   +{-}{-} simulation/             \# Production dataset generation and TRACK execution orchestration}
\NormalTok{|   +{-}{-} track/                  \# TRACK executable wrapper}
\NormalTok{|   +{-}{-} utils/                  \# Calibration, RL, TRACK runtime, diagnostics, and parallel helpers}
\NormalTok{|   +{-}{-} visualizations/         \# Scanner, phase{-}space, and plotting utilities}
\NormalTok{+{-}{-} tests/                      \# Unit, integration, notebook hygiene, and physics/ML regression tests}
\end{Highlighting}
\end{Shaded}

\hypertarget{current-maintained-workflows}{%
\paragraph{Current Maintained
Workflows}\label{current-maintained-workflows}}

\begin{itemize}
\tightlist
\item
  \texttt{RAON/LEBT\_notebooks/} mirrors the maintained LEBT
  calibration, orbit-correction, and extended surrogate notebooks.
\item
  \texttt{RAON/MEBT\_notebooks/} contains MEBT lattice creation/loading
  and segmented TRACK run notebooks. Segmented TRACK runs advance
  scratch distribution numbers deterministically so each section reads
  the previous section output.
\item
  \texttt{notebooks/} contains the main staged LEBT no-offset Sobol
  production, validation, MLP training, physics-regularized training,
  model-comparison, and publication-metric notebooks.
\item
  \texttt{models/lebt\_no\_offset\_model\_comparison/publication\_metrics/}
  stores publication-oriented diagnostic, component-wise state, derived
  Courant-Snyder/emittance, and covariance-validity tables.
\item
  Per-run training artifacts use \texttt{training\_run\_artifacts\_v2}:
  \texttt{artifact\_manifest.json} records stable paths for checkpoints,
  metrics, canonical state-physics reports, compatibility reports,
  plots, summaries, and configuration.
\item
  HDF5 readers detect supported legacy LEBT, no-offset Sobol, and MEBT
  Sobol layouts through \texttt{synapticTrack.data.dataset\_compat};
  they report unavailable fields explicitly and never fabricate missing
  state or fixed-input data.
\item
  The GitHub Actions workflow covers \texttt{devel}, \texttt{main}, and
  \texttt{master}; package metadata includes \texttt{PyYAML} so CI
  installs the YAML dependency used by production configs.
\item
  \texttt{docs/site/} is generated from
  \texttt{docs/paper/synapticTrack\_consolidated.md} and no longer loads
  \texttt{polyfill.io}.
\end{itemize}

\hypertarget{maintained-lebt-production-workflow}{%
\paragraph{Maintained LEBT Production
Workflow}\label{maintained-lebt-production-workflow}}

The current production dataset workflow is documented in:

\begin{itemize}
\tightlist
\item
  \texttt{docs/lebt\_no\_offset\_sobol\_workflow.md}
\item
  \texttt{docs/notebook\_hygiene.md}
\item
  \texttt{docs/TRACK\_SETUP.md}
\item
  \texttt{docs/mebt\_surrogate\_input\_schema.md}
\item
  \texttt{docs/beam\_state\_physics\_api.md}
\item
  \texttt{docs/hdf5\_dataset\_compatibility.md}
\end{itemize}

Notebook-facing production APIs are maintained through direct workflow
submodule imports such as
\texttt{synapticTrack.simulation.no\_offset\_sobol\_production}. Package
initializers retain lazy compatibility exports, but new code should
avoid importing complete workflow stacks from \texttt{synapticTrack},
\texttt{synapticTrack.simulation}, or \texttt{synapticTrack.ml}.

\hypertarget{author}{%
\paragraph{Author}\label{author}}

\textbf{Chong Shik Park}\\
Accelerator Physicist / Machine Learning Researcher\\
Email: {[}kuphy@korea.ac.kr{]}\\
Affiliation: {[}Korea University, Sejong{]}

\hypertarget{status}{%
\paragraph{Status}\label{status}}

\textbf{In development} -- contributions, suggestions, and feedback are
welcome.

\hypertarget{program-roadmap-and-milestones}{%
\section{Program Roadmap and
Milestones}\label{program-roadmap-and-milestones}}

Source Markdown: \texttt{CODEX\_TASKS\_2026.md}

\hypertarget{codex-milestones}{%
\subsection{Codex Milestones}\label{codex-milestones}}

Use one branch or reviewable checkpoint per task.

\hypertarget{task-0-repository-survey}{%
\subsubsection{Task 0 --- Repository
survey}\label{task-0-repository-survey}}

Read \texttt{AGENTS.md} and
\texttt{notebooks/00\_program\_specification.ipynb}. Inspect beam
classes, backends, scanner analysis, units, packaging, and tests. Do not
modify code. Propose a file map and list ambiguities.

\hypertarget{task-1-gaussian-beam-utilities}{%
\subsubsection{Task 1 --- Gaussian beam
utilities}\label{task-1-gaussian-beam-utilities}}

Implement Twiss-to-covariance conversion and deterministic PyTorch
Gaussian sampling. Add tests for covariance, sampled moments, invalid
beta/emittance, dtype/device, and reproducibility. Add a notebook that
imports the module.

\hypertarget{task-2-lebt-simulation-adapter}{%
\subsubsection{Task 2 --- LEBT simulation
adapter}\label{task-2-lebt-simulation-adapter}}

Implement backend-neutral request/result schemas and an adapter for the
existing LEBT backend. Capture state, controls, diagnostics, last-WS
state, exit state, units, seed, validity mask, and failure reason.

\hypertarget{task-3-dataset-generator}{%
\subsubsection{Task 3 --- Dataset
generator}\label{task-3-dataset-generator}}

Implement Sobol-sampled, resumable LEBT data generation with validated
configuration, deterministic batch seeds, dry-run mode, metadata, and a
smoke-test dataset.

\hypertarget{task-4-mlp-surrogate}{%
\subsubsection{Task 4 --- MLP surrogate}\label{task-4-mlp-surrogate}}

Implement normalization, multi-output PyTorch MLP training, grouped
splits, checkpoints, metrics, plots, and reload tests.

\hypertarget{task-5-physics-regularized-surrogate}{%
\subsubsection{Task 5 --- Physics-regularized
surrogate}\label{task-5-physics-regularized-surrogate}}

Add positive-definite covariance outputs and moment--Twiss consistency
losses. Compare with the MLP using identical data and splits.

\hypertarget{task-6-beam-state-inference}{%
\subsubsection{Task 6 --- Beam-state
inference}\label{task-6-beam-state-inference}}

Implement a permutation-invariant set encoder and bounded differentiable
MAP refinement through a frozen surrogate. Test synthetic recovery and
identifiability.

\hypertarget{task-7-mebt-surrogate}{%
\subsubsection{Task 7 --- MEBT surrogate}\label{task-7-mebt-surrogate}}

Implement MEBT data generation and surrogate modeling with injected beam
state, corrector settings, and quadrupole settings. Include correction
trajectories.

\hypertarget{task-8-rl-environment-and-algorithms}{%
\subsubsection{Task 8 --- RL environment and
algorithms}\label{task-8-rl-environment-and-algorithms}}

Implement a Gymnasium-compatible surrogate environment, normalized
target-centered rewards, bounded incremental actions, loss termination,
SAC, TD3, response-matrix/SVD baseline, and optimizer baseline.

\hypertarget{task-9-single-ws-inversion}{%
\subsubsection{Task 9 --- Single-WS
inversion}\label{task-9-single-ws-inversion}}

Implement direct inversion from one WS image/profile plus settings to
the Gaussian beam state, followed by optional differentiable refinement.

\hypertarget{task-10-high-fidelity-differentiable-tracking}{%
\subsubsection{Task 10 --- High-fidelity differentiable
tracking}\label{task-10-high-fidelity-differentiable-tracking}}

Implement a deterministic external-tracker adapter and compare
finite-difference, local-neural, and custom-backward gradient methods.

\hypertarget{task-11-generative-reconstruction}{%
\subsubsection{Task 11 --- Generative
reconstruction}\label{task-11-generative-reconstruction}}

Implement a conditional VAE baseline, then conditional diffusion with
forward-consistency evaluation and uncertainty ensembles.

Source Markdown: \texttt{docs/codex\_tasks/README\_CODEX\_TASKS.md}

\hypertarget{codex-markdown-task-documents-for-lebt-surrogate-development}{%
\subsection{Codex Markdown Task Documents for LEBT Surrogate
Development}\label{codex-markdown-task-documents-for-lebt-surrogate-development}}

This folder contains separate Codex-ready markdown documents derived
from the latest discussion about LEBT surrogate modeling.

\hypertarget{documents}{%
\subsubsection{Documents}\label{documents}}

\begin{enumerate}
\def\labelenumi{\arabic{enumi}.}
\item
  \texttt{01\_input\_dimension\_strategy.md}

  Explains how to reduce or split the input dimension. Defines optics,
  orbit-response, and full surrogate model variants.
\item
  \texttt{02\_output\_storage\_strategy.md}

  Defines how to store original diagnostic outputs and final beam-state
  outputs separately so that different models can use different targets.
\item
  \texttt{03\_input\_parameter\_ranges.md}

  Defines recommended ranges for Twiss parameters, emittances, offsets,
  angles, current, steerers, and electrostatic quadrupoles.
\item
  \texttt{04\_training\_sample\_requirements.md}

  Defines recommended sample counts, dataset phases, learning curves,
  and validation split strategy.
\end{enumerate}

\hypertarget{recommended-repository-placement}{%
\subsubsection{Recommended repository
placement}\label{recommended-repository-placement}}

Copy this folder into:

\begin{Shaded}
\begin{Highlighting}[]
\NormalTok{synapticTrack/docs/codex\_tasks/}
\end{Highlighting}
\end{Shaded}

Recommended full structure:

\begin{Shaded}
\begin{Highlighting}[]
\NormalTok{synapticTrack/}
\NormalTok{+{-}{-} AGENTS.md}
\NormalTok{+{-}{-} CODEX\_TASKS\_2026.md}
\NormalTok{+{-}{-} docs/}
\NormalTok{|   +{-}{-} codex\_tasks/}
\NormalTok{|       +{-}{-} README\_CODEX\_TASKS.md}
\NormalTok{|       +{-}{-} 01\_input\_dimension\_strategy.md}
\NormalTok{|       +{-}{-} 02\_output\_storage\_strategy.md}
\NormalTok{|       +{-}{-} 03\_input\_parameter\_ranges.md}
\NormalTok{|       +{-}{-} 04\_training\_sample\_requirements.md}
\NormalTok{+{-}{-} notebooks/}
\end{Highlighting}
\end{Shaded}

\hypertarget{recommended-codex-usage}{%
\subsubsection{Recommended Codex usage}\label{recommended-codex-usage}}

Give Codex one document at a time.

Start with:

\begin{Shaded}
\begin{Highlighting}[]
\NormalTok{Read AGENTS.md and docs/codex\_tasks/01\_input\_dimension\_strategy.md.}
\NormalTok{Inspect the existing dataset generation and training pipeline.}
\NormalTok{Do not modify files yet. Summarize how the current code maps to the requested input{-}view strategy.}
\end{Highlighting}
\end{Shaded}

Then allow implementation only after reviewing Codex's survey.

\hypertarget{input-parameters-output-storage-and-sampling-requirements}{%
\section{Input Parameters, Output Storage, and Sampling
Requirements}\label{input-parameters-output-storage-and-sampling-requirements}}

Source Markdown:
\texttt{docs/codex\_tasks/01\_input\_dimension\_strategy.md}

\hypertarget{input-dimension-strategy-for-lebt-surrogate-models}{%
\subsection{01 --- Input Dimension Strategy for LEBT Surrogate
Models}\label{input-dimension-strategy-for-lebt-surrogate-models}}

\hypertarget{purpose-1}{%
\subsubsection{Purpose}\label{purpose-1}}

This document defines how to handle the LEBT surrogate input dimension
after Task 3 dataset generation.

Current situation:

\begin{itemize}
\tightlist
\item
  Current input dimension: approximately 24D.
\item
  Current output: diagnostic quantities, currently 12D.
\item
  The current input includes initial beam Twiss parameters, emittances,
  initial offsets/angles, beam current, steering magnet strengths, and
  electrostatic quadrupole strengths.
\end{itemize}

Main question:

\begin{quote}
Should \texttt{x0}, \texttt{xp0}, \texttt{y0}, \texttt{yp0}, and
\texttt{I\_beam} be removed from the surrogate input?
\end{quote}

\hypertarget{recommendation}{%
\subsubsection{Recommendation}\label{recommendation}}

Do not force one large 24D model to serve all purposes. Instead, split
the work into separate surrogate families:

\begin{enumerate}
\def\labelenumi{\arabic{enumi}.}
\tightlist
\item
  optics surrogate,
\item
  orbit-response surrogate,
\item
  optional full absolute diagnostic surrogate.
\end{enumerate}

This reduces sample complexity and makes each model physically
interpretable.

\begin{center}\rule{0.5\linewidth}{0.5pt}\end{center}

\hypertarget{can-x0-xp0-y0-yp0-be-removed}{%
\subsubsection{\texorpdfstring{1. Can \texttt{x0}, \texttt{xp0},
\texttt{y0}, \texttt{yp0} be
removed?}{1. Can x0, xp0, y0, yp0 be removed?}}\label{can-x0-xp0-y0-yp0-be-removed}}

Yes, but only for models that do not need to predict absolute beam
centroids from arbitrary injected offsets.

The variables

\[
x_0,\quad x'_0,\quad y_0,\quad y'_0
\]

mainly control centroid/orbit propagation. For linear optics and central
moments, centroid motion and rms-size evolution are approximately
separable.

Therefore:

\begin{itemize}
\tightlist
\item
  For rms-size and Twiss propagation, they can often be removed.
\item
  For absolute centroid prediction, they should either be kept or the
  target should be changed to centroid shift.
\end{itemize}

\begin{center}\rule{0.5\linewidth}{0.5pt}\end{center}

\hypertarget{recommended-model-split}{%
\subsubsection{2. Recommended model
split}\label{recommended-model-split}}

\hypertarget{model-a-optics-surrogate}{%
\paragraph{Model A: Optics surrogate}\label{model-a-optics-surrogate}}

Purpose:

\begin{itemize}
\tightlist
\item
  predict rms beam size,
\item
  predict output covariance/moments,
\item
  predict Twiss/emittance propagation,
\item
  support phase-space reconstruction.
\end{itemize}

Suggested input:

\[
[
\alpha_x,\beta_x,\epsilon_x,
\alpha_y,\beta_y,\epsilon_y,
I_{\mathrm{beam}},
\mathrm{equad}_1,\ldots,\mathrm{equad}_{N_q}
]
\]

If current is fixed or weakly influential, use:

\[
[
\alpha_x,\beta_x,\epsilon_x,
\alpha_y,\beta_y,\epsilon_y,
\mathrm{equad}_1,\ldots,\mathrm{equad}_{N_q}
]
\]

Do not include:

\[
x_0,\quad x'_0,\quad y_0,\quad y'_0
\]

unless strong nonlinear orbit-size coupling or aperture clipping is
important.

Suggested output:

\begin{itemize}
\tightlist
\item
  rms sizes at diagnostic locations,
\item
  final beam moments,
\item
  Twiss/emittance at last wire scanner and LEBT exit.
\end{itemize}

\begin{center}\rule{0.5\linewidth}{0.5pt}\end{center}

\hypertarget{model-b-orbit-response-surrogate}{%
\paragraph{Model B: Orbit-response
surrogate}\label{model-b-orbit-response-surrogate}}

Purpose:

\begin{itemize}
\tightlist
\item
  support LEBT orbit correction,
\item
  avoid dependence on unknown injected centroid/angle,
\item
  support real experimental validation where initial state may be
  unknown.
\end{itemize}

Suggested input:

\[
[
\Delta \mathrm{steerer}_1,\ldots,\Delta \mathrm{steerer}_{N_s},
\mathrm{equad}_1,\ldots,\mathrm{equad}_{N_q}
]
\]

or

\[
[
\mathrm{steerer}_1,\ldots,\mathrm{steerer}_{N_s},
\mathrm{equad}_1,\ldots,\mathrm{equad}_{N_q}
]
\]

Suggested output:

\[
[
\Delta x_c(s_1),\Delta y_c(s_1),\ldots,
\Delta x_c(s_N),\Delta y_c(s_N)
]
\]

where

\[
\Delta x_c(s_i)
=
x_c(s_i;\mathbf{u})
-
x_c(s_i;\mathbf{u}_{\mathrm{ref}})
\]

and similarly for \(y\).

This model does not require explicit initial offsets if the reference
orbit is measured.

\begin{center}\rule{0.5\linewidth}{0.5pt}\end{center}

\hypertarget{model-c-full-absolute-diagnostic-surrogate}{%
\paragraph{Model C: Full absolute diagnostic
surrogate}\label{model-c-full-absolute-diagnostic-surrogate}}

Purpose:

\begin{itemize}
\tightlist
\item
  full simulation replacement,
\item
  absolute prediction of centroid and rms diagnostics,
\item
  useful for synthetic-data studies and model completeness.
\end{itemize}

Suggested input:

\[
[
\alpha_x,\beta_x,\epsilon_x,
\alpha_y,\beta_y,\epsilon_y,
x_0,x'_0,y_0,y'_0,
I_{\mathrm{beam}},
\mathrm{steerers},
\mathrm{equads}
]
\]

Suggested output:

\[
[
x_c(s_i),y_c(s_i),\sigma_x(s_i),\sigma_y(s_i)
]_{i=1}^{N}
\]

This is the most complete model but requires more training samples.

\begin{center}\rule{0.5\linewidth}{0.5pt}\end{center}

\hypertarget{what-to-do-with-beam-current}{%
\subsubsection{3. What to do with beam
current?}\label{what-to-do-with-beam-current}}

Beam current \(I_{\mathrm{beam}}\) should be handled according to
physics relevance.

\hypertarget{option-1-fixed-current-model}{%
\paragraph{Option 1: fixed-current
model}\label{option-1-fixed-current-model}}

Use this if all simulation and experiment are done near one nominal
current:

\[
I = I_{\mathrm{nom}}
\]

Do not include current as input.

\hypertarget{option-2-narrow-current-model}{%
\paragraph{Option 2: narrow current
model}\label{option-2-narrow-current-model}}

Use this if current varies mildly:

\[
I \in [0.8I_{\mathrm{nom}},1.2I_{\mathrm{nom}}]
\]

Include current as input only in optics/state models.

\hypertarget{option-3-current-indexed-models}{%
\paragraph{Option 3: current-indexed
models}\label{option-3-current-indexed-models}}

Train separate models:

\[
F_{\theta}^{I_1},\quad F_{\theta}^{I_2},\quad F_{\theta}^{I_3}
\]

This is useful if power-supply or beam-current settings are discrete.

\hypertarget{option-4-full-current-dependent-model}{%
\paragraph{Option 4: full current-dependent
model}\label{option-4-full-current-dependent-model}}

Use if current strongly affects space charge and rms evolution:

\[
I \in [0.5I_{\mathrm{nom}},1.5I_{\mathrm{nom}}]
\]

This requires more data.

\begin{center}\rule{0.5\linewidth}{0.5pt}\end{center}

\hypertarget{recommended-implementation-decision}{%
\subsubsection{4. Recommended implementation
decision}\label{recommended-implementation-decision}}

For the next Codex milestone, implement model-selection support in the
dataset loader.

Do not delete existing input columns. Instead, allow different input
views:

\begin{Shaded}
\begin{Highlighting}[]
\NormalTok{input\_view }\OperatorTok{=} \StringTok{"full"}
\NormalTok{input\_view }\OperatorTok{=} \StringTok{"optics"}
\NormalTok{input\_view }\OperatorTok{=} \StringTok{"orbit\_response"}
\end{Highlighting}
\end{Shaded}

Suggested behavior:

\begin{itemize}
\tightlist
\item
  \texttt{full}: all available inputs.
\item
  \texttt{optics}: Twiss, emittance, optional current, equads.
\item
  \texttt{orbit\_response}: steerers, equads, possibly delta steerers.
\item
  \texttt{custom}: user-provided list of input columns.
\end{itemize}

\begin{center}\rule{0.5\linewidth}{0.5pt}\end{center}

\hypertarget{suggested-repository-changes}{%
\subsubsection{5. Suggested repository
changes}\label{suggested-repository-changes}}

Possible files:

\begin{Shaded}
\begin{Highlighting}[]
\NormalTok{src/synaptictrack/data/input\_views.py}
\NormalTok{src/synaptictrack/data/schemas.py}
\NormalTok{src/synaptictrack/data/datasets.py}
\NormalTok{tests/test\_input\_views.py}
\end{Highlighting}
\end{Shaded}

Suggested functionality:

\begin{Shaded}
\begin{Highlighting}[]
\NormalTok{get\_input\_columns(view: }\BuiltInTok{str}\NormalTok{, }\OperatorTok{*}\NormalTok{, include\_current: }\BuiltInTok{bool} \OperatorTok{=} \VariableTok{True}\NormalTok{) }\OperatorTok{{-}\textgreater{}} \BuiltInTok{list}\NormalTok{[}\BuiltInTok{str}\NormalTok{]}
\NormalTok{select\_input\_view(dataframe, view: }\BuiltInTok{str}\NormalTok{, include\_current: }\BuiltInTok{bool} \OperatorTok{=} \VariableTok{True}\NormalTok{)}
\end{Highlighting}
\end{Shaded}

\begin{center}\rule{0.5\linewidth}{0.5pt}\end{center}

\hypertarget{required-tests}{%
\subsubsection{6. Required tests}\label{required-tests}}

Codex should add tests that verify:

\begin{enumerate}
\def\labelenumi{\arabic{enumi}.}
\tightlist
\item
  the original full input view still works;
\item
  the optics input view removes \texttt{x0}, \texttt{xp0}, \texttt{y0},
  \texttt{yp0};
\item
  the optics input view optionally includes or excludes
  \texttt{I\_beam};
\item
  the orbit-response view includes steering and equad variables;
\item
  input column ordering is deterministic;
\item
  invalid view names raise a clear error;
\item
  training code can accept different input dimensions.
\end{enumerate}

\begin{center}\rule{0.5\linewidth}{0.5pt}\end{center}

\hypertarget{codex-prompt}{%
\subsubsection{7. Codex prompt}\label{codex-prompt}}

Use the following prompt for Codex:

\begin{Shaded}
\begin{Highlighting}[]
\NormalTok{Read AGENTS.md, CODEX\_TASKS\_2026.md, and this document:}
\NormalTok{docs/codex\_tasks/01\_input\_dimension\_strategy.md.}

\NormalTok{Task:}
\NormalTok{Implement input{-}view support for LEBT surrogate datasets.}

\NormalTok{Requirements:}
\NormalTok{1. Do not remove existing dataset columns.}
\NormalTok{2. Add reusable input{-}view utilities.}
\NormalTok{3. Support at least:}
\NormalTok{   {-} full}
\NormalTok{   {-} optics}
\NormalTok{   {-} orbit\_response}
\NormalTok{   {-} custom}
\NormalTok{4. In the optics view, allow \textasciigrave{}include\_current=True/False\textasciigrave{}.}
\NormalTok{5. Preserve deterministic column ordering.}
\NormalTok{6. Add tests for all input views.}
\NormalTok{7. Update the relevant notebook to demonstrate training with the reduced optics input.}
\NormalTok{8. Do not modify the simulation backend.}
\NormalTok{9. Do not implement MEBT or RL in this task.}

\NormalTok{After implementation, summarize changed files, tests run, and remaining assumptions.}
\end{Highlighting}
\end{Shaded}

Source Markdown:
\texttt{docs/codex\_tasks/02\_output\_storage\_strategy.md}

\hypertarget{output-storage-strategy-for-diagnostics-and-beam-state-targets}{%
\subsection{02 --- Output Storage Strategy for Diagnostics and
Beam-State
Targets}\label{output-storage-strategy-for-diagnostics-and-beam-state-targets}}

\hypertarget{purpose-2}{%
\subsubsection{Purpose}\label{purpose-2}}

This document defines how to store original diagnostic outputs and
additional beam-state outputs so that different surrogate models can use
different targets.

Main question:

\begin{quote}
Can original diagnostic outputs and beam-state outputs be stored
separately and reused by different models?
\end{quote}

Answer:

Yes. This is strongly recommended.

\begin{center}\rule{0.5\linewidth}{0.5pt}\end{center}

\hypertarget{core-recommendation}{%
\subsubsection{1. Core recommendation}\label{core-recommendation}}

Do not store all targets only as one flat array.

Instead, store outputs in separate named groups:

\begin{Shaded}
\begin{Highlighting}[]
\NormalTok{output/}
\NormalTok{+{-}{-} diagnostics/}
\NormalTok{+{-}{-} orbit\_response/}
\NormalTok{+{-}{-} state\_last\_ws/}
\NormalTok{+{-}{-} state\_lebt\_exit/}
\end{Highlighting}
\end{Shaded}

This allows the same simulation dataset to support multiple models:

\begin{itemize}
\tightlist
\item
  diagnostic-only surrogate,
\item
  orbit-response surrogate,
\item
  final-state surrogate,
\item
  multi-task surrogate,
\item
  phase-space inference model.
\end{itemize}

\begin{center}\rule{0.5\linewidth}{0.5pt}\end{center}

\hypertarget{recommended-output-groups}{%
\subsubsection{2. Recommended output
groups}\label{recommended-output-groups}}

\hypertarget{y_diag}{%
\paragraph{\texorpdfstring{2.1
\texttt{y\_diag}}{2.1 y\_diag}}\label{y_diag}}

Original diagnostic outputs.

Possible contents:

\[
[
x_c(s_1),y_c(s_1),\sigma_x(s_1),\sigma_y(s_1),
\ldots,
x_c(s_N),y_c(s_N),\sigma_x(s_N),\sigma_y(s_N)
]
\]

If the current project output is 12D, keep it exactly as-is for backward
compatibility.

Do not break the existing training pipeline.

\begin{center}\rule{0.5\linewidth}{0.5pt}\end{center}

\hypertarget{y_orbit_response}{%
\paragraph{\texorpdfstring{2.2
\texttt{y\_orbit\_response}}{2.2 y\_orbit\_response}}\label{y_orbit_response}}

Centroid changes relative to a reference setting.

\[
\Delta x_c(s_i)
=
x_c(s_i;\mathbf{u})
-
x_c(s_i;\mathbf{u}_{\mathrm{ref}})
\]

\[
\Delta y_c(s_i)
=
y_c(s_i;\mathbf{u})
-
y_c(s_i;\mathbf{u}_{\mathrm{ref}})
\]

Suggested array:

\[
[
\Delta x_c(s_1),\Delta y_c(s_1),\ldots,
\Delta x_c(s_N),\Delta y_c(s_N)
]
\]

This target is useful for orbit correction because it avoids requiring
exact knowledge of the initial injected centroid.

\begin{center}\rule{0.5\linewidth}{0.5pt}\end{center}

\hypertarget{y_state_last_ws}{%
\paragraph{\texorpdfstring{2.3
\texttt{y\_state\_last\_ws}}{2.3 y\_state\_last\_ws}}\label{y_state_last_ws}}

Beam state at the last wire scanner.

Recommended contents:

\[
[
x_c,\ x'_c,\ y_c,\ y'_c,
\langle x^2\rangle,\langle xx'\rangle,\langle x'^2\rangle,
\langle y^2\rangle,\langle yy'\rangle,\langle y'^2\rangle
]
\]

This is a 10D state target.

Also store derived Twiss values as either metadata or optional target:

\[
[
\alpha_x,\beta_x,\epsilon_x,
\alpha_y,\beta_y,\epsilon_y
]
\]

\begin{center}\rule{0.5\linewidth}{0.5pt}\end{center}

\hypertarget{y_state_lebt_exit}{%
\paragraph{\texorpdfstring{2.4
\texttt{y\_state\_lebt\_exit}}{2.4 y\_state\_lebt\_exit}}\label{y_state_lebt_exit}}

Beam state at the LEBT exit.

Use the same 10D structure:

\[
[
x_c,\ x'_c,\ y_c,\ y'_c,
\langle x^2\rangle,\langle xx'\rangle,\langle x'^2\rangle,
\langle y^2\rangle,\langle yy'\rangle,\langle y'^2\rangle
]
\]

This is the most important output for the later MEBT surrogate because
it becomes the injected beam state.

\begin{center}\rule{0.5\linewidth}{0.5pt}\end{center}

\hypertarget{why-moments-are-preferred-over-twiss-as-training-targets}{%
\subsubsection{3. Why moments are preferred over Twiss as training
targets}\label{why-moments-are-preferred-over-twiss-as-training-targets}}

The central second moments are

\[
\langle x^2\rangle,\quad
\langle xx'\rangle,\quad
\langle x'^2\rangle.
\]

From these,

\[
\epsilon_x
=
\sqrt{
\langle x^2\rangle\langle x'^2\rangle
-
\langle xx'\rangle^2
}
\]

\[
\beta_x
=
\frac{\langle x^2\rangle}{\epsilon_x}
\]

\[
\alpha_x
=
-\frac{\langle xx'\rangle}{\epsilon_x}.
\]

Moments are closer to raw particle statistics and are generally more
stable as supervised-learning labels.

Therefore:

\begin{itemize}
\tightlist
\item
  train primarily on moments;
\item
  derive Twiss for analysis;
\item
  optionally include Twiss in metadata or auxiliary targets.
\end{itemize}

\begin{center}\rule{0.5\linewidth}{0.5pt}\end{center}

\hypertarget{dataset-schema-recommendation}{%
\subsubsection{4. Dataset schema
recommendation}\label{dataset-schema-recommendation}}

A sample should support this structure:

\begin{Shaded}
\begin{Highlighting}[]
\NormalTok{sample/}
\NormalTok{+{-}{-} input/}
\NormalTok{|   +{-}{-} twiss}
\NormalTok{|   +{-}{-} emittance}
\NormalTok{|   +{-}{-} initial\_centroid}
\NormalTok{|   +{-}{-} current}
\NormalTok{|   +{-}{-} steerers}
\NormalTok{|   +{-}{-} equads}
\NormalTok{+{-}{-} output/}
\NormalTok{|   +{-}{-} diagnostics}
\NormalTok{|   +{-}{-} orbit\_response}
\NormalTok{|   +{-}{-} state\_last\_ws}
\NormalTok{|   +{-}{-} state\_lebt\_exit}
\NormalTok{+{-}{-} metadata/}
\NormalTok{|   +{-}{-} units}
\NormalTok{|   +{-}{-} seed}
\NormalTok{|   +{-}{-} lattice\_version}
\NormalTok{|   +{-}{-} tracking\_code}
\NormalTok{|   +{-}{-} valid\_flag}
\NormalTok{|   +{-}{-} failure\_reason}
\end{Highlighting}
\end{Shaded}

If using HDF5 or Zarr, use groups. If using tabular formats, use column
prefixes such as:

\begin{Shaded}
\begin{Highlighting}[]
\NormalTok{diag/x\_c\_WS01}
\NormalTok{diag/y\_c\_WS01}
\NormalTok{diag/sigma\_x\_WS01}
\NormalTok{diag/sigma\_y\_WS01}

\NormalTok{state\_exit/x\_c}
\NormalTok{state\_exit/xp\_c}
\NormalTok{state\_exit/y\_c}
\NormalTok{state\_exit/yp\_c}
\NormalTok{state\_exit/x2}
\NormalTok{state\_exit/xxp}
\NormalTok{state\_exit/xp2}
\NormalTok{state\_exit/y2}
\NormalTok{state\_exit/yyp}
\NormalTok{state\_exit/yp2}
\end{Highlighting}
\end{Shaded}

\begin{center}\rule{0.5\linewidth}{0.5pt}\end{center}

\hypertarget{target-selection-api}{%
\subsubsection{5. Target selection API}\label{target-selection-api}}

The dataset loader should support different target views:

\begin{Shaded}
\begin{Highlighting}[]
\NormalTok{target\_view }\OperatorTok{=} \StringTok{"diagnostics"}
\NormalTok{target\_view }\OperatorTok{=} \StringTok{"orbit\_response"}
\NormalTok{target\_view }\OperatorTok{=} \StringTok{"state\_exit"}
\NormalTok{target\_view }\OperatorTok{=} \StringTok{"state\_last\_ws"}
\NormalTok{target\_view }\OperatorTok{=} \StringTok{"diag\_plus\_exit\_state"}
\NormalTok{target\_view }\OperatorTok{=} \StringTok{"diag\_plus\_all\_states"}
\NormalTok{target\_view }\OperatorTok{=} \StringTok{"custom"}
\end{Highlighting}
\end{Shaded}

Examples:

\begin{Shaded}
\begin{Highlighting}[]
\NormalTok{dataset }\OperatorTok{=}\NormalTok{ LEBTSurrogateDataset(}
\NormalTok{    path}\OperatorTok{=}\StringTok{"data/lebt\_dataset\_v3.h5"}\NormalTok{,}
\NormalTok{    input\_view}\OperatorTok{=}\StringTok{"optics"}\NormalTok{,}
\NormalTok{    target\_view}\OperatorTok{=}\StringTok{"state\_exit"}\NormalTok{,}
\NormalTok{)}
\end{Highlighting}
\end{Shaded}

\begin{Shaded}
\begin{Highlighting}[]
\NormalTok{dataset }\OperatorTok{=}\NormalTok{ LEBTSurrogateDataset(}
\NormalTok{    path}\OperatorTok{=}\StringTok{"data/lebt\_dataset\_v3.h5"}\NormalTok{,}
\NormalTok{    input\_view}\OperatorTok{=}\StringTok{"full"}\NormalTok{,}
\NormalTok{    target\_view}\OperatorTok{=}\StringTok{"diag\_plus\_all\_states"}\NormalTok{,}
\NormalTok{)}
\end{Highlighting}
\end{Shaded}

\begin{center}\rule{0.5\linewidth}{0.5pt}\end{center}

\hypertarget{multi-task-output-strategy}{%
\subsubsection{6. Multi-task output
strategy}\label{multi-task-output-strategy}}

A multi-task surrogate can predict:

\[
\mathbf{y}
=
[
\mathbf{y}_{\mathrm{diag}},
\mathbf{y}_{\mathrm{state,lastWS}},
\mathbf{y}_{\mathrm{state,exit}}
]
\]

The training loss should be grouped:

\[
\mathcal{L}
=
\lambda_{\mathrm{diag}}\mathcal{L}_{\mathrm{diag}}
+
\lambda_{\mathrm{lastWS}}\mathcal{L}_{\mathrm{lastWS}}
+
\lambda_{\mathrm{exit}}\mathcal{L}_{\mathrm{exit}}.
\]

Do not force all outputs to have equal weight. Normalize each group
separately.

\begin{center}\rule{0.5\linewidth}{0.5pt}\end{center}

\hypertarget{required-tests-1}{%
\subsubsection{7. Required tests}\label{required-tests-1}}

Codex should add tests that verify:

\begin{enumerate}
\def\labelenumi{\arabic{enumi}.}
\tightlist
\item
  original diagnostic output path remains unchanged;
\item
  \texttt{y\_diag} can be loaded independently;
\item
  \texttt{y\_state\_last\_ws} can be loaded independently;
\item
  \texttt{y\_state\_lebt\_exit} can be loaded independently;
\item
  combined target views concatenate in deterministic order;
\item
  derived Twiss parameters from moments are physically valid;
\item
  covariance determinants are positive for generated labels;
\item
  old notebooks still run with diagnostic-only output.
\end{enumerate}

\begin{center}\rule{0.5\linewidth}{0.5pt}\end{center}

\hypertarget{codex-prompt-1}{%
\subsubsection{8. Codex prompt}\label{codex-prompt-1}}

Use the following prompt for Codex:

\begin{Shaded}
\begin{Highlighting}[]
\NormalTok{Read AGENTS.md, CODEX\_TASKS\_2026.md, and this document:}
\NormalTok{docs/codex\_tasks/02\_output\_storage\_strategy.md.}

\NormalTok{Task:}
\NormalTok{Add separate output{-}target storage and target{-}view loading for LEBT surrogate datasets.}

\NormalTok{Requirements:}
\NormalTok{1. Preserve the existing diagnostic output exactly.}
\NormalTok{2. Add separate groups or column prefixes for:}
\NormalTok{   {-} diagnostics}
\NormalTok{   {-} orbit\_response}
\NormalTok{   {-} state\_last\_ws}
\NormalTok{   {-} state\_lebt\_exit}
\NormalTok{3. For beam{-}state outputs, store central moments and centroids/angles.}
\NormalTok{4. Store derived Twiss/emittance as metadata or optional targets.}
\NormalTok{5. Add target{-}view selection in the dataset loader.}
\NormalTok{6. Support combined target views for multi{-}task learning.}
\NormalTok{7. Add tests for shape, ordering, backward compatibility, and physical validity.}
\NormalTok{8. Update the notebook to demonstrate diagnostic{-}only and diagnostic+state training.}
\NormalTok{9. Do not implement RL or MEBT in this task.}

\NormalTok{After implementation, summarize changed files, tests run, and remaining assumptions.}
\end{Highlighting}
\end{Shaded}

Source Markdown:
\texttt{docs/codex\_tasks/03\_input\_parameter\_ranges.md}

\hypertarget{recommended-input-parameter-ranges-for-lebt-surrogate-data-generation}{%
\subsection{03 --- Recommended Input Parameter Ranges for LEBT Surrogate
Data
Generation}\label{recommended-input-parameter-ranges-for-lebt-surrogate-data-generation}}

\hypertarget{purpose-3}{%
\subsubsection{Purpose}\label{purpose-3}}

This document defines recommended sampling ranges for LEBT surrogate
model data generation.

Main question:

\begin{quote}
What ranges should be used for input parameters?
\end{quote}

The exact values should ultimately come from the RAON LEBT nominal
optics, hardware limits, and trusted simulation acceptance. Until those
are fully fixed, use relative ranges around nominal values.

\begin{center}\rule{0.5\linewidth}{0.5pt}\end{center}

\hypertarget{range-strategy}{%
\subsubsection{1. Range strategy}\label{range-strategy}}

Use three nested ranges:

\begin{enumerate}
\def\labelenumi{\arabic{enumi}.}
\tightlist
\item
  nominal range,
\item
  training range,
\item
  stress-test range.
\end{enumerate}

Do not use one wide uniform box for everything.

The recommended strategy is:

\begin{itemize}
\tightlist
\item
  train densely near normal operation;
\item
  include enough mismatch for robustness;
\item
  reserve extreme cases for validation and stress testing;
\item
  avoid wasting samples in always-lost or physically irrelevant regions.
\end{itemize}

\begin{center}\rule{0.5\linewidth}{0.5pt}\end{center}

\hypertarget{twiss-parameter-ranges}{%
\subsubsection{2. Twiss parameter ranges}\label{twiss-parameter-ranges}}

Let the nominal input Twiss parameters be:

\[
\alpha_{x0},\quad \beta_{x0},\quad \epsilon_{x0}
\]

and similarly for the \(y\) plane.

\hypertarget{training-range}{%
\paragraph{2.1 Training range}\label{training-range}}

Use:

\[
\alpha_x \in [\alpha_{x0}-2,\alpha_{x0}+2]
\]

\[
\beta_x \in [0.5\beta_{x0},2.0\beta_{x0}]
\]

\[
\epsilon_x \in [0.5\epsilon_{x0},2.0\epsilon_{x0}]
\]

Apply the same structure to \(y\):

\[
\alpha_y \in [\alpha_{y0}-2,\alpha_{y0}+2]
\]

\[
\beta_y \in [0.5\beta_{y0},2.0\beta_{y0}]
\]

\[
\epsilon_y \in [0.5\epsilon_{y0},2.0\epsilon_{y0}]
\]

\hypertarget{stress-test-range}{%
\paragraph{2.2 Stress-test range}\label{stress-test-range}}

Use:

\[
\alpha_x \in [\alpha_{x0}-4,\alpha_{x0}+4]
\]

\[
\beta_x \in [0.25\beta_{x0},4.0\beta_{x0}]
\]

\[
\epsilon_x \in [0.25\epsilon_{x0},3.0\epsilon_{x0}]
\]

Again, apply the same structure to \(y\).

Stress-test samples should be used primarily for validation, not as the
dominant training population.

\begin{center}\rule{0.5\linewidth}{0.5pt}\end{center}

\hypertarget{initial-offsets-and-angles}{%
\subsubsection{3. Initial offsets and
angles}\label{initial-offsets-and-angles}}

If included in the full surrogate:

\hypertarget{nominal-training-range}{%
\paragraph{3.1 Nominal training range}\label{nominal-training-range}}

\[
x_0,y_0 \in [-2,2]\ \mathrm{mm}
\]

\[
x'_0,y'_0 \in [-10,10]\ \mathrm{mrad}
\]

\hypertarget{stress-test-range-1}{%
\paragraph{3.2 Stress-test range}\label{stress-test-range-1}}

\[
x_0,y_0 \in [-5,5]\ \mathrm{mm}
\]

\[
x'_0,y'_0 \in [-30,30]\ \mathrm{mrad}
\]

If these values exceed the realistic LEBT acceptance, reduce them. Do
not oversample regions where the beam is always lost.

\begin{center}\rule{0.5\linewidth}{0.5pt}\end{center}

\hypertarget{beam-current}{%
\subsubsection{4. Beam current}\label{beam-current}}

If current is fixed, do not include it as an input.

If included, use one of the following.

\hypertarget{narrow-current-model}{%
\paragraph{4.1 Narrow current model}\label{narrow-current-model}}

\[
I \in [0.8I_{\mathrm{nom}},1.2I_{\mathrm{nom}}]
\]

Recommended for the first current-dependent surrogate.

\hypertarget{general-current-model}{%
\paragraph{4.2 General current model}\label{general-current-model}}

\[
I \in [0.5I_{\mathrm{nom}},1.5I_{\mathrm{nom}}]
\]

Use only after the narrow model works.

\hypertarget{stress-test-current-range}{%
\paragraph{4.3 Stress-test current
range}\label{stress-test-current-range}}

\[
I \in [0,2I_{\mathrm{nom}}]
\]

Use for validation, not first training.

\begin{center}\rule{0.5\linewidth}{0.5pt}\end{center}

\hypertarget{steering-magnet-ranges}{%
\subsubsection{5. Steering magnet ranges}\label{steering-magnet-ranges}}

If simulation uses kick angles:

\[
\theta_i \in [-\theta_{i,\max},\theta_{i,\max}]
\]

However, the training set should be denser in the operational region:

\[
\theta_i \in [-0.5\theta_{i,\max},0.5\theta_{i,\max}]
\]

For real experimental validation, prefer calibrated hardware variables:

\[
V_i \in [V_{i,\min},V_{i,\max}]
\]

or

\[
I_i \in [I_{i,\min},I_{i,\max}]
\]

with conversion to kick angle handled by an actuator model.

\begin{center}\rule{0.5\linewidth}{0.5pt}\end{center}

\hypertarget{electrostatic-quadrupole-ranges}{%
\subsubsection{6. Electrostatic quadrupole
ranges}\label{electrostatic-quadrupole-ranges}}

Let \(G_{j0}\) be the nominal equad strength.

\hypertarget{first-training-range}{%
\paragraph{6.1 First training range}\label{first-training-range}}

\[
G_j \in [0.8G_{j0},1.2G_{j0}]
\]

\hypertarget{wider-training-range}{%
\paragraph{6.2 Wider training range}\label{wider-training-range}}

\[
G_j \in [0.6G_{j0},1.4G_{j0}]
\]

\hypertarget{hardware-limit-range}{%
\paragraph{6.3 Hardware-limit range}\label{hardware-limit-range}}

\[
G_j \in [G_{j,\min},G_{j,\max}]
\]

Use the hardware-limit range only when needed for recovery or robustness
studies.

\begin{center}\rule{0.5\linewidth}{0.5pt}\end{center}

\hypertarget{recommended-dataset-versions}{%
\subsubsection{7. Recommended dataset
versions}\label{recommended-dataset-versions}}

\hypertarget{dataset-v1-optics-surrogate}{%
\paragraph{Dataset v1: optics
surrogate}\label{dataset-v1-optics-surrogate}}

Purpose:

\begin{itemize}
\tightlist
\item
  rms and beam-state propagation.
\end{itemize}

Input:

\[
[
\alpha_x,\beta_x,\epsilon_x,
\alpha_y,\beta_y,\epsilon_y,
\mathrm{equads}
]
\]

Optional:

\[
I_{\mathrm{beam}}
\]

Ranges:

\[
\alpha = \alpha_0\pm 2
\]

\[
\beta = 0.5\text{--}2.0\times\beta_0
\]

\[
\epsilon = 0.5\text{--}2.0\times\epsilon_0
\]

\[
G = 0.8\text{--}1.2\times G_0
\]

\begin{center}\rule{0.5\linewidth}{0.5pt}\end{center}

\hypertarget{dataset-v2-orbit-response-surrogate}{%
\paragraph{Dataset v2: orbit-response
surrogate}\label{dataset-v2-orbit-response-surrogate}}

Purpose:

\begin{itemize}
\tightlist
\item
  LEBT orbit correction.
\end{itemize}

Input:

\[
[
\mathrm{steerers},
\mathrm{equads}
]
\]

Output:

\[
[
\Delta x_c,\Delta y_c
]
\]

Ranges:

\[
\theta = 0.5\text{--}1.0\times \text{operational steering range}
\]

\[
G = 0.8\text{--}1.2\times G_0
\]

\begin{center}\rule{0.5\linewidth}{0.5pt}\end{center}

\hypertarget{dataset-v3-full-surrogate}{%
\paragraph{Dataset v3: full surrogate}\label{dataset-v3-full-surrogate}}

Purpose:

\begin{itemize}
\tightlist
\item
  complete absolute diagnostic prediction.
\end{itemize}

Input:

\[
[
\alpha,\beta,\epsilon,
x_0,x'_0,y_0,y'_0,
I,
\mathrm{steerers},
\mathrm{equads}
]
\]

Use after v1 and v2 are validated.

\begin{center}\rule{0.5\linewidth}{0.5pt}\end{center}

\hypertarget{sampling-method}{%
\subsubsection{8. Sampling method}\label{sampling-method}}

Use Sobol or Latin-hypercube sampling.

Recommended mixture:

\begin{itemize}
\tightlist
\item
  60\% nominal/training range,
\item
  25\% mismatch/recovery range,
\item
  15\% stress-test or boundary range.
\end{itemize}

For each simulation sample, store:

\begin{itemize}
\tightlist
\item
  input vector,
\item
  physical parameter values before normalization,
\item
  validity flag,
\item
  failure reason,
\item
  aperture/loss status,
\item
  random seed,
\item
  lattice version,
\item
  units.
\end{itemize}

\begin{center}\rule{0.5\linewidth}{0.5pt}\end{center}

\hypertarget{required-tests-2}{%
\subsubsection{9. Required tests}\label{required-tests-2}}

Codex should add tests that verify:

\begin{enumerate}
\def\labelenumi{\arabic{enumi}.}
\tightlist
\item
  generated ranges obey configuration bounds;
\item
  beta and emittance are always positive;
\item
  current is non-negative;
\item
  steering and equad settings obey bounds;
\item
  stress-test ranges are separated from training ranges;
\item
  invalid range configurations raise clear errors;
\item
  sampling is reproducible for a fixed seed;
\item
  generated samples can be passed to the simulation adapter.
\end{enumerate}

\begin{center}\rule{0.5\linewidth}{0.5pt}\end{center}

\hypertarget{codex-prompt-2}{%
\subsubsection{10. Codex prompt}\label{codex-prompt-2}}

Use the following prompt for Codex:

\begin{Shaded}
\begin{Highlighting}[]
\NormalTok{Read AGENTS.md, CODEX\_TASKS\_2026.md, and this document:}
\NormalTok{docs/codex\_tasks/03\_input\_parameter\_ranges.md.}

\NormalTok{Task:}
\NormalTok{Implement configurable input{-}parameter range definitions for LEBT surrogate}
\NormalTok{data generation.}

\NormalTok{Requirements:}
\NormalTok{1. Define range configuration objects for:}
\NormalTok{   {-} Twiss parameters}
\NormalTok{   {-} emittances}
\NormalTok{   {-} initial offsets and angles}
\NormalTok{   {-} beam current}
\NormalTok{   {-} steering magnets}
\NormalTok{   {-} electrostatic quadrupoles}
\NormalTok{2. Support nominal, training, and stress{-}test ranges.}
\NormalTok{3. Support dataset presets:}
\NormalTok{   {-} optics}
\NormalTok{   {-} orbit\_response}
\NormalTok{   {-} full}
\NormalTok{4. Use Sobol or Latin{-}hypercube sampling if already available; otherwise}
\NormalTok{   implement a clean interface and use a deterministic fallback.}
\NormalTok{5. Validate all physical bounds.}
\NormalTok{6. Add tests for reproducibility, bounds, invalid configs, and dataset presets.}
\NormalTok{7. Do not modify the tracking backend.}
\NormalTok{8. Do not train neural networks in this task.}

\NormalTok{After implementation, summarize changed files, tests run, and remaining assumptions.}
\end{Highlighting}
\end{Shaded}

Source Markdown:
\texttt{docs/codex\_tasks/04\_training\_sample\_requirements.md}

\hypertarget{training-sample-requirements-for-lebt-surrogate-models}{%
\subsection{04 --- Training Sample Requirements for LEBT Surrogate
Models}\label{training-sample-requirements-for-lebt-surrogate-models}}

\hypertarget{purpose-4}{%
\subsubsection{Purpose}\label{purpose-4}}

This document estimates how much simulation data is needed for training
the LEBT surrogate models.

Main question:

\begin{quote}
How many samples are needed for training the surrogate model?
\end{quote}

The answer depends on:

\begin{itemize}
\tightlist
\item
  input dimension,
\item
  output dimension,
\item
  nonlinearities,
\item
  model type,
\item
  sampling range,
\item
  whether absolute outputs or response shifts are predicted,
\item
  simulation noise from macroparticle statistics.
\end{itemize}

\begin{center}\rule{0.5\linewidth}{0.5pt}\end{center}

\hypertarget{summary-recommendation}{%
\subsubsection{1. Summary recommendation}\label{summary-recommendation}}

Do not generate a huge full 24D production dataset immediately.

Use a staged data-generation strategy:

\begin{enumerate}
\def\labelenumi{\arabic{enumi}.}
\tightlist
\item
  smoke-test dataset,
\item
  first real dataset,
\item
  production reduced-input dataset,
\item
  production full-input dataset,
\item
  active-learning refinement.
\end{enumerate}

\begin{center}\rule{0.5\linewidth}{0.5pt}\end{center}

\hypertarget{recommended-sample-counts}{%
\subsubsection{2. Recommended sample
counts}\label{recommended-sample-counts}}

\hypertarget{smoke-test-dataset}{%
\paragraph{2.1 Smoke-test dataset}\label{smoke-test-dataset}}

Purpose:

\begin{itemize}
\tightlist
\item
  verify code,
\item
  verify dataset schema,
\item
  verify normalization,
\item
  verify training loop,
\item
  verify notebook execution.
\end{itemize}

Recommended size:

\[
N = 500\text{--}2{,}000
\]

Do not use this for scientific conclusions.

\begin{center}\rule{0.5\linewidth}{0.5pt}\end{center}

\hypertarget{first-real-model}{%
\paragraph{2.2 First real model}\label{first-real-model}}

Purpose:

\begin{itemize}
\tightlist
\item
  determine whether the input/output definition is reasonable,
\item
  compare MLP and physics-regularized MLP,
\item
  inspect error distributions.
\end{itemize}

Recommended size:

\[
N = 10{,}000\text{--}30{,}000
\]

This is suitable for first useful model comparisons.

\begin{center}\rule{0.5\linewidth}{0.5pt}\end{center}

\hypertarget{reduced-input-optics-surrogate}{%
\paragraph{2.3 Reduced-input optics
surrogate}\label{reduced-input-optics-surrogate}}

Purpose:

\begin{itemize}
\tightlist
\item
  rms-size and beam-state propagation,
\item
  Twiss/emittance dependence,
\item
  phase-space reconstruction support.
\end{itemize}

Approximate input dimension:

\[
6 + N_{\mathrm{equad}}
\]

or

\[
7 + N_{\mathrm{equad}}
\]

if current is included.

Recommended size:

\[
N = 50{,}000\text{--}100{,}000
\]

\begin{center}\rule{0.5\linewidth}{0.5pt}\end{center}

\hypertarget{orbit-response-surrogate}{%
\paragraph{2.4 Orbit-response
surrogate}\label{orbit-response-surrogate}}

Purpose:

\begin{itemize}
\tightlist
\item
  orbit correction,
\item
  steering response,
\item
  experimental RL pretraining.
\end{itemize}

Approximate input dimension:

\[
N_{\mathrm{steerer}} + N_{\mathrm{equad}}
\]

Recommended size:

\[
N = 10{,}000\text{--}50{,}000
\]

If the response is close to linear, the lower end may be sufficient.

\begin{center}\rule{0.5\linewidth}{0.5pt}\end{center}

\hypertarget{full-24d-surrogate}{%
\paragraph{2.5 Full 24D surrogate}\label{full-24d-surrogate}}

Purpose:

\begin{itemize}
\tightlist
\item
  absolute diagnostic prediction over a broad operating range.
\end{itemize}

Recommended size:

\[
N = 100{,}000\text{--}300{,}000
\]

If the model includes wide current variation, large offsets, strong
nonlinearities, and beam-state outputs, use more data.

\begin{center}\rule{0.5\linewidth}{0.5pt}\end{center}

\hypertarget{full-model-with-diagnostic-and-state-outputs}{%
\paragraph{2.6 Full model with diagnostic and state
outputs}\label{full-model-with-diagnostic-and-state-outputs}}

Purpose:

\begin{itemize}
\tightlist
\item
  publication-quality full surrogate,
\item
  diagnostic prediction,
\item
  final-state prediction,
\item
  uncertainty analysis.
\end{itemize}

Recommended size:

\[
N = 200{,}000\text{--}500{,}000
\]

Use this only after reduced models are validated.

\begin{center}\rule{0.5\linewidth}{0.5pt}\end{center}

\hypertarget{why-the-full-24d-model-needs-many-samples}{%
\subsubsection{3. Why the full 24D model needs many
samples}\label{why-the-full-24d-model-needs-many-samples}}

A 24D input space is large. Uniformly covering it is inefficient.

The required sample size increases when:

\begin{itemize}
\tightlist
\item
  many inputs are weakly relevant but included anyway,
\item
  output includes both centroids and beam-state moments,
\item
  current-dependent dynamics are included,
\item
  large offset/angle ranges cause nonlinear or aperture effects,
\item
  the training region includes both nominal and stress-test conditions.
\end{itemize}

Therefore, reducing the input dimension by model purpose is strongly
recommended.

\begin{center}\rule{0.5\linewidth}{0.5pt}\end{center}

\hypertarget{active-learning-strategy}{%
\subsubsection{4. Active-learning
strategy}\label{active-learning-strategy}}

Use iterative data generation instead of one large blind dataset.

Recommended loop:

\begin{enumerate}
\def\labelenumi{\arabic{enumi}.}
\tightlist
\item
  generate \(N=5{,}000\) samples;
\item
  train a quick MLP;
\item
  evaluate error versus input parameters;
\item
  identify high-error regions;
\item
  add \(N=20{,}000\) Sobol samples;
\item
  retrain;
\item
  add focused samples near high-error or high-gradient regions;
\item
  repeat until validation error saturates.
\end{enumerate}

This is more efficient than generating hundreds of thousands of samples
immediately.

\begin{center}\rule{0.5\linewidth}{0.5pt}\end{center}

\hypertarget{learning-curve-requirement}{%
\subsubsection{5. Learning-curve
requirement}\label{learning-curve-requirement}}

Codex should implement tools to generate learning curves.

Train the same model with:

\[
N =
[
1{,}000,
3{,}000,
10{,}000,
30{,}000,
100{,}000
]
\]

when available.

Plot:

\[
\mathrm{MAE}(N),
\qquad
\mathrm{RMSE}(N),
\qquad
R^2(N),
\qquad
\mathrm{physics\ violation}(N).
\]

This determines whether more data are still useful.

\begin{center}\rule{0.5\linewidth}{0.5pt}\end{center}

\hypertarget{validation-split-strategy}{%
\subsubsection{6. Validation split
strategy}\label{validation-split-strategy}}

Do not use only random splits.

Use:

\begin{enumerate}
\def\labelenumi{\arabic{enumi}.}
\tightlist
\item
  random interpolation split;
\item
  held-out Twiss/emittance region;
\item
  held-out equad range;
\item
  held-out steering range;
\item
  stress-test region;
\item
  closed-loop correction trajectory validation.
\end{enumerate}

For each split, report separate metrics.

\begin{center}\rule{0.5\linewidth}{0.5pt}\end{center}

\hypertarget{suggested-dataset-phases}{%
\subsubsection{7. Suggested dataset
phases}\label{suggested-dataset-phases}}

\hypertarget{phase-a-smoke-dataset}{%
\paragraph{Phase A: smoke dataset}\label{phase-a-smoke-dataset}}

\[
N=1{,}000
\]

Use in CI tests and notebooks.

\hypertarget{phase-b-pilot-dataset}{%
\paragraph{Phase B: pilot dataset}\label{phase-b-pilot-dataset}}

\[
N=10{,}000\text{--}30{,}000
\]

Use for first MLP and physics-regularized comparison.

\hypertarget{phase-c-reduced-production-dataset}{%
\paragraph{Phase C: reduced production
dataset}\label{phase-c-reduced-production-dataset}}

\[
N=50{,}000\text{--}100{,}000
\]

Use for optics and orbit-response models.

\hypertarget{phase-d-full-production-dataset}{%
\paragraph{Phase D: full production
dataset}\label{phase-d-full-production-dataset}}

\[
N=100{,}000\text{--}300{,}000
\]

Use only after the target schema and model split are stable.

\hypertarget{phase-e-active-learning-refinement}{%
\paragraph{Phase E: active-learning
refinement}\label{phase-e-active-learning-refinement}}

Add focused samples in high-error regions.

\begin{center}\rule{0.5\linewidth}{0.5pt}\end{center}

\hypertarget{required-tests-and-utilities}{%
\subsubsection{8. Required tests and
utilities}\label{required-tests-and-utilities}}

Codex should implement:

\begin{enumerate}
\def\labelenumi{\arabic{enumi}.}
\tightlist
\item
  configurable dataset-size presets;
\item
  smoke-test dataset generator;
\item
  learning-curve training script;
\item
  grouped validation split utilities;
\item
  error-by-region reporting;
\item
  reproducibility check for fixed seed;
\item
  runtime metadata logging.
\end{enumerate}

\begin{center}\rule{0.5\linewidth}{0.5pt}\end{center}

\hypertarget{codex-prompt-3}{%
\subsubsection{9. Codex prompt}\label{codex-prompt-3}}

Use the following prompt for Codex:

\begin{Shaded}
\begin{Highlighting}[]
\NormalTok{Read AGENTS.md, CODEX\_TASKS\_2026.md, and this document:}
\NormalTok{docs/codex\_tasks/04\_training\_sample\_requirements.md.}

\NormalTok{Task:}
\NormalTok{Add dataset{-}size presets, learning{-}curve utilities, and validation split}
\NormalTok{support for LEBT surrogate model development.}

\NormalTok{Requirements:}
\NormalTok{1. Support dataset{-}size presets:}
\NormalTok{   {-} smoke: 500{-}2,000 samples}
\NormalTok{   {-} pilot: 10,000{-}30,000 samples}
\NormalTok{   {-} reduced\_production: 50,000{-}100,000 samples}
\NormalTok{   {-} full\_production: 100,000{-}300,000 samples}
\NormalTok{2. Add grouped validation split utilities:}
\NormalTok{   {-} random}
\NormalTok{   {-} held{-}out Twiss/emittance region}
\NormalTok{   {-} held{-}out equad range}
\NormalTok{   {-} held{-}out steering range}
\NormalTok{   {-} stress{-}test region}
\NormalTok{3. Add a learning{-}curve configuration that can train on increasing subsets.}
\NormalTok{4. Add metrics grouped by output target and validation region.}
\NormalTok{5. Ensure the smoke preset is small enough for CI or quick notebook execution.}
\NormalTok{6. Do not generate a large dataset by default.}
\NormalTok{7. Do not implement MEBT or RL in this task.}

\NormalTok{After implementation, summarize changed files, tests run, and remaining assumptions.}
\end{Highlighting}
\end{Shaded}

\hypertarget{physics-models-and-machine-conventions}{%
\section{Physics Models and Machine
Conventions}\label{physics-models-and-machine-conventions}}

Source Markdown:
\texttt{docs/LEBT\_STEERING\_CURRENT\_KICK\_CONVERSION.md}

\hypertarget{lebt-steering-current-to-kick-conversion}{%
\subsection{LEBT Steering Current to Kick
Conversion}\label{lebt-steering-current-to-kick-conversion}}

\hypertarget{current-internal-convention}{%
\subsubsection{Current Internal
Convention}\label{current-internal-convention}}

The existing synapticTrack LEBT simulation, no-offset Sobol dataset,
validation workflow, and MLP training pipeline use steering magnet kick
angles in \texttt{mrad}.

The existing 19D sampled input remains:

\begin{Shaded}
\begin{Highlighting}[]
\NormalTok{alpha\_x, beta\_x, emittance\_x,}
\NormalTok{alpha\_y, beta\_y, emittance\_y,}
\NormalTok{SM0\_H, SM1\_H, SM2\_H, SM3\_H, SM4\_H,}
\NormalTok{SM0\_V, SM1\_V, SM2\_V, SM3\_V, SM4\_V,}
\NormalTok{alpha\_0, alpha\_1, alpha\_2}
\end{Highlighting}
\end{Shaded}

The steering entries \texttt{SM0\_H\ ...\ SM4\_V} are kick values in
\texttt{mrad}.

\hypertarget{machine-current-convention}{%
\subsubsection{Machine Current
Convention}\label{machine-current-convention}}

The real RAON LEBT steering magnets are controlled by power-supply
current in amperes. The current assumed linear calibration is:

\begin{Shaded}
\begin{Highlighting}[]
\NormalTok{1 mrad kick = 0.168 A}
\end{Highlighting}
\end{Shaded}

The nominal power-supply current step is:

\begin{Shaded}
\begin{Highlighting}[]
\NormalTok{0.01 A}
\end{Highlighting}
\end{Shaded}

Therefore:

\begin{Shaded}
\begin{Highlighting}[]
\NormalTok{kick\_mrad = current\_A / 0.168}
\NormalTok{current\_A = 0.168 * kick\_mrad}
\NormalTok{0.01 A = 0.0595238 mrad}
\end{Highlighting}
\end{Shaded}

\hypertarget{conversion-equations}{%
\subsubsection{Conversion Equations}\label{conversion-equations}}

The implemented sign and offset convention is:

\begin{Shaded}
\begin{Highlighting}[]
\NormalTok{kick\_mrad = polarity * (current\_A {-} current\_offset\_A) / amp\_per\_mrad}
\NormalTok{current\_A = current\_offset\_A + polarity * amp\_per\_mrad * kick\_mrad}
\end{Highlighting}
\end{Shaded}

The default calibration is:

\begin{Shaded}
\begin{Highlighting}[]
\NormalTok{amp\_per\_mrad = 0.168}
\NormalTok{current\_step\_A = 0.01}
\NormalTok{current\_offset\_A = 0.0}
\NormalTok{polarity = +1}
\end{Highlighting}
\end{Shaded}

Each steering magnet can override:

\begin{Shaded}
\begin{Highlighting}[]
\NormalTok{amp\_per\_mrad}
\NormalTok{current\_offset\_A}
\NormalTok{current\_min\_A}
\NormalTok{current\_max\_A}
\NormalTok{current\_step\_A}
\NormalTok{polarity}
\end{Highlighting}
\end{Shaded}

The optional config file is:

\begin{Shaded}
\begin{Highlighting}[]
\NormalTok{configs/hardware/lebt\_steering\_calibration.yaml}
\end{Highlighting}
\end{Shaded}

If that file is missing, built-in defaults are used.

\hypertarget{quantization}{%
\subsubsection{Quantization}\label{quantization}}

Current quantization is:

\begin{Shaded}
\begin{Highlighting}[]
\NormalTok{I\_quantized = round(I / current\_step\_A) * current\_step\_A}
\end{Highlighting}
\end{Shaded}

Quantized kick is computed by converting the quantized current:

\begin{Shaded}
\begin{Highlighting}[]
\NormalTok{kick\_quantized\_mrad = I\_quantized / amp\_per\_mrad}
\end{Highlighting}
\end{Shaded}

With the default calibration:

\begin{Shaded}
\begin{Highlighting}[]
\NormalTok{kick\_step\_mrad = 0.01 / 0.168 = 0.0595238 mrad}
\end{Highlighting}
\end{Shaded}

Quantization is opt-in. Simulation-only dataset generation and MLP
training should normally leave quantization disabled.

\hypertarget{policy-options}{%
\subsubsection{Policy Options}\label{policy-options}}

\hypertarget{option-a-simulation-native}{%
\paragraph{Option A:
Simulation-Native}\label{option-a-simulation-native}}

Training and simulation use steering kick in \texttt{mrad}. Machine
current is converted only at the hardware boundary.

This is the current default and the recommended policy for the existing
no-offset Sobol dataset and MLP surrogate training.

\hypertarget{option-b-hardware-native}{%
\paragraph{Option B: Hardware-Native}\label{option-b-hardware-native}}

ML or RL actions use power-supply current in \texttt{A}. Before
simulation or surrogate evaluation, currents are converted to kicks in
\texttt{mrad}.

This is recommended for future real-machine RL deployment because the
action variable matches the physical actuator.

\hypertarget{option-c-dual-representation}{%
\paragraph{Option C: Dual
Representation}\label{option-c-dual-representation}}

Datasets may explicitly store both:

\begin{Shaded}
\begin{Highlighting}[]
\NormalTok{steering\_kick\_mrad}
\NormalTok{steering\_current\_A}
\end{Highlighting}
\end{Shaded}

Training can then choose a representation through an explicit opt-in
input view. This is useful for future experiment/simulation alignment,
but it is not enabled by default.

\hypertarget{current-training-recommendation}{%
\subsubsection{Current Training
Recommendation}\label{current-training-recommendation}}

Keep the current production dataset and MLP training in kick units. The
simulation naturally accepts mrad kicks, and the existing 19D Sobol
dataset is already defined that way.

Use:

\begin{Shaded}
\begin{Highlighting}[]
\ImportTok{from}\NormalTok{ synapticTrack.hardware }\ImportTok{import}\NormalTok{ SteeringActuatorAdapter}

\NormalTok{adapter }\OperatorTok{=}\NormalTok{ SteeringActuatorAdapter()}
\NormalTok{kick\_dict }\OperatorTok{=}\NormalTok{ adapter.currents\_to\_kicks(current\_dict, quantize}\OperatorTok{=}\VariableTok{True}\NormalTok{)}
\end{Highlighting}
\end{Shaded}

only when translating machine-facing current settings to
simulation-facing kicks.

\hypertarget{future-real-machine-rl-recommendation}{%
\subsubsection{Future Real-Machine RL
Recommendation}\label{future-real-machine-rl-recommendation}}

Use current in \texttt{A} as the RL action variable at the hardware
boundary. Convert to \texttt{mrad} before simulation or surrogate
evaluation:

\begin{Shaded}
\begin{Highlighting}[]
\NormalTok{kick\_dict }\OperatorTok{=}\NormalTok{ adapter.currents\_to\_kicks(current\_dict, quantize}\OperatorTok{=}\VariableTok{True}\NormalTok{)}
\end{Highlighting}
\end{Shaded}

This keeps the policy action space aligned with the real power supplies
while preserving synapticTrack's internal mrad convention.

\hypertarget{limitations}{%
\subsubsection{Limitations}\label{limitations}}

The current implementation assumes a linear current-to-kick calibration.
Future machine studies should add:

\begin{itemize}
\tightlist
\item
  beam-based calibration per magnet;
\item
  hysteresis characterization;
\item
  polarity verification;
\item
  current limits from hardware;
\item
  reproducibility checks after cycling magnets;
\item
  time-dependent or nonlinear correction terms if needed.
\end{itemize}

Those effects are deliberately not assumed in the current layer.

Source Markdown: \texttt{docs/lebt\_production\_dataset\_configs.md}

\hypertarget{lebt-production-dataset-configurations}{%
\subsection{LEBT Production Dataset
Configurations}\label{lebt-production-dataset-configurations}}

This document records the first versioned production dataset-generation
configurations for the Task-3 LEBT surrogate workflow. These configs are
intentionally staged and do not launch full production generation by
default.

\hypertarget{current-status-relative-to-sobol-workflow}{%
\subsubsection{Current Status Relative to Sobol
Workflow}\label{current-status-relative-to-sobol-workflow}}

These early Task-3 production configs remain useful as design records
and for legacy/generic LEBT dataset generation. The maintained no-offset
production workflow is now
\texttt{lebt\_production\_no\_offset\_multitarget\_sobol\_v1},
documented in \texttt{docs/lebt\_no\_offset\_sobol\_workflow.md}, with a
19D sampled input and 44D target output.

\hypertarget{pilot-audit-decision}{%
\subsubsection{Pilot Audit Decision}\label{pilot-audit-decision}}

The latest pilot audit did not find a pilot-sized current-schema dataset
in the workspace. The only current-schema TRACK dataset had 8 samples,
and the larger available files were legacy orbit-correction datasets
without schema metadata, validity masks, units, failed-sample reasons,
or beam-state labels.

Decision for v1 production configs:

\begin{itemize}
\tightlist
\item
  keep the Task-3 \texttt{training} ranges instead of widening ranges;
\item
  split production datasets by model type;
\item
  require a 100-sample preflight for each config before full generation;
\item
  include metadata that records the pilot-audit blocker.
\end{itemize}

\hypertarget{versioned-configs}{%
\subsubsection{Versioned Configs}\label{versioned-configs}}

\begin{longtable}[]{@{}
  >{\raggedright\arraybackslash}p{(\columnwidth - 10\tabcolsep) * \real{0.1429}}
  >{\raggedleft\arraybackslash}p{(\columnwidth - 10\tabcolsep) * \real{0.1905}}
  >{\raggedleft\arraybackslash}p{(\columnwidth - 10\tabcolsep) * \real{0.1905}}
  >{\raggedright\arraybackslash}p{(\columnwidth - 10\tabcolsep) * \real{0.1429}}
  >{\raggedleft\arraybackslash}p{(\columnwidth - 10\tabcolsep) * \real{0.1905}}
  >{\raggedright\arraybackslash}p{(\columnwidth - 10\tabcolsep) * \real{0.1429}}@{}}
\toprule\noalign{}
\begin{minipage}[b]{\linewidth}\raggedright
Config ID
\end{minipage} & \begin{minipage}[b]{\linewidth}\raggedleft
Purpose
\end{minipage} & \begin{minipage}[b]{\linewidth}\raggedleft
Samples
\end{minipage} & \begin{minipage}[b]{\linewidth}\raggedright
Preset
\end{minipage} & \begin{minipage}[b]{\linewidth}\raggedleft
State Outputs
\end{minipage} & \begin{minipage}[b]{\linewidth}\raggedright
Default Output
\end{minipage} \\
\midrule\noalign{}
\endhead
\bottomrule\noalign{}
\endlastfoot
\texttt{lebt\_production\_diag\_v1} & full-input diagnostic target &
50,000 & \texttt{full} & no &
\texttt{data/production/lebt\_production\_diag\_v1.h5} \\
\texttt{lebt\_production\_optics\_state\_v1} & optics/state target &
50,000 & \texttt{optics} & yes &
\texttt{data/production/lebt\_production\_optics\_state\_v1.h5} \\
\texttt{lebt\_production\_orbit\_response\_v1} & steerer/equad orbit
response & 20,000 & \texttt{orbit\_response} & no &
\texttt{data/production/lebt\_production\_orbit\_response\_v1.h5} \\
\end{longtable}

Config files live in \texttt{configs/lebt/}.

\hypertarget{metadata-saved}{%
\subsubsection{Metadata Saved}\label{metadata-saved}}

The config loader injects these fields into dataset metadata:

\begin{itemize}
\tightlist
\item
  \texttt{production\_config\_id}
\item
  \texttt{production\_config\_version}
\item
  \texttt{git\_commit\_hash}
\item
  \texttt{lattice\_version}
\item
  \texttt{backend\_name}
\item
  \texttt{units}
\item
  \texttt{random\_seed}
\item
  \texttt{sampling\_method}
\item
  \texttt{parameter\_ranges}
\item
  \texttt{generation\_timestamp}
\item
  \texttt{pilot\_audit\_status}
\end{itemize}

The full resolved \texttt{LEBTDatasetConfig} is also saved in generated
HDF5 \texttt{config\_json} by the dataset generator. Input-only dry runs
save the same resolved config metadata.

\hypertarget{dry-run-input-tables}{%
\subsubsection{Dry-Run Input Tables}\label{dry-run-input-tables}}

Dry-run mode for production configs writes only sampled inputs and
summary statistics. It does not create target arrays and does not call
TRACK.

\begin{Shaded}
\begin{Highlighting}[]
\ExtensionTok{python3} \AttributeTok{{-}m}\NormalTok{ synapticTrack.cli lebt{-}production{-}dry{-}run lebt\_production\_diag\_v1}
\ExtensionTok{python3} \AttributeTok{{-}m}\NormalTok{ synapticTrack.cli lebt{-}production{-}dry{-}run lebt\_production\_optics\_state\_v1}
\ExtensionTok{python3} \AttributeTok{{-}m}\NormalTok{ synapticTrack.cli lebt{-}production{-}dry{-}run lebt\_production\_orbit\_response\_v1}
\end{Highlighting}
\end{Shaded}

Optional overrides:

\begin{Shaded}
\begin{Highlighting}[]
\ExtensionTok{python3} \AttributeTok{{-}m}\NormalTok{ synapticTrack.cli lebt{-}production{-}dry{-}run lebt\_production\_diag\_v1 }\DataTypeTok{\textbackslash{}}
  \AttributeTok{{-}{-}n{-}samples}\NormalTok{ 1000 }\DataTypeTok{\textbackslash{}}
  \AttributeTok{{-}{-}output{-}file}\NormalTok{ data/dry\_run/lebt\_production\_diag\_v1\_inputs\_1000.h5 }\DataTypeTok{\textbackslash{}}
  \AttributeTok{{-}{-}summary{-}file}\NormalTok{ data/dry\_run/lebt\_production\_diag\_v1\_inputs\_1000.csv}
\end{Highlighting}
\end{Shaded}

\hypertarget{preflight-commands}{%
\subsubsection{Preflight Commands}\label{preflight-commands}}

These commands print a 100-sample preflight command by default:

\begin{Shaded}
\begin{Highlighting}[]
\ExtensionTok{python3} \AttributeTok{{-}m}\NormalTok{ synapticTrack.cli lebt{-}production{-}preflight lebt\_production\_diag\_v1}
\ExtensionTok{python3} \AttributeTok{{-}m}\NormalTok{ synapticTrack.cli lebt{-}production{-}preflight lebt\_production\_optics\_state\_v1}
\ExtensionTok{python3} \AttributeTok{{-}m}\NormalTok{ synapticTrack.cli lebt{-}production{-}preflight lebt\_production\_orbit\_response\_v1}
\end{Highlighting}
\end{Shaded}

To actually run the preflight after reviewing dry-run inputs and TRACK
availability, add \texttt{-\/-execute}:

\begin{Shaded}
\begin{Highlighting}[]
\ExtensionTok{python3} \AttributeTok{{-}m}\NormalTok{ synapticTrack.cli lebt{-}production{-}preflight lebt\_production\_diag\_v1 }\AttributeTok{{-}{-}execute}
\ExtensionTok{python3} \AttributeTok{{-}m}\NormalTok{ synapticTrack.cli lebt{-}production{-}preflight lebt\_production\_optics\_state\_v1 }\AttributeTok{{-}{-}execute}
\ExtensionTok{python3} \AttributeTok{{-}m}\NormalTok{ synapticTrack.cli lebt{-}production{-}preflight lebt\_production\_orbit\_response\_v1 }\AttributeTok{{-}{-}execute}
\end{Highlighting}
\end{Shaded}

Do not run full production generation until each preflight has been
audited for validity fraction, failed/lost rows, output ranges, units,
and beam-state moment validity where applicable.

\hypertarget{full-production-commands}{%
\subsubsection{Full Production
Commands}\label{full-production-commands}}

Full generation is intentionally not wrapped as a one-line default
command in this task. After successful preflight review, run production
through the config loader or a reviewed operations script that calls
\texttt{generate\_lebt\_dataset(config.to\_lebt\_config(dry\_run=False))}
for the selected config.

\hypertarget{simulation-code-structure-and-external-codes}{%
\section{Simulation Code Structure and External
Codes}\label{simulation-code-structure-and-external-codes}}

Source Markdown: \texttt{docs/TRACK\_SETUP.md}

\hypertarget{track-setup}{%
\subsection{TRACK Setup}\label{track-setup}}

\texttt{synapticTrack} does not commit or vendor the external TRACK
executable. The repository-local convention is:

\begin{Shaded}
\begin{Highlighting}[]
\NormalTok{synapticTrack/}
\NormalTok{+{-}{-} TRACK/}
\NormalTok{|   +{-}{-} TRACKv39C.exe}
\NormalTok{+{-}{-} bin/}
\NormalTok{|   +{-}{-} helper commands}
\end{Highlighting}
\end{Shaded}

\texttt{TRACK/} is local-only and ignored by Git. Create it yourself:

\begin{Shaded}
\begin{Highlighting}[]
\FunctionTok{mkdir} \AttributeTok{{-}p}\NormalTok{ TRACK}
\CommentTok{\#\# place TRACKv39C.exe and any local runtime files under TRACK/}
\FunctionTok{chmod}\NormalTok{ +x TRACK/TRACKv39C.exe}
\end{Highlighting}
\end{Shaded}

The default resolved executable is:

\begin{Shaded}
\begin{Highlighting}[]
\NormalTok{TRACK/TRACKv39C.exe}
\end{Highlighting}
\end{Shaded}

Override it with an environment variable:

\begin{Shaded}
\begin{Highlighting}[]
\BuiltInTok{export} \VariableTok{SYNAPTICTRACK\_TRACK\_EXE}\OperatorTok{=}\NormalTok{/custom/path/to/TRACKv39C.exe}
\ExtensionTok{bin/syntrack{-}track{-}check}
\end{Highlighting}
\end{Shaded}

Or use a config override:

\begin{Shaded}
\begin{Highlighting}[]
\FunctionTok{track\_exe\_path}\KeywordTok{:}\AttributeTok{ /custom/path/to/TRACKv39C.exe}
\end{Highlighting}
\end{Shaded}

Resolution precedence is:

\begin{enumerate}
\def\labelenumi{\arabic{enumi}.}
\tightlist
\item
  explicit config or command-line path;
\item
  \texttt{SYNAPTICTRACK\_TRACK\_EXE};
\item
  repository-local \texttt{TRACK/TRACKv39C.exe}.
\end{enumerate}

\hypertarget{helper-commands}{%
\subsubsection{Helper Commands}\label{helper-commands}}

Check TRACK path:

\begin{Shaded}
\begin{Highlighting}[]
\ExtensionTok{bin/syntrack{-}track{-}check}
\ExtensionTok{bin/syntrack{-}track{-}check} \AttributeTok{{-}{-}track{-}exe}\NormalTok{ TRACK/TRACKv39C.exe}
\end{Highlighting}
\end{Shaded}

Inspect preflight settings without running production:

\begin{Shaded}
\begin{Highlighting}[]
\ExtensionTok{bin/syntrack{-}lebt{-}preflight}
\end{Highlighting}
\end{Shaded}

Run only the 240-sample preflight:

\begin{Shaded}
\begin{Highlighting}[]
\ExtensionTok{bin/syntrack{-}lebt{-}preflight} \AttributeTok{{-}{-}run}
\end{Highlighting}
\end{Shaded}

Validate an existing no-offset Sobol dataset:

\begin{Shaded}
\begin{Highlighting}[]
\ExtensionTok{bin/syntrack{-}lebt{-}validate}
\ExtensionTok{bin/syntrack{-}lebt{-}validate}\NormalTok{ data/processed/lebt\_production\_no\_offset\_multitarget\_sobol\_v1 }\AttributeTok{{-}{-}write{-}reports}
\end{Highlighting}
\end{Shaded}

Installed console scripts expose the same commands when the package is
installed:

\begin{Shaded}
\begin{Highlighting}[]
\ExtensionTok{syntrack{-}track{-}check}
\ExtensionTok{syntrack{-}lebt{-}preflight}
\ExtensionTok{syntrack{-}lebt{-}validate}
\end{Highlighting}
\end{Shaded}

\hypertarget{segmented-track-runs}{%
\subsubsection{Segmented TRACK Runs}\label{segmented-track-runs}}

RAON MEBT notebooks use segmented TRACK execution. The initial beam
distribution number must match the source lattice convention, and each
segment must write the next scratch distribution number for the
following segment. For example, an MEBT run starting from
\texttt{scratch.\#07} writes \texttt{scratch.\#08} after the first
segment, then the next segment reads \texttt{scratch.\#08}. This
preserves full-lattice ordering while exposing intermediate diagnostic
outputs.

\hypertarget{production-notebooks}{%
\subsubsection{Production Notebooks}\label{production-notebooks}}

Production notebooks display the resolved TRACK path and whether it
exists before generation. They should not contain machine-specific TRACK
paths in committed source or outputs.

\hypertarget{beam-current-unit-contract}{%
\subsubsection{Beam Current Unit
Contract}\label{beam-current-unit-contract}}

The repository does not contain an authoritative TRACK manual statement
for the \texttt{current} namelist unit. \texttt{synapticTrack} therefore
uses \texttt{uA} as its canonical beam-current unit and exposes
\texttt{track\_current\_per\_synaptic\_current} as an explicit
conversion boundary. Its current default is \texttt{1.0}, solely to
preserve existing \texttt{track.dat} values such as
\texttt{current\ =\ 0.05}; it is marked
\texttt{unverified\_local\_repository} and must not be interpreted as a
verified physical unit equivalence. Generated LEBT/MEBT metadata records
the canonical value, the TRACK input value, the conversion factor, and
the verification status.

Source Markdown: \texttt{docs/refactor\_baseline\_r0.md}

\hypertarget{r0-baseline-characterization-production-dataset-pipeline}{%
\subsection{R0 Baseline Characterization: Production Dataset
Pipeline}\label{r0-baseline-characterization-production-dataset-pipeline}}

Date: 2026-07-15

Branch inspected: \texttt{devel}

Baseline note: this report was created after the local R1 and R2 commits
had already been applied:

\begin{itemize}
\tightlist
\item
  \texttt{e06fd2e} --- R1 schema centralization
\item
  \texttt{66c720e} --- R2 production config policy and validation
\end{itemize}

The behavioral baseline below records the current branch state at report
time. No production generation, neural-network training, or tracking
backend execution was performed for this report.

\hypertarget{scope-inspected}{%
\subsubsection{Scope Inspected}\label{scope-inspected}}

Required files inspected:

\begin{itemize}
\tightlist
\item
  \texttt{src/synapticTrack/simulation/no\_offset\_sobol\_production.py}
\item
  \texttt{src/synapticTrack/simulation/dataset\_generator.py}
\item
  \texttt{src/synapticTrack/data/input\_views.py}
\item
  \texttt{src/synapticTrack/data/target\_views.py}
\item
  \texttt{src/synapticTrack/simulation/no\_offset\_sobol\_validation.py}
\item
  \texttt{notebooks/04\_lebt\_no\_offset\_sobol\_dataset\_production.ipynb}
\item
  \texttt{notebooks/05\_lebt\_no\_offset\_sobol\_dataset\_validation.ipynb}
\end{itemize}

Task documents inspected:

\begin{itemize}
\tightlist
\item
  \texttt{AGENTS.md}
\item
  \texttt{CODEX\_TASKS\_2026.md}
\item
  \texttt{docs/codex\_tasks/01\_input\_dimension\_strategy.md}
\item
  \texttt{docs/codex\_tasks/02\_output\_storage\_strategy.md}
\item
  \texttt{docs/codex\_tasks/03\_input\_parameter\_ranges.md}
\item
  \texttt{docs/codex\_tasks/04\_training\_sample\_requirements.md}
\end{itemize}

\hypertarget{current-expected-behavior}{%
\subsubsection{Current Expected
Behavior}\label{current-expected-behavior}}

The no-offset Sobol LEBT production workflow currently targets:

\begin{itemize}
\tightlist
\item
  Dataset name:
  \texttt{lebt\_production\_no\_offset\_multitarget\_sobol\_v1}
\item
  Sampled input dimension: 19D
\item
  Fixed-zero parameters: 5

  \begin{itemize}
  \tightlist
  \item
    \texttt{x0}
  \item
    \texttt{xp0}
  \item
    \texttt{y0}
  \item
    \texttt{yp0}
  \item
    \texttt{beam\_current}
  \end{itemize}
\item
  Effective full simulation parameter dimension: 24D
\item
  Target groups:

  \begin{itemize}
  \tightlist
  \item
    \texttt{diagnostics}: 24D
  \item
    \texttt{state\_last\_ws}: 10D
  \item
    \texttt{state\_lebt\_exit}: 10D
  \end{itemize}
\item
  Total target dimension: 44D
\item
  Target view for multi-task production:
  \texttt{diag\_plus\_all\_states}
\item
  Input view for no-offset production:
  \texttt{full\_19d\_fixed\_injection}
\item
  Sobol continuation stages:

  \begin{itemize}
  \tightlist
  \item
    200,000 samples
  \item
    300,000 samples
  \item
    500,000 samples
  \end{itemize}
\item
  Batch-size default for production config: 2,400
\item
  Preflight config: 240 samples, batch size 120
\end{itemize}

Sobol continuation is based on global sample indices. Existing completed
indices are discovered from batch files, duplicates are rejected, and
missing contiguous index ranges are planned as new batches.

\hypertarget{current-modules-and-responsibilities}{%
\subsubsection{Current Modules And
Responsibilities}\label{current-modules-and-responsibilities}}

\hypertarget{no_offset_sobol_production.py}{%
\paragraph{\texorpdfstring{\texttt{no\_offset\_sobol\_production.py}}{no\_offset\_sobol\_production.py}}\label{no_offset_sobol_production.py}}

Primary responsibility: specialized no-offset Sobol LEBT production
orchestration.

Current responsibilities include:

\begin{itemize}
\tightlist
\item
  \texttt{SobolConfig} and \texttt{NoOffsetSobolProductionConfig}
\item
  production/preflight config loading and validation
\item
  Sobol unit-cube sampling by global index
\item
  physical input mapping
\item
  19D to full 24D expansion with fixed-zero injection parameters
\item
  duplicate global-index detection
\item
  batch planning and resume/extension behavior
\item
  preflight configuration conversion
\item
  batch execution through \texttt{SimulationBackend}
\item
  per-sample work-dir policy and progress printing
\item
  HDF5 batch writing
\item
  config, metadata, generation log, Sobol index, audit, and
  validity-summary writing
\item
  runtime estimation from preflight output
\end{itemize}

The module is still broad: about 988 lines at inspection time.

\hypertarget{dataset_generator.py}{%
\paragraph{\texorpdfstring{\texttt{dataset\_generator.py}}{dataset\_generator.py}}\label{dataset_generator.py}}

Primary responsibility: general LEBT dataset generation path used by
smoke/pilot/test workflows.

Current responsibilities include:

\begin{itemize}
\tightlist
\item
  LEBT parameter range objects
\item
  dataset presets and range profile resolution
\item
  beamline-neutral \texttt{DatasetGenerationConfig}
\item
  LEBT-specific \texttt{LEBTDatasetConfig}
\item
  dry-run sampling validation and summary tables
\item
  deterministic seed utilities
\item
  Sobol/LHS input generation
\item
  generic calibrated LEBT dataset generation
\item
  input view materialization
\item
  target group storage for diagnostics, orbit response, and state
  outputs
\item
  beam-state moment/Twiss conversion and validation
\item
  dataset audit support
\end{itemize}

The module is about 1,473 lines and overlaps with the production
generator in config validation, worker execution, deterministic
sampling, HDF5 writing, progress handling, cleanup policy, masks,
failure reasons, and metadata.

\hypertarget{input_views.py}{%
\paragraph{\texorpdfstring{\texttt{input\_views.py}}{input\_views.py}}\label{input_views.py}}

Primary responsibility: deterministic input-column selection.

Public behavior:

\begin{itemize}
\tightlist
\item
  \texttt{full}: all available input columns
\item
  \texttt{optics}: Twiss, emittance, optional current, and equads
\item
  \texttt{orbit\_response}: steerers plus equads
\item
  \texttt{custom}: user-supplied ordered column list
\end{itemize}

The helper validates missing/duplicated columns and preserves
deterministic ordering.

\hypertarget{target_views.py}{%
\paragraph{\texorpdfstring{\texttt{target\_views.py}}{target\_views.py}}\label{target_views.py}}

Primary responsibility: target-view loading from HDF5 datasets.

Public behavior:

\begin{itemize}
\tightlist
\item
  \texttt{diagnostics}
\item
  \texttt{orbit\_response}
\item
  \texttt{state\_last\_ws}
\item
  \texttt{state\_lebt\_exit}
\item
  \texttt{diag\_plus\_exit\_state}
\item
  \texttt{diag\_plus\_all\_states}
\item
  \texttt{custom}
\end{itemize}

It supports current grouped output paths and legacy flat datasets such
as \texttt{y\_diag}, \texttt{y\_state\_last\_ws}, and
\texttt{y\_state\_exit}.

\hypertarget{no_offset_sobol_validation.py}{%
\paragraph{\texorpdfstring{\texttt{no\_offset\_sobol\_validation.py}}{no\_offset\_sobol\_validation.py}}\label{no_offset_sobol_validation.py}}

Primary responsibility: validate no-offset Sobol production datasets and
write validation reports.

Current responsibilities include:

\begin{itemize}
\tightlist
\item
  required-file checks
\item
  metadata/config loading
\item
  Sobol index uniqueness and continuity checks
\item
  batch integrity checks
\item
  fixed-zero parameter checks
\item
  output dimension checks
\item
  diagnostic finite/RMS positivity checks
\item
  moment-to-Twiss conversion
\item
  beam-state physical validity checks
\item
  validation report JSON/Markdown writing
\end{itemize}

\hypertarget{public-functions-and-classes-in-inspected-modules}{%
\subsubsection{Public Functions And Classes In Inspected
Modules}\label{public-functions-and-classes-in-inspected-modules}}

\hypertarget{no-offset-sobol-production}{%
\paragraph{No-offset Sobol
production}\label{no-offset-sobol-production}}

\begin{itemize}
\tightlist
\item
  \texttt{SobolConfig}
\item
  \texttt{NoOffsetSobolProductionConfig}
\item
  \texttt{replace\_config}
\item
  \texttt{load\_no\_offset\_sobol\_config}
\item
  \texttt{load\_no\_offset\_sobol\_production\_config}
\item
  \texttt{load\_no\_offset\_sobol\_preflight\_config}
\item
  \texttt{no\_offset\_sobol\_config\_from\_mapping}
\item
  \texttt{no\_offset\_sobol\_config\_to\_dict}
\item
  \texttt{sampled\_bounds}
\item
  \texttt{sobol\_unit\_samples\_for\_indices}
\item
  \texttt{physical\_samples\_for\_indices}
\item
  \texttt{expand\_19d\_to\_full\_24d}
\item
  \texttt{duplicate\_global\_indices}
\item
  \texttt{BatchPlanItem}
\item
  \texttt{plan\_batches}
\item
  \texttt{existing\_completed\_indices}
\item
  \texttt{assert\_sobol\_config\_compatible}
\item
  \texttt{make\_generation\_summary}
\item
  \texttt{print\_generation\_summary}
\item
  \texttt{write\_config\_files}
\item
  \texttt{run\_preflight}
\item
  \texttt{run\_generation}
\item
  \texttt{generate\_batch}
\item
  \texttt{write\_generation\_log}
\item
  \texttt{write\_sobol\_index}
\item
  \texttt{audit\_dataset\_batches}
\item
  \texttt{write\_validity\_summary}
\item
  \texttt{estimate\_runtime\_from\_preflight}
\end{itemize}

\hypertarget{general-lebt-dataset-generation}{%
\paragraph{General LEBT dataset
generation}\label{general-lebt-dataset-generation}}

\begin{itemize}
\tightlist
\item
  \texttt{ParameterRange}
\item
  \texttt{TwissRangeConfig}
\item
  \texttt{EmittanceRangeConfig}
\item
  \texttt{InitialOrbitRangeConfig}
\item
  \texttt{BeamCurrentRangeConfig}
\item
  \texttt{SteeringMagnetRangeConfig}
\item
  \texttt{ElectrostaticQuadrupoleRangeConfig}
\item
  \texttt{LEBTParameterRangeConfig}
\item
  \texttt{default\_lebt\_parameter\_ranges}
\item
  \texttt{resolve\_lebt\_range\_preset}
\item
  \texttt{DatasetGenerationConfig}
\item
  \texttt{LEBTDatasetConfig}
\item
  \texttt{LEBTDryRunValidationResult}
\item
  \texttt{validate\_lebt\_dry\_run\_config}
\item
  \texttt{format\_input\_summary\_table}
\item
  \texttt{LEBTDatasetGenerationResult}
\item
  \texttt{deterministic\_seed}
\item
  \texttt{generate\_control\_samples}
\item
  \texttt{generate\_input\_samples}
\item
  \texttt{generate\_sobol\_controls}
\item
  \texttt{generate\_lhs\_controls}
\item
  \texttt{generate\_lebt\_dataset}
\item
  \texttt{beam\_state\_output\_vector}
\item
  \texttt{derive\_twiss\_from\_state\_moments}
\item
  \texttt{validate\_state\_moment\_labels}
\end{itemize}

\hypertarget{input-and-target-views}{%
\paragraph{Input and target views}\label{input-and-target-views}}

\begin{itemize}
\tightlist
\item
  \texttt{InputViewSpec}
\item
  \texttt{get\_input\_columns}
\item
  \texttt{select\_input\_view}
\item
  \texttt{resolve\_target\_view\_names}
\item
  \texttt{load\_target\_view}
\end{itemize}

\hypertarget{no-offset-validation}{%
\paragraph{No-offset validation}\label{no-offset-validation}}

\begin{itemize}
\tightlist
\item
  \texttt{ValidationIssue}
\item
  \texttt{load\_metadata}
\item
  \texttt{load\_json\_file}
\item
  \texttt{load\_config\_yaml}
\item
  \texttt{discover\_batch\_files}
\item
  \texttt{validate\_required\_files}
\item
  \texttt{validate\_config\_values}
\item
  \texttt{validate\_sobol\_index}
\item
  \texttt{read\_batch\_global\_indices}
\item
  \texttt{check\_fixed\_zero\_parameters}
\item
  \texttt{check\_output\_dimensions}
\item
  \texttt{check\_diagnostic\_rms\_positive}
\item
  \texttt{moments\_to\_twiss}
\item
  \texttt{check\_state\_physics}
\item
  \texttt{validate\_batch\_integrity}
\item
  \texttt{validate\_no\_offset\_sobol\_dataset}
\item
  \texttt{write\_validation\_reports}
\item
  \texttt{format\_validation\_report\_markdown}
\end{itemize}

\hypertarget{notebook-usage}{%
\subsubsection{Notebook Usage}\label{notebook-usage}}

\texttt{04\_lebt\_no\_offset\_sobol\_dataset\_production.ipynb}
currently uses package code for:

\begin{itemize}
\tightlist
\item
  loading no-offset Sobol config
\item
  displaying resolved config
\item
  generation summary
\item
  preflight execution
\item
  optional preflight scaling runs
\item
  production generation
\item
  extension generation
\item
  post-generation audit
\end{itemize}

Important safety behavior:

\begin{itemize}
\tightlist
\item
  full production is gated by \texttt{RUN\_FULL\_PRODUCTION\ =\ False}
\item
  extension is gated by \texttt{RUN\_EXTENSION\ =\ False}
\item
  scale testing is gated by
  \texttt{RUN\_PREFLIGHT\_SCALE\_TEST\ =\ False}
\end{itemize}

\texttt{05\_lebt\_no\_offset\_sobol\_dataset\_validation.ipynb}
currently uses package code for:

\begin{itemize}
\tightlist
\item
  loading metadata/config/Sobol index
\item
  validating required files
\item
  validating batch integrity
\item
  validating fixed-zero parameters
\item
  validating diagnostics and beam-state physics
\item
  writing validation reports
\item
  optional plots
\end{itemize}

\hypertarget{preflight-path-inspection}{%
\subsubsection{Preflight Path
Inspection}\label{preflight-path-inspection}}

The no-offset Sobol preflight path is available as:

\begin{Shaded}
\begin{Highlighting}[]
\NormalTok{run\_preflight(config)}
\end{Highlighting}
\end{Shaded}

Current behavior:

\begin{itemize}
\tightlist
\item
  derives
  \texttt{preflight\_dir\ =\ config.dataset\_path\ /\ "preflight"}
\item
  replaces the requested config with:

  \begin{itemize}
  \tightlist
  \item
    \texttt{target\_total\_samples\ =\ 240}
  \item
    \texttt{mode\ =\ "preflight"}
  \item
    \texttt{batch\_size\ =\ 120}
  \item
    \texttt{overwrite\ =\ True}
  \item
    \texttt{resume\ =\ False}
  \end{itemize}
\item
  calls \texttt{run\_generation(preflight,\ backend=backend)}
\end{itemize}

For this R0 report, the preflight path was inspected and config-loaded
but not executed with TRACK. The loaded versioned configs reported:

\begin{Shaded}
\begin{Highlighting}[]
\NormalTok{production: dataset=lebt\_production\_no\_offset\_multitarget\_sobol\_v1, mode=production, samples=200000, batch\_size=2400, n\_workers=20}
\NormalTok{preflight:  dataset=lebt\_production\_no\_offset\_multitarget\_sobol\_v1, mode=preflight,  samples=240,    batch\_size=120,  n\_workers=20}
\NormalTok{dimensions: sampled=19, fixed=5, full=24, targets=\{diagnostics: 24, state\_last\_ws: 10, state\_lebt\_exit: 10, total: 44\}}
\end{Highlighting}
\end{Shaded}

\hypertarget{duplicated-logic}{%
\subsubsection{Duplicated Logic}\label{duplicated-logic}}

The main duplication is between \texttt{dataset\_generator.py} and
\texttt{no\_offset\_sobol\_production.py}:

\begin{itemize}
\tightlist
\item
  configuration validation
\item
  deterministic seed/sample policy
\item
  Sobol sampling
\item
  parallel worker loops
\item
  progress printing
\item
  work directory cleanup
\item
  HDF5 dataset writing
\item
  metadata/config writing
\item
  success/completed masks
\item
  failure reason handling
\item
  state output storage
\item
  state moment/Twiss validation
\item
  audit/report summaries
\end{itemize}

There is also some overlap between validation helpers and generation
audit code:

\begin{itemize}
\tightlist
\item
  diagnostic finite checks
\item
  RMS positivity checks
\item
  state moment validity checks
\item
  target dimension checks
\item
  duplicate index checks
\end{itemize}

\hypertarget{current-test-coverage}{%
\subsubsection{Current Test Coverage}\label{current-test-coverage}}

Relevant tests exercised for this report:

\begin{Shaded}
\begin{Highlighting}[]
\NormalTok{python3 {-}m pytest tests/test\_no\_offset\_sobol\_production.py tests/test\_no\_offset\_sobol\_validation.py tests/test\_lebt\_dataset\_generator.py tests/test\_input\_views.py tests/test\_dataset\_schema.py {-}q}
\end{Highlighting}
\end{Shaded}

Result:

\begin{Shaded}
\begin{Highlighting}[]
\NormalTok{76 passed}
\end{Highlighting}
\end{Shaded}

Full suite exercised:

\begin{Shaded}
\begin{Highlighting}[]
\NormalTok{python3 {-}m pytest {-}q}
\end{Highlighting}
\end{Shaded}

Result:

\begin{Shaded}
\begin{Highlighting}[]
\NormalTok{all tests passed, 3 skipped}
\end{Highlighting}
\end{Shaded}

Expected skips:

\begin{itemize}
\tightlist
\item
  CUDA test skipped because CUDA is not available.
\item
  JuTrack integration skipped unless
  \texttt{SYNAPTICTRACK\_RUN\_JUTRACK=1}.
\item
  Notebook execution skipped unless
  \texttt{SYNAPTICTRACK\_RUN\_NOTEBOOKS=1}.
\end{itemize}

Coverage strengths:

\begin{itemize}
\tightlist
\item
  Sobol determinism and continuation
\item
  non-power-of-two Sobol range warning prevention
\item
  duplicate global sample index detection
\item
  batch resume and extension
\item
  fixed-zero parameter placement
\item
  dry-run 19D/24D/44D shapes
\item
  target groups and failure handling
\item
  validation report writing
\item
  input view selection and ordering
\item
  target view loading and combined target views
\item
  state moment/Twiss physical checks on generated labels
\item
  config load and invalid sample-count policy
\end{itemize}

Coverage gaps:

\begin{itemize}
\tightlist
\item
  production notebook execution is not run by default
\item
  actual TRACK-backed preflight is not run in CI
\item
  large production resume/extension is covered by small mock tests, not
  full-scale files
\item
  performance/scaling behavior is notebook-driven and not unit-tested
  beyond config plumbing
\item
  local machine override behavior is tested through temporary files, not
  the default ignored local path
\end{itemize}

\hypertarget{known-risky-areas}{%
\subsubsection{Known Risky Areas}\label{known-risky-areas}}

\begin{enumerate}
\def\labelenumi{\arabic{enumi}.}
\item
  \texttt{no\_offset\_sobol\_production.py} remains a large specialized
  orchestration module. Its responsibilities include sampling, config,
  execution, I/O, audit, resume, and reports.
\item
  \texttt{dataset\_generator.py} and no-offset production generation
  still have overlapping execution machinery. This can cause future
  drift when adding MEBT or other accelerators.
\item
  HDF5 schema compatibility is critical. Existing notebooks/tests expect
  legacy and grouped names to keep working.
\item
  Global Sobol index semantics are high-risk. Production extension must
  never regenerate, duplicate, or reshuffle previous points.
\item
  Failed/lost-beam samples must remain explicit. Any refactor must
  preserve masks, \texttt{success}, \texttt{completed}, and
  \texttt{failure\_reason}.
\item
  Beam-state moment physics checks are sensitive to numerical
  tolerances. Refactors should preserve \texttt{PHYS\_TOL}-style
  near-zero handling.
\item
  Runtime paths and local config overrides are environment-sensitive.
  The production notebook must avoid hard-coded machine paths and must
  show the resolved config before generation.
\item
  The current no-offset Sobol workflow is specialized. Refactors should
  not make general accelerator abstractions depend on LEBT-specific
  names like \texttt{state\_last\_ws}.
\end{enumerate}

\hypertarget{recommended-refactor-order}{%
\subsubsection{Recommended Refactor
Order}\label{recommended-refactor-order}}

\begin{enumerate}
\def\labelenumi{\arabic{enumi}.}
\item
  Preserve and expand schema centralization. Keep names, dimensions,
  masks, aliases, and target groups behind schema objects.
\item
  Extract config loading and validation policy. Keep
  production/preflight/test modes explicit and reject invalid production
  sample counts loudly.
\item
  Split no-offset Sobol production into focused modules:

  \begin{itemize}
  \tightlist
  \item
    sampling and Sobol index handling
  \item
    batch planning and resume policy
  \item
    batch HDF5 I/O
  \item
    execution orchestration
  \item
    audit/report generation
  \end{itemize}
\item
  Extract shared generation runtime utilities from
  \texttt{dataset\_generator.py} and
  \texttt{no\_offset\_sobol\_production.py}:

  \begin{itemize}
  \tightlist
  \item
    worker loop
  \item
    work directory policy
  \item
    progress reporting
  \item
    failure result normalization
  \end{itemize}
\item
  Generalize accelerator/beamline schema support. Add future schemas
  without forcing MEBT or general accelerator workflows into
  LEBT-specific naming.
\item
  Move notebook logic steadily into package modules. Notebooks should
  load config, display config, call package functions, and visualize
  outputs.
\item
  Only after the above, consider broader API cleanup. Keep compatibility
  shims for existing notebooks/tests until replacement workflows are
  stable.
\end{enumerate}

\hypertarget{acceptance-status}{%
\subsubsection{Acceptance Status}\label{acceptance-status}}

\begin{itemize}
\tightlist
\item
  Behavior changes: none introduced by this R0 report.
\item
  Production notebook changes: none introduced by this R0 report.
\item
  Production generation: not run.
\item
  TRACK preflight generation: not run; path inspected and configs
  loaded.
\item
  Relevant tests: passed.
\item
  Full test suite: passed with expected skips.
\end{itemize}

\hypertarget{simulation-procedures-and-production-dataset-workflow}{%
\section{Simulation Procedures and Production Dataset
Workflow}\label{simulation-procedures-and-production-dataset-workflow}}

Source Markdown: \texttt{docs/lebt\_no\_offset\_sobol\_workflow.md}

\hypertarget{lebt-no-offset-sobol-dataset-workflow}{%
\subsection{LEBT No-Offset Sobol Dataset
Workflow}\label{lebt-no-offset-sobol-dataset-workflow}}

This is the maintained production workflow for the no-offset LEBT
surrogate dataset:

\begin{Shaded}
\begin{Highlighting}[]
\NormalTok{lebt\_production\_no\_offset\_multitarget\_sobol\_v1}
\NormalTok{19D sampled input {-}\textgreater{} 44D target output}
\end{Highlighting}
\end{Shaded}

The notebook-facing API is
\texttt{synapticTrack.simulation.no\_offset\_sobol\_production}. Its
public imports are maintained for production notebooks, while
implementation details live in focused modules such as schema, sampling,
batch planning, batch I/O, audit, and runtime helpers.

Beam current metadata uses canonical \texttt{uA} values and an explicit
\texttt{track\_current\_per\_synaptic\_current} conversion factor. The
default factor is \texttt{1.0} only for backward compatibility; TRACK
current units remain unverified by local project documentation.

\hypertarget{fixed-parameter-contract}{%
\subsubsection{Fixed-Parameter
Contract}\label{fixed-parameter-contract}}

The sampled input is 19D. A full 24D simulation input is accepted by the
LEBT surrogate only when its fixed parameters match the dataset schema:
\texttt{x0}, \texttt{xp0}, \texttt{y0}, and \texttt{yp0} are
\texttt{0.0}, while \texttt{beam\_current} is fixed at \texttt{0.05} in
the current TRACK input convention. Full-input validation checks these
expected values with a configurable tolerance; it does not treat every
fixed parameter as zero.

\hypertarget{configure}{%
\subsubsection{Configure}\label{configure}}

Use the versioned config:

\begin{Shaded}
\begin{Highlighting}[]
\NormalTok{configs/datasets/lebt\_no\_offset\_sobol\_production.yaml}
\end{Highlighting}
\end{Shaded}

Put machine-specific paths in the local override file, which is
git-ignored:

\begin{Shaded}
\begin{Highlighting}[]
\NormalTok{configs/local/lebt\_no\_offset\_sobol\_local.yaml}
\end{Highlighting}
\end{Shaded}

Typical local override fields:

\begin{Shaded}
\begin{Highlighting}[]
\FunctionTok{track\_exe\_path}\KeywordTok{:}\AttributeTok{ TRACK/TRACKv39C.exe}
\FunctionTok{source\_dir}\KeywordTok{:}\AttributeTok{ examples/lebt\_orbit\_corrections/orbit\_corrections}
\FunctionTok{calibration\_dir}\KeywordTok{:}\AttributeTok{ examples/lebt\_orbit\_corrections/orbit\_corrections/calibrations}
\FunctionTok{work\_root}\KeywordTok{:}\AttributeTok{ run.production/lebt\_production\_no\_offset\_multitarget\_sobol\_v1}
\end{Highlighting}
\end{Shaded}

\hypertarget{production-notebook}{%
\subsubsection{Production Notebook}\label{production-notebook}}

Run:

\begin{Shaded}
\begin{Highlighting}[]
\NormalTok{notebooks/04\_lebt\_no\_offset\_sobol\_dataset\_production.ipynb}
\end{Highlighting}
\end{Shaded}

Default behavior is safe:

\begin{Shaded}
\begin{Highlighting}[]
\NormalTok{RUN\_PREFLIGHT }\OperatorTok{=} \VariableTok{True}
\NormalTok{RUN\_PREFLIGHT\_SCALE\_TEST }\OperatorTok{=} \VariableTok{False}
\NormalTok{RUN\_PRODUCTION\_100K }\OperatorTok{=} \VariableTok{False}
\NormalTok{RUN\_EXTENSION\_200K }\OperatorTok{=} \VariableTok{False}
\NormalTok{RUN\_EXTENSION\_300K }\OperatorTok{=} \VariableTok{False}
\NormalTok{RUN\_EXTENSION\_400K }\OperatorTok{=} \VariableTok{False}
\NormalTok{RUN\_EXTENSION\_500K }\OperatorTok{=} \VariableTok{False}
\NormalTok{RUN\_EXTENSION\_600K }\OperatorTok{=} \VariableTok{False}
\NormalTok{RUN\_EXTENSION\_700K }\OperatorTok{=} \VariableTok{False}
\NormalTok{RUN\_EXTENSION\_800K }\OperatorTok{=} \VariableTok{False}
\NormalTok{RUN\_EXTENSION\_900K }\OperatorTok{=} \VariableTok{False}
\NormalTok{RUN\_EXTENSION\_1000K }\OperatorTok{=} \VariableTok{False}
\end{Highlighting}
\end{Shaded}

Only enable a production or extension stage after preflight passes. The
parallel scale-test cell is optional because it launches multiple TRACK
preflights and can select the fastest \texttt{n\_workers} for staged
production.

\hypertarget{sobol-extension}{%
\subsubsection{Sobol Extension}\label{sobol-extension}}

The supported production totals are:

\begin{Shaded}
\begin{Highlighting}[]
\NormalTok{100\_000 {-}\textgreater{} 200\_000 {-}\textgreater{} ... {-}\textgreater{} 900\_000 {-}\textgreater{} 1\_000\_000}
\end{Highlighting}
\end{Shaded}

To extend, keep the same dataset name, Sobol dimension, scramble flag,
and scramble seed. Enable one stage cell at a time:

\begin{Shaded}
\begin{Highlighting}[]
\NormalTok{RUN\_PRODUCTION\_100K }\OperatorTok{=} \VariableTok{True}
\NormalTok{RUN\_EXTENSION\_200K }\OperatorTok{=} \VariableTok{True}
\NormalTok{RUN\_EXTENSION\_300K }\OperatorTok{=} \VariableTok{True}
\NormalTok{RUN\_EXTENSION\_400K }\OperatorTok{=} \VariableTok{True}
\NormalTok{RUN\_EXTENSION\_500K }\OperatorTok{=} \VariableTok{True}
\NormalTok{RUN\_EXTENSION\_600K }\OperatorTok{=} \VariableTok{True}
\NormalTok{RUN\_EXTENSION\_700K }\OperatorTok{=} \VariableTok{True}
\NormalTok{RUN\_EXTENSION\_800K }\OperatorTok{=} \VariableTok{True}
\NormalTok{RUN\_EXTENSION\_900K }\OperatorTok{=} \VariableTok{True}
\NormalTok{RUN\_EXTENSION\_1000K }\OperatorTok{=} \VariableTok{True}
\NormalTok{OVERWRITE }\OperatorTok{=} \VariableTok{False}
\NormalTok{RESUME }\OperatorTok{=} \VariableTok{True}
\end{Highlighting}
\end{Shaded}

The generator plans missing global Sobol indices only. It must not
reshuffle or regenerate completed samples.

\hypertarget{consolidated-stage-files}{%
\subsubsection{Consolidated Stage
Files}\label{consolidated-stage-files}}

Generation writes resumable batch files under:

\begin{Shaded}
\begin{Highlighting}[]
\NormalTok{data/processed/lebt\_production\_no\_offset\_multitarget\_sobol\_v1/batches/}
\end{Highlighting}
\end{Shaded}

After a stage completes, you can optionally write one read-optimized
HDF5 file for that stage without regenerating samples:

\begin{Shaded}
\begin{Highlighting}[]
\ImportTok{from}\NormalTok{ synapticTrack.simulation.no\_offset\_sobol\_production }\ImportTok{import}\NormalTok{ consolidate\_dataset\_stage}

\ControlFlowTok{for}\NormalTok{ stage\_total }\KeywordTok{in} \BuiltInTok{range}\NormalTok{(}\DecValTok{100\_000}\NormalTok{, }\DecValTok{1\_000\_001}\NormalTok{, }\DecValTok{100\_000}\NormalTok{):}
\NormalTok{    consolidate\_dataset\_stage(config, target\_total\_samples}\OperatorTok{=}\NormalTok{stage\_total)}
\end{Highlighting}
\end{Shaded}

Default file names are:

\begin{Shaded}
\begin{Highlighting}[]
\NormalTok{data/processed/lebt\_production\_no\_offset\_multitarget\_sobol\_v1/lebt\_production\_no\_offset\_multitarget\_sobol\_v1\_100000.h5}
\NormalTok{...}
\NormalTok{data/processed/lebt\_production\_no\_offset\_multitarget\_sobol\_v1/lebt\_production\_no\_offset\_multitarget\_sobol\_v1\_1000000.h5}
\end{Highlighting}
\end{Shaded}

The production notebook exposes \texttt{RUN\_CONSOLIDATION} and
\texttt{CONSOLIDATION\_STAGE\_TOTALS} for this step. Consolidation
checks duplicate indices and requires a contiguous completed range from
global Sobol index 0 to the requested stage total.

\hypertarget{validation-notebook}{%
\subsubsection{Validation Notebook}\label{validation-notebook}}

Run:

\begin{Shaded}
\begin{Highlighting}[]
\NormalTok{notebooks/05\_lebt\_no\_offset\_sobol\_dataset\_validation.ipynb}
\end{Highlighting}
\end{Shaded}

The notebook calls the pure audit API first:

\begin{Shaded}
\begin{Highlighting}[]
\ImportTok{from}\NormalTok{ synapticTrack.data.production\_audit }\ImportTok{import}\NormalTok{ compute\_validation\_report}
\NormalTok{report }\OperatorTok{=}\NormalTok{ compute\_validation\_report(DATASET\_DIR)}
\end{Highlighting}
\end{Shaded}

Reports are written only when requested:

\begin{Shaded}
\begin{Highlighting}[]
\NormalTok{WRITE\_REPORTS }\OperatorTok{=} \VariableTok{True}
\end{Highlighting}
\end{Shaded}

Outputs:

\begin{Shaded}
\begin{Highlighting}[]
\NormalTok{validation\_report.json}
\NormalTok{validation\_report.md}
\end{Highlighting}
\end{Shaded}

\hypertarget{training-and-publication-metrics}{%
\subsubsection{Training and Publication
Metrics}\label{training-and-publication-metrics}}

The staged no-offset 44D surrogate workflow includes plain MLP and
physics-regularized MLP training notebooks, model-comparison notebooks,
and publication-metric table generation. Paper-oriented metrics should
not use a single mixed-unit 10D state RMSE as the main physical metric.
\texttt{compute\_state\_physics\_metrics()} is the canonical report
shared by plain and physics-regularized MLP evaluation. Use diagnostic
RMSE in millimetres, component-wise state errors, normalized state
errors, derived Courant-Snyder/emittance errors, and covariance-validity
fractions. The mixed-unit value remains a marked legacy comparison
metric only.

Current publication tables are generated under:

\begin{Shaded}
\begin{Highlighting}[]
\NormalTok{models/lebt\_no\_offset\_model\_comparison/publication\_metrics/}
\end{Highlighting}
\end{Shaded}

The best diagnostics-only configuration from the available staged runs
is the plain MLP at 1M samples. The physics-regularized 900k model is
preferred when covariance-valid beam-state predictions are required.

\hypertarget{api-status}{%
\subsubsection{API Status}\label{api-status}}

Maintained production APIs:

\begin{Shaded}
\begin{Highlighting}[]
\ImportTok{from}\NormalTok{ synapticTrack.simulation.no\_offset\_sobol\_production }\ImportTok{import}\NormalTok{ (}
\NormalTok{    load\_no\_offset\_sobol\_config,}
\NormalTok{    run\_preflight,}
\NormalTok{    run\_generation,}
\NormalTok{    make\_generation\_summary,}
\NormalTok{    plan\_batches,}
\NormalTok{    audit\_dataset\_batches,}
\NormalTok{)}
\end{Highlighting}
\end{Shaded}

Maintained shared modules:

\begin{Shaded}
\begin{Highlighting}[]
\NormalTok{synapticTrack.data.schema}
\NormalTok{synapticTrack.data.batch\_io}
\NormalTok{synapticTrack.data.production\_audit}
\NormalTok{synapticTrack.sampling.sobol}
\NormalTok{synapticTrack.sampling.latin\_hypercube}
\NormalTok{synapticTrack.simulation.batch\_planning}
\NormalTok{synapticTrack.simulation.generation\_runtime}
\NormalTok{synapticTrack.utils.parallel}
\end{Highlighting}
\end{Shaded}

Legacy/generic LEBT APIs in
\texttt{synapticTrack.simulation.dataset\_generator} are still supported
for smoke datasets, dry runs, and development workflows. They should not
be treated as the production no-offset Sobol pipeline.

\hypertarget{cli-note}{%
\subsubsection{CLI note}\label{cli-note}}

The current \texttt{synapticTrack} CLI \texttt{lebt-production-*}
commands target the legacy/generic LEBT production-config path. They are
not the maintained no-offset Sobol production workflow. Use the
production notebook or package APIs above for
\texttt{lebt\_production\_no\_offset\_multitarget\_sobol\_v1}.

\hypertarget{artifact-and-dataset-compatibility}{%
\subsubsection{Artifact And Dataset
Compatibility}\label{artifact-and-dataset-compatibility}}

Each new training run writes \texttt{artifact\_manifest.json} with
format \texttt{training\_run\_artifacts\_v2}. It records the stable
paths for checkpoints, configuration, timing, plots, summaries, standard
metrics,
\texttt{state\_physics\_metrics\_\textless{}split\textgreater{}.json},
compatibility
\texttt{physics\_metrics\_\textless{}split\textgreater{}.json}, derived
metrics, and covariance-validity metrics. Existing run directories
without that file remain readable as an explicit legacy layout through
\texttt{discover\_run\_artifacts()}.

No-offset loaders detect the HDF5 format before reading layout-specific
paths. Current no-offset batch and consolidated formats are normalized
through \texttt{read\_hdf5\_dataset\_compat()}. See
\texttt{docs/hdf5\_dataset\_compatibility.md} for supported legacy LEBT
and MEBT formats and the future opt-in migration path.

Source Markdown:
\texttt{docs/lebt\_full\_production\_simulation\_and\_analysis.md}

\hypertarget{lebt-full-production-simulation-and-analysis-runner}{%
\subsection{LEBT Full Production Simulation and Analysis
Runner}\label{lebt-full-production-simulation-and-analysis-runner}}

This document describes the one-file shell runner:

\begin{Shaded}
\begin{Highlighting}[]
\ExtensionTok{scripts/run\_lebt\_full\_production\_and\_analysis.sh}
\end{Highlighting}
\end{Shaded}

The script runs the maintained no-offset Sobol LEBT production workflow
and then runs dataset analysis/validation. It is intended for long
full-production runs on a new machine, local workstation, Codex
environment, or Antigravity prompt where screen output should be
controlled.

\hypertarget{what-the-script-runs}{%
\subsubsection{What the Script Runs}\label{what-the-script-runs}}

The script is organized as a sequence of Python-backed execution stages.
Each stage uses package APIs from \texttt{synapticTrack}; production
logic is not duplicated inside notebooks or ad hoc scripts.

\begin{longtable}[]{@{}
  >{\raggedright\arraybackslash}p{(\columnwidth - 6\tabcolsep) * \real{0.2500}}
  >{\raggedright\arraybackslash}p{(\columnwidth - 6\tabcolsep) * \real{0.2500}}
  >{\raggedright\arraybackslash}p{(\columnwidth - 6\tabcolsep) * \real{0.2500}}
  >{\raggedright\arraybackslash}p{(\columnwidth - 6\tabcolsep) * \real{0.2500}}@{}}
\toprule\noalign{}
\begin{minipage}[b]{\linewidth}\raggedright
Stage
\end{minipage} & \begin{minipage}[b]{\linewidth}\raggedright
Shell label
\end{minipage} & \begin{minipage}[b]{\linewidth}\raggedright
Notebook section
\end{minipage} & \begin{minipage}[b]{\linewidth}\raggedright
Purpose
\end{minipage} \\
\midrule\noalign{}
\endhead
\bottomrule\noalign{}
\endlastfoot
0 & \texttt{00\_manifest} & \texttt{00\ Manifest} & Record run
parameters, folders, Python/platform metadata, and local Git
metadata. \\
1 & \texttt{01\_preflight} & \texttt{01\ Preflight} & Run the 240-sample
preflight using the selected scan and runtime options. \\
2 & \texttt{02\_production\_stages} & \texttt{02\ Production\ Stages} &
Run cumulative Sobol production stages, normally
\texttt{100000,200000,300000,400000,500000}. \\
3 & \texttt{03\_consolidation} & \texttt{03\ Consolidation} & Write
read-optimized consolidated HDF5 files for requested stages. \\
4 & \texttt{04\_analysis\_validation} &
\texttt{04\ Analysis\ and\ Validation} & Compute final validation
reports and write JSON/Markdown summaries. \\
\end{longtable}

\hypertarget{default-output-layout}{%
\subsubsection{Default Output Layout}\label{default-output-layout}}

Every execution uses a run-specific ID. If \texttt{-\/-run-id} is not
provided, the script uses a timestamp.

\begin{Shaded}
\begin{Highlighting}[]
\NormalTok{data/processed/full\_production\_runs/\textless{}run{-}id\textgreater{}/dataset/}
\NormalTok{run.production/full\_production\_runs/\textless{}run{-}id\textgreater{}/work/}
\NormalTok{run.production/full\_production\_runs/\textless{}run{-}id\textgreater{}/logs/}
\end{Highlighting}
\end{Shaded}

This keeps simulation outputs, temporary TRACK work files, and logs in
separate folders. Override these with \texttt{-\/-dataset-dir},
\texttt{-\/-work-root}, and \texttt{-\/-log-dir}.

\hypertarget{parallel-execution}{%
\subsubsection{Parallel Execution}\label{parallel-execution}}

By default, the script computes the number of online CPU cores with:

\begin{Shaded}
\begin{Highlighting}[]
\ExtensionTok{getconf}\NormalTok{ \_NPROCESSORS\_ONLN}
\end{Highlighting}
\end{Shaded}

It then uses 90 percent of those cores for production workers. Override
this with either:

\begin{Shaded}
\begin{Highlighting}[]
\ExtensionTok{{-}{-}n{-}workers}\NormalTok{ 32}
\ExtensionTok{{-}{-}cpu{-}percent}\NormalTok{ 75}
\end{Highlighting}
\end{Shaded}

\texttt{-\/-n-workers} takes precedence over \texttt{-\/-cpu-percent}.

\hypertarget{verbose-and-quiet-modes}{%
\subsubsection{Verbose and Quiet Modes}\label{verbose-and-quiet-modes}}

Quiet mode is the default. It disables sample progress and TRACK
verbosity on screen and redirects each Python stage into a run-specific
log file. This is the recommended mode for Codex or Antigravity prompt
execution because it avoids large token consumption during full
simulation.

\begin{Shaded}
\begin{Highlighting}[]
\ExtensionTok{scripts/run\_lebt\_full\_production\_and\_analysis.sh} \AttributeTok{{-}{-}quiet}
\end{Highlighting}
\end{Shaded}

Verbose mode prints progress on screen and uses the requested TRACK
verbosity:

\begin{Shaded}
\begin{Highlighting}[]
\ExtensionTok{scripts/run\_lebt\_full\_production\_and\_analysis.sh} \AttributeTok{{-}{-}verbose} \AttributeTok{{-}{-}track{-}verbose}\NormalTok{ 1}
\end{Highlighting}
\end{Shaded}

\hypertarget{full-production-example}{%
\subsubsection{Full Production Example}\label{full-production-example}}

\begin{Shaded}
\begin{Highlighting}[]
\ExtensionTok{scripts/run\_lebt\_full\_production\_and\_analysis.sh} \DataTypeTok{\textbackslash{}}
  \AttributeTok{{-}{-}run{-}id}\NormalTok{ lebt\_no\_offset\_500k\_v1 }\DataTypeTok{\textbackslash{}}
  \AttributeTok{{-}{-}stages}\NormalTok{ 100000,200000,300000,400000,500000 }\DataTypeTok{\textbackslash{}}
  \AttributeTok{{-}{-}quiet}
\end{Highlighting}
\end{Shaded}

The script applies
\texttt{configs/local/lebt\_no\_offset\_sobol\_local.yaml} by default if
that file exists. Use it for machine-specific TRACK paths and source
directories. To ignore local overrides:

\begin{Shaded}
\begin{Highlighting}[]
\ExtensionTok{scripts/run\_lebt\_full\_production\_and\_analysis.sh} \AttributeTok{{-}{-}no{-}local{-}override}
\end{Highlighting}
\end{Shaded}

\hypertarget{test-sized-pipeline-check}{%
\subsubsection{Test-Sized Pipeline
Check}\label{test-sized-pipeline-check}}

Use \texttt{-\/-dry-run} and \texttt{-\/-allow-nonstandard} for a small
local check that does not call TRACK:

\begin{Shaded}
\begin{Highlighting}[]
\ExtensionTok{scripts/run\_lebt\_full\_production\_and\_analysis.sh} \DataTypeTok{\textbackslash{}}
  \AttributeTok{{-}{-}run{-}id}\NormalTok{ dry\_run\_240 }\DataTypeTok{\textbackslash{}}
  \AttributeTok{{-}{-}stages}\NormalTok{ 240 }\DataTypeTok{\textbackslash{}}
  \AttributeTok{{-}{-}batch{-}size}\NormalTok{ 120 }\DataTypeTok{\textbackslash{}}
  \AttributeTok{{-}{-}n{-}workers}\NormalTok{ 2 }\DataTypeTok{\textbackslash{}}
  \AttributeTok{{-}{-}dry{-}run} \DataTypeTok{\textbackslash{}}
  \AttributeTok{{-}{-}allow{-}nonstandard} \DataTypeTok{\textbackslash{}}
  \AttributeTok{{-}{-}skip{-}preflight} \DataTypeTok{\textbackslash{}}
  \AttributeTok{{-}{-}quiet}
\end{Highlighting}
\end{Shaded}

\hypertarget{analysis-outputs}{%
\subsubsection{Analysis Outputs}\label{analysis-outputs}}

The analysis stage writes the maintained validation reports into the
dataset directory:

\begin{Shaded}
\begin{Highlighting}[]
\NormalTok{validation\_report.json}
\NormalTok{validation\_report.md}
\NormalTok{final\_analysis\_summary.json}
\end{Highlighting}
\end{Shaded}

Generation also writes production metadata and audit files, including:

\begin{Shaded}
\begin{Highlighting}[]
\NormalTok{run\_manifest.json}
\NormalTok{generation\_log.csv}
\NormalTok{sobol\_index.json}
\NormalTok{audit\_report.json}
\NormalTok{validity\_summary.json}
\end{Highlighting}
\end{Shaded}

\hypertarget{matching-notebook}{%
\subsubsection{Matching Notebook}\label{matching-notebook}}

The matching notebook is:

\begin{Shaded}
\begin{Highlighting}[]
\NormalTok{notebooks/20\_lebt\_full\_production\_simulation\_and\_analysis.ipynb}
\end{Highlighting}
\end{Shaded}

It uses the same parameter names, output folder layout, stage order, and
package APIs as the shell runner. Use the shell script for unattended
full production and the notebook for step-by-step inspection, review,
and manuscript-supporting analysis notes.

\hypertarget{repository-operation-reminder}{%
\subsubsection{Repository Operation
Reminder}\label{repository-operation-reminder}}

Do not run \texttt{git\ push} unless the repository owner explicitly
requests it. Remote GitHub Actions should be checked only when
explicitly requested for a CI/debugging task. Local GitHub Actions
workflow files may be edited in the local repository when needed.

\hypertarget{analysis-methods-and-analysis-procedure}{%
\section{Analysis Methods and Analysis
Procedure}\label{analysis-methods-and-analysis-procedure}}

Source Markdown: \texttt{docs/scanner\_measurement\_analysis.md}

\hypertarget{scanner-measurement-analysis}{%
\subsection{Scanner Measurement
Analysis}\label{scanner-measurement-analysis}}

This note documents how scanner measurement data is manipulated by the
current analysis path in \texttt{synapticTrack}.

Relevant files:

\begin{itemize}
\tightlist
\item
  \texttt{src/synapticTrack/io/scanner\_io.py}
\item
  \texttt{src/synapticTrack/beam/beam\_scanner.py}
\item
  \texttt{src/synapticTrack/analysis/scanner\_analysis.py}
\item
  \texttt{src/synapticTrack/utils/stats.py}
\item
  \texttt{src/synapticTrack/utils/math\_functions.py}
\end{itemize}

\hypertarget{file-loading}{%
\subsubsection{File Loading}\label{file-loading}}

Scanner files are read with whitespace-separated columns and one skipped
header row.

Wire scanner files are read as:

\begin{Shaded}
\begin{Highlighting}[]
\NormalTok{x\_pos x\_current y\_pos y\_current d\_pos d\_current}
\end{Highlighting}
\end{Shaded}

The required analysis columns are:

\begin{Shaded}
\begin{Highlighting}[]
\NormalTok{x\_pos, x\_current, y\_pos, y\_current}
\end{Highlighting}
\end{Shaded}

Allison scanner files are read as:

\begin{Shaded}
\begin{Highlighting}[]
\NormalTok{x xp x\_current hv y\_current}
\end{Highlighting}
\end{Shaded}

The required analysis columns are:

\begin{Shaded}
\begin{Highlighting}[]
\NormalTok{x, xp, x\_current}
\end{Highlighting}
\end{Shaded}

The file stem is stored as \texttt{scan\_id}.

\hypertarget{wire-scanner-manipulation}{%
\subsubsection{Wire Scanner
Manipulation}\label{wire-scanner-manipulation}}

Wire scanner analysis is implemented by \texttt{analyze\_wire\_scanner}.

Each plane is processed independently:

\begin{Shaded}
\begin{Highlighting}[]
\NormalTok{x\_pos, x\_current}
\NormalTok{y\_pos, y\_current}
\end{Highlighting}
\end{Shaded}

The preprocessing function is
\texttt{\_preprocess\_profile(position,\ current,\ shift,\ cutoff,\ cutoff\_frac)}.

Processing steps:

\begin{enumerate}
\def\labelenumi{\arabic{enumi}.}
\item
  Convert \texttt{position} and \texttt{current} to floating-point NumPy
  arrays.
\item
  Require matching shapes.
\item
  Remove all samples where either position or current is NaN or Inf.
\item
  If \texttt{shift=True}, subtract the minimum current from the current
  profile:

\begin{Shaded}
\begin{Highlighting}[]
\NormalTok{current \textless{}{-} current {-} min(current)}
\end{Highlighting}
\end{Shaded}

  This is a baseline shift. It makes the smallest remaining current
  value zero.
\item
  If \texttt{cutoff=True}, remove low-current samples after the optional
  shift:

\begin{Shaded}
\begin{Highlighting}[]
\NormalTok{threshold = cutoff\_frac * max(current)}
\NormalTok{keep current \textgreater{}= threshold}
\end{Highlighting}
\end{Shaded}
\item
  Require at least 4 points after filtering.
\item
  Require positive total current weight:

\begin{Shaded}
\begin{Highlighting}[]
\NormalTok{sum(current) \textgreater{} 0}
\end{Highlighting}
\end{Shaded}
\end{enumerate}

Default behavior:

\begin{Shaded}
\begin{Highlighting}[]
\NormalTok{shift=False}
\NormalTok{cutoff=False}
\NormalTok{cutoff\_frac=0.05}
\end{Highlighting}
\end{Shaded}

So by default the analysis removes non-finite rows only. Baseline
subtraction and cutoff filtering are opt-in.

\hypertarget{wire-scanner-derived-quantities}{%
\subsubsection{Wire Scanner Derived
Quantities}\label{wire-scanner-derived-quantities}}

After preprocessing, weighted centroid and RMS size are computed with
current as the weight.

For one plane:

\begin{Shaded}
\begin{Highlighting}[]
\NormalTok{center = sum(current\_i * position\_i) / sum(current\_i)}
\NormalTok{rms    = sqrt(sum(current\_i * (position\_i {-} center)\^{}2) / sum(current\_i))}
\end{Highlighting}
\end{Shaded}

The returned weighted quantities are:

\begin{Shaded}
\begin{Highlighting}[]
\NormalTok{x\_center}
\NormalTok{y\_center}
\NormalTok{sigma\_x}
\NormalTok{sigma\_y}
\end{Highlighting}
\end{Shaded}

A separate Gaussian fit is also performed for each plane:

\begin{Shaded}
\begin{Highlighting}[]
\NormalTok{current(position) = a * exp({-}(position {-} x0)\^{}2 / (2 * sigma\^{}2)) + offset}
\end{Highlighting}
\end{Shaded}

Initial fit guesses are:

\begin{Shaded}
\begin{Highlighting}[]
\NormalTok{a      = max(current)}
\NormalTok{x0     = weighted center}
\NormalTok{sigma  = weighted rms}
\NormalTok{offset = 0}
\end{Highlighting}
\end{Shaded}

The returned fit quantities are:

\begin{Shaded}
\begin{Highlighting}[]
\NormalTok{gaussian\_fit\_x\_center}
\NormalTok{gaussian\_fit\_y\_center}
\NormalTok{gaussian\_fit\_sigma\_x}
\NormalTok{gaussian\_fit\_sigma\_y}
\end{Highlighting}
\end{Shaded}

The function also returns metadata about preprocessing:

\begin{Shaded}
\begin{Highlighting}[]
\NormalTok{shift\_applied}
\NormalTok{cutoff\_applied}
\NormalTok{cutoff\_fraction}
\NormalTok{x\_points\_used}
\NormalTok{y\_points\_used}
\end{Highlighting}
\end{Shaded}

If plotting is enabled, the plot uses the processed profiles, not the
original raw profiles.

\hypertarget{allison-scanner-manipulation}{%
\subsubsection{Allison Scanner
Manipulation}\label{allison-scanner-manipulation}}

Allison scanner analysis is implemented by
\texttt{analyze\_allison\_scanner\_2d}.

Current behavior uses the loaded arrays directly:

\begin{Shaded}
\begin{Highlighting}[]
\NormalTok{x}
\NormalTok{xp}
\NormalTok{x\_current}
\end{Highlighting}
\end{Shaded}

Unlike wire scanner analysis, there is currently no explicit finite
filtering, baseline shifting, or cutoff filtering in
\texttt{analyze\_allison\_scanner\_2d}.

Weighted centroids and RMS values are computed using \texttt{x\_current}
as the weight:

\begin{Shaded}
\begin{Highlighting}[]
\NormalTok{x\_center  = weighted mean of x}
\NormalTok{xp\_center = weighted mean of xp}
\NormalTok{sigma\_x   = weighted RMS of x}
\NormalTok{sigma\_xp  = weighted RMS of xp}
\end{Highlighting}
\end{Shaded}

The weighted covariance is:

\begin{Shaded}
\begin{Highlighting}[]
\NormalTok{covariance\_x\_xp =}
\NormalTok{    sum(x\_current\_i * (x\_i {-} x\_center) * (xp\_i {-} xp\_center))}
\NormalTok{    / sum(x\_current\_i)}
\end{Highlighting}
\end{Shaded}

The RMS emittance is:

\begin{Shaded}
\begin{Highlighting}[]
\NormalTok{emittance\_rms = sqrt(sigma\_x\^{}2 * sigma\_xp\^{}2 {-} covariance\_x\_xp\^{}2)}
\end{Highlighting}
\end{Shaded}

The current geometric-emittance estimate is:

\begin{Shaded}
\begin{Highlighting}[]
\NormalTok{emittance\_geometric = sigma\_x * sigma\_xp}
\end{Highlighting}
\end{Shaded}

Separate 1D Gaussian fits are performed for \texttt{x} and \texttt{xp}
projections using the same current weights as measured values:

\begin{Shaded}
\begin{Highlighting}[]
\NormalTok{x\_current(x)  = a * exp({-}(x  {-} x0)\^{}2  / (2 * sigma\_x\^{}2))  + offset}
\NormalTok{x\_current(xp) = a * exp({-}(xp {-} xp0)\^{}2 / (2 * sigma\_xp\^{}2)) + offset}
\end{Highlighting}
\end{Shaded}

Initial fit guesses are:

\begin{Shaded}
\begin{Highlighting}[]
\NormalTok{a      = max(x\_current)}
\NormalTok{center = weighted center}
\NormalTok{sigma  = weighted RMS}
\NormalTok{offset = 0}
\end{Highlighting}
\end{Shaded}

Returned Allison quantities:

\begin{Shaded}
\begin{Highlighting}[]
\NormalTok{x\_center}
\NormalTok{xp\_center}
\NormalTok{sigma\_x}
\NormalTok{sigma\_xp}
\NormalTok{gaussian\_fit\_x\_center}
\NormalTok{gaussian\_fit\_xp\_center}
\NormalTok{gaussian\_fit\_sigma\_x}
\NormalTok{gaussian\_fit\_sigma\_xp}
\NormalTok{covariance\_x\_xp}
\NormalTok{emittance\_rms}
\NormalTok{emittance\_geometric}
\end{Highlighting}
\end{Shaded}

If plotting is enabled, the plot uses the original loaded Allison
arrays.

\hypertarget{important-assumptions-and-caveats}{%
\subsubsection{Important Assumptions And
Caveats}\label{important-assumptions-and-caveats}}

\begin{itemize}
\tightlist
\item
  Wire scanner preprocessing is optional except for NaN/Inf removal.
\item
  Allison scanner analysis currently does not remove NaN/Inf values
  before computing moments or fitting.
\item
  Current values are used directly as weights. Negative current values
  are not clipped unless \texttt{shift=True} is used for wire scanner
  profiles.
\item
  Wire scanner Gaussian-fit sigmas are returned directly from
  \texttt{curve\_fit}; the code does not force them positive with
  \texttt{abs}.
\item
  The RMS moment calculation and Gaussian fit are separate estimates.
  Downstream code should be explicit about which one it uses.
\item
  Plot files reflect processed wire-scanner data when preprocessing is
  enabled.
\end{itemize}

Source Markdown: \texttt{docs/notebook\_hygiene.md}

\hypertarget{notebook-hygiene}{%
\subsection{Notebook Hygiene}\label{notebook-hygiene}}

Committed notebooks should contain source cells only:

\begin{itemize}
\tightlist
\item
  no execution counts;
\item
  no embedded outputs;
\item
  no absolute machine-specific paths;
\item
  no local TRACK paths in source cells.
\end{itemize}

Use package code from \texttt{src/synapticTrack}; do not redefine
production functions inside notebooks.

\hypertarget{covered-notebook-roots}{%
\subsubsection{Covered Notebook Roots}\label{covered-notebook-roots}}

Notebook hygiene tests cover notebooks in \texttt{examples/},
\texttt{notebooks/}, and \texttt{RAON/}. This includes mirrored RAON
LEBT notebooks, extended LEBT training notebooks, and MEBT TRACK
notebooks. Generated run folders and scratch/output files should be
ignored locally rather than committed.

\hypertarget{clear-outputs}{%
\subsubsection{Clear Outputs}\label{clear-outputs}}

Recommended workflow before committing notebooks:

\begin{Shaded}
\begin{Highlighting}[]
\ExtensionTok{python3}\NormalTok{ scripts/strip\_notebook\_outputs.py}
\ExtensionTok{python3}\NormalTok{ scripts/strip\_notebook\_outputs.py }\AttributeTok{{-}{-}check}
\end{Highlighting}
\end{Shaded}

By default, the script processes \texttt{examples/},
\texttt{notebooks/}, and \texttt{RAON/}, matching the notebook roots
covered by the tests, and skips \texttt{.ipynb\_checkpoints}. Pass
explicit files or directories to limit the scope:

\begin{Shaded}
\begin{Highlighting}[]
\ExtensionTok{python3}\NormalTok{ scripts/strip\_notebook\_outputs.py RAON/MEBT\_notebooks}
\end{Highlighting}
\end{Shaded}

\texttt{nbstripout} is also listed in \texttt{pyproject.toml} and can be
installed as a git filter, but the project script is dependency-free and
covers all notebook roots checked by the tests.

\hypertarget{local-paths}{%
\subsubsection{Local Paths}\label{local-paths}}

Use config files and environment variables for machine-specific paths.
For example:

\begin{Shaded}
\begin{Highlighting}[]
\NormalTok{TRACK\_EXE\_PATH }\OperatorTok{=}\NormalTok{ Path(os.environ.get(}\StringTok{\textquotesingle{}TRACK\_EXE\_PATH\textquotesingle{}}\NormalTok{, }\StringTok{\textquotesingle{}../TRACK/TRACKv39C.exe\textquotesingle{}}\NormalTok{)).expanduser().resolve()}
\end{Highlighting}
\end{Shaded}

For production Sobol generation, prefer
\texttt{configs/local/lebt\_no\_offset\_sobol\_local.yaml}; files under
\texttt{configs/local/} are git-ignored except examples.

\hypertarget{checks}{%
\subsubsection{Checks}\label{checks}}

The test suite includes notebook hygiene checks for project notebooks:

\begin{Shaded}
\begin{Highlighting}[]
\ExtensionTok{python3} \AttributeTok{{-}m}\NormalTok{ pytest tests/test\_notebooks.py }\AttributeTok{{-}q}
\end{Highlighting}
\end{Shaded}

\hypertarget{new-user-installation-and-local-operation}{%
\section{New User Installation and Local
Operation}\label{new-user-installation-and-local-operation}}

Source Markdown: \texttt{docs/INSTALL.md}

\hypertarget{installation-guide}{%
\subsection{Installation Guide}\label{installation-guide}}

This guide is for a new user setting up \texttt{synapticTrack} on a new
machine.

\hypertarget{prerequisites}{%
\subsubsection{Prerequisites}\label{prerequisites}}

\begin{itemize}
\tightlist
\item
  Git
\item
  Python 3.8 or newer
\item
  A working C/C++ build toolchain if your platform needs to compile
  Python wheels
\item
  Network access to the Python package index used by \texttt{pip}
\end{itemize}

\texttt{synapticTrack} is a Python package. Install it from the
repository root with \texttt{pip\ install\ -e\ .}; the project
dependencies are declared in \texttt{pyproject.toml}.

\hypertarget{quick-install}{%
\subsubsection{Quick Install}\label{quick-install}}

Replace \texttt{\textless{}repo-url\textgreater{}} with the actual Git
URL for this repository:

\begin{Shaded}
\begin{Highlighting}[]
\FunctionTok{git}\NormalTok{ clone }\OperatorTok{\textless{}}\NormalTok{repo{-}url}\OperatorTok{\textgreater{}}\NormalTok{ synapticTrack}
\BuiltInTok{cd}\NormalTok{ synapticTrack}

\ExtensionTok{python3} \AttributeTok{{-}m}\NormalTok{ venv .venv}
\BuiltInTok{source}\NormalTok{ .venv/bin/activate}

\ExtensionTok{python} \AttributeTok{{-}m}\NormalTok{ pip install }\AttributeTok{{-}{-}upgrade}\NormalTok{ pip setuptools wheel}
\ExtensionTok{python} \AttributeTok{{-}m}\NormalTok{ pip install }\AttributeTok{{-}e}\NormalTok{ .}
\ExtensionTok{python} \AttributeTok{{-}m}\NormalTok{ pytest}
\end{Highlighting}
\end{Shaded}

If your system uses \texttt{python} instead of \texttt{python3},
substitute that command in the \texttt{venv} creation step.

\hypertarget{scripted-install}{%
\subsubsection{Scripted Install}\label{scripted-install}}

The repository includes a bootstrap script that performs the same steps:

\begin{Shaded}
\begin{Highlighting}[]
\FunctionTok{bash}\NormalTok{ scripts/bootstrap\_install.sh }\OperatorTok{\textless{}}\NormalTok{repo{-}url}\OperatorTok{\textgreater{}}
\end{Highlighting}
\end{Shaded}

By default, the script clones into \texttt{synapticTrack}. To choose a
different target directory:

\begin{Shaded}
\begin{Highlighting}[]
\FunctionTok{bash}\NormalTok{ scripts/bootstrap\_install.sh }\OperatorTok{\textless{}}\NormalTok{repo{-}url}\OperatorTok{\textgreater{}}\NormalTok{ /path/to/synapticTrack}
\end{Highlighting}
\end{Shaded}

The script:

\begin{enumerate}
\def\labelenumi{\arabic{enumi}.}
\tightlist
\item
  checks that Git and Python are available;
\item
  clones the repository if the target directory does not already exist;
\item
  creates a virtual environment at \texttt{.venv};
\item
  installs the package in editable mode;
\item
  runs the full pytest suite.
\end{enumerate}

\hypertarget{reusing-an-existing-clone}{%
\subsubsection{Reusing an Existing
Clone}\label{reusing-an-existing-clone}}

If the target directory already exists, the script does not reclone it.
It reuses the existing checkout and continues with virtual environment
creation, package installation, and tests:

\begin{Shaded}
\begin{Highlighting}[]
\FunctionTok{bash}\NormalTok{ scripts/bootstrap\_install.sh }\OperatorTok{\textless{}}\NormalTok{repo{-}url}\OperatorTok{\textgreater{}}\NormalTok{ /existing/path/to/synapticTrack}
\end{Highlighting}
\end{Shaded}

\hypertarget{dependency-metadata}{%
\subsubsection{Dependency Metadata}\label{dependency-metadata}}

Continuous integration installs the package with:

\begin{Shaded}
\begin{Highlighting}[]
\ExtensionTok{python} \AttributeTok{{-}m}\NormalTok{ pip install }\AttributeTok{{-}e}\NormalTok{ .}\PreprocessorTok{[}\SpecialStringTok{dev}\PreprocessorTok{]}
\end{Highlighting}
\end{Shaded}

The install metadata in \texttt{pyproject.toml} and \texttt{setup.py} is
therefore the source of truth for CI dependencies. \texttt{PyYAML} is
declared there because the maintained configuration loaders import
\texttt{yaml} for local and production dataset configuration files.

\hypertarget{optional-external-track-setup}{%
\subsubsection{Optional External TRACK
Setup}\label{optional-external-track-setup}}

Some workflows require the external TRACK executable, which is not
committed to this repository. See \texttt{docs/TRACK\_SETUP.md} after
the Python package is installed.

\hypertarget{verify-console-commands}{%
\subsubsection{Verify Console Commands}\label{verify-console-commands}}

After installation, these package commands should be available inside
the activated virtual environment:

\begin{Shaded}
\begin{Highlighting}[]
\ExtensionTok{synapticTrack} \AttributeTok{{-}{-}help}
\ExtensionTok{syntrack{-}track{-}check} \AttributeTok{{-}{-}help}
\ExtensionTok{syntrack{-}lebt{-}preflight} \AttributeTok{{-}{-}help}
\ExtensionTok{syntrack{-}lebt{-}validate} \AttributeTok{{-}{-}help}
\end{Highlighting}
\end{Shaded}

\hypertarget{development-and-agent-operating-instructions}{%
\section{Development and Agent Operating
Instructions}\label{development-and-agent-operating-instructions}}

Source Markdown: \texttt{CODEX.md}

\hypertarget{synaptictrack-1}{%
\subsection{synapticTrack}\label{synaptictrack-1}}

\hypertarget{overview}{%
\subsubsection{Overview}\label{overview}}

synapticTrack is a Python framework for accelerator modeling, beam
dynamics analysis, machine learning, optimization, and accelerator
control applications.

The framework is intended to support both simulation and future online
accelerator operation.

Primary application areas include:

\begin{itemize}
\tightlist
\item
  Beam dynamics simulations
\item
  Accelerator lattice modeling
\item
  Beam diagnostics analysis
\item
  Scanner data processing
\item
  Machine learning surrogate modeling
\item
  Reinforcement learning based control
\item
  Multi-objective optimization
\item
  Bayesian optimization
\item
  Digital twin development
\item
  EPICS/TANGO integration
\item
  Accelerator AI applications
\end{itemize}

Current accelerator applications include:

\begin{itemize}
\tightlist
\item
  RAON LEBT orbit correction
\item
  RAON MEBT studies
\item
  Electron injector optimization
\item
  Space-charge studies
\item
  Beam phase-space reconstruction
\item
  Accelerator tuning and control
\end{itemize}

\begin{center}\rule{0.5\linewidth}{0.5pt}\end{center}

\hypertarget{software-architecture}{%
\subsubsection{Software Architecture}\label{software-architecture}}

\hypertarget{core-modules}{%
\paragraph{Core Modules}\label{core-modules}}

\hypertarget{beam}{%
\subparagraph{beam}\label{beam}}

Beam representations:

\begin{itemize}
\tightlist
\item
  Beam
\item
  BeamWS
\item
  BeamAS
\item
  Twiss
\end{itemize}

Responsibilities:

\begin{itemize}
\tightlist
\item
  Particle coordinates
\item
  Beam statistics
\item
  Centroids
\item
  RMS sizes
\item
  Emittances
\item
  Twiss parameters
\end{itemize}

\begin{center}\rule{0.5\linewidth}{0.5pt}\end{center}

\hypertarget{io}{%
\subparagraph{io}\label{io}}

Input/output interfaces:

\begin{itemize}
\tightlist
\item
  TRACK
\item
  JuTrack
\item
  OPAL
\item
  FLAME
\item
  User-defined formats
\end{itemize}

Responsibilities:

\begin{itemize}
\tightlist
\item
  Read beam files
\item
  Write beam files
\item
  Conversion between simulation codes
\end{itemize}

\begin{center}\rule{0.5\linewidth}{0.5pt}\end{center}

\hypertarget{lattice}{%
\subparagraph{lattice}\label{lattice}}

Accelerator lattice definitions.

Responsibilities:

\begin{itemize}
\tightlist
\item
  Element models
\item
  TRACK lattice generation
\item
  Lattice manipulation
\item
  Future MAD-X support
\end{itemize}

\begin{center}\rule{0.5\linewidth}{0.5pt}\end{center}

\hypertarget{analysis}{%
\subparagraph{analysis}\label{analysis}}

Beam diagnostics and data analysis.

Responsibilities:

\begin{itemize}
\tightlist
\item
  Wire scanner analysis
\item
  Allison scanner analysis
\item
  Centroid extraction
\item
  RMS calculations
\item
  Emittance calculations
\item
  Beam quality metrics
\end{itemize}

\begin{center}\rule{0.5\linewidth}{0.5pt}\end{center}

\hypertarget{visualization}{%
\subparagraph{visualization}\label{visualization}}

Plotting and diagnostics.

Responsibilities:

\begin{itemize}
\tightlist
\item
  Phase-space plots
\item
  Beam distributions
\item
  Orbit plots
\item
  Scanner plots
\item
  Optimization histories
\end{itemize}

\begin{center}\rule{0.5\linewidth}{0.5pt}\end{center}

\hypertarget{optimization}{%
\subparagraph{optimization}\label{optimization}}

Optimization algorithms.

Current and planned methods:

\begin{itemize}
\tightlist
\item
  Differential Evolution
\item
  Genetic Algorithms
\item
  Bayesian Optimization
\item
  Multi-objective Optimization
\item
  Reinforcement Learning
\end{itemize}

\begin{center}\rule{0.5\linewidth}{0.5pt}\end{center}

\hypertarget{machine_learning}{%
\subparagraph{machine\_learning}\label{machine_learning}}

Data-driven accelerator modeling.

Current methods:

\begin{itemize}
\tightlist
\item
  Fully Connected Neural Networks
\item
  Surrogate Models
\end{itemize}

Planned methods:

\begin{itemize}
\tightlist
\item
  Physics-Informed Neural Networks
\item
  Graph Neural Networks
\item
  Transformer Models
\item
  Differentiable Beam Dynamics
\end{itemize}

\begin{center}\rule{0.5\linewidth}{0.5pt}\end{center}

\hypertarget{reinforcement_learning}{%
\subparagraph{reinforcement\_learning}\label{reinforcement_learning}}

Closed-loop accelerator control.

Current methods:

\begin{itemize}
\tightlist
\item
  PPO
\end{itemize}

Planned methods:

\begin{itemize}
\tightlist
\item
  SAC
\item
  TD3
\item
  DDPG
\item
  Model-based RL
\end{itemize}

Applications:

\begin{itemize}
\tightlist
\item
  Orbit correction
\item
  Beam tuning
\item
  Autonomous operation
\end{itemize}

\begin{center}\rule{0.5\linewidth}{0.5pt}\end{center}

\hypertarget{development-principles}{%
\subsubsection{Development Principles}\label{development-principles}}

\hypertarget{avoid-circular-imports}{%
\paragraph{Avoid Circular Imports}\label{avoid-circular-imports}}

Do not import high-level modules inside low-level modules.

Preferred dependency direction:

beam → io → analysis → visualization

optimization → machine\_learning → reinforcement\_learning

Avoid:

analysis → reinforcement\_learning

reinforcement\_learning → analysis

Use local imports when necessary.

\begin{center}\rule{0.5\linewidth}{0.5pt}\end{center}

\hypertarget{public-apis}{%
\paragraph{Public APIs}\label{public-apis}}

Do not break existing public APIs unless explicitly requested.

Maintain backward compatibility whenever possible.

\begin{center}\rule{0.5\linewidth}{0.5pt}\end{center}

\hypertarget{testing}{%
\paragraph{Testing}\label{testing}}

Run before every commit:

python -m pytest tests -v

For import-related changes:

python -c ``import synapticTrack''

python -c ``import synapticTrack.utils.sm\_utils''

python -c ``import synapticTrack.utils.rl\_utils''

\begin{center}\rule{0.5\linewidth}{0.5pt}\end{center}

\hypertarget{machine-learning-conventions}{%
\paragraph{Machine Learning
Conventions}\label{machine-learning-conventions}}

RAON LEBT orbit correction:

Inputs:

{[}FH1,FH2,FH3,FH4,FH5, FV1,FV2,FV3,FV4,FV5, alpha1,alpha2,alpha3{]}

Outputs:

{[}WS1\_x, AS\_x, WS2\_x, WS3\_x, WS4\_x, END\_x, WS1\_y, AS\_y, WS2\_y,
WS3\_y, WS4\_y, END\_y{]}

Steering magnet units:

mrad

Alpha parameters:

dimensionless scaling factors

\begin{center}\rule{0.5\linewidth}{0.5pt}\end{center}

\hypertarget{code-quality}{%
\paragraph{Code Quality}\label{code-quality}}

Prefer:

\begin{itemize}
\tightlist
\item
  Small functions
\item
  Unit tests
\item
  Explicit typing
\item
  Clear documentation
\end{itemize}

Avoid:

\begin{itemize}
\tightlist
\item
  Hidden global state
\item
  Wildcard imports
\item
  Circular dependencies
\item
  Hard-coded paths
\end{itemize}

Source Markdown: \texttt{AGENTS.md}

\hypertarget{agents.md-synaptictrack-aiml-development}{%
\subsection{AGENTS.md --- synapticTrack AI/ML
Development}\label{agents.md-synaptictrack-aiml-development}}

\hypertarget{scope}{%
\subsubsection{Scope}\label{scope}}

Implement the staged 2026 program for LEBT surrogate modeling, Gaussian
beam-state inference, MEBT surrogate modeling, RL orbit correction, and
later advanced reconstruction.

\hypertarget{required-principles}{%
\subsubsection{Required principles}\label{required-principles}}

\begin{itemize}
\tightlist
\item
  Use PyTorch for ML and differentiable inference.
\item
  The core initial transverse distribution is an uncoupled 4D Gaussian.
\item
  Do not add a separate shape latent vector to the core model.
\item
  LEBT inputs: Twiss parameters, emittances, offsets, angles, current,
  steering strengths, and electrostatic-quadrupole strengths.
\item
  LEBT outputs: centroids, rms sizes, and beam states at the last WS and
  LEBT exit.
\item
  Build the MEBT surrogate before RL orbit correction.
\item
  Include injected beam state in the MEBT surrogate inputs.
\item
  Put reusable logic in Python modules; notebooks import package code.
\end{itemize}

\hypertarget{scientific-rules}{%
\subsubsection{Scientific rules}\label{scientific-rules}}

\begin{itemize}
\tightlist
\item
  State units explicitly and reuse existing synapticTrack conventions.
\item
  Validate positive beta, positive emittance, and non-negative current.
\item
  Preserve covariance positive definiteness.
\item
  Store failed/lost-beam cases with validity masks and reasons.
\item
  Use deterministic seeds and record code/configuration metadata.
\end{itemize}

\hypertarget{testing-1}{%
\subsubsection{Testing}\label{testing-1}}

Every milestone adds tests for relevant physics, tensor shapes,
dtype/device behavior, determinism, serialization, bounds, reward
behavior, and termination conditions. Run the relevant existing tests
before completing a task.

\hypertarget{codex-behavior}{%
\subsubsection{Codex behavior}\label{codex-behavior}}

\begin{itemize}
\tightlist
\item
  Inspect existing code before adding modules.
\item
  Reuse existing beam and backend abstractions.
\item
  Implement only the requested milestone.
\item
  Avoid unrelated refactors.
\item
  Report assumptions, changed files, commands, test outcomes, and
  blockers.
\end{itemize}

\hypertarget{repository-operations}{%
\subsubsection{Repository operations}\label{repository-operations}}

\begin{itemize}
\tightlist
\item
  Do not run \texttt{git\ push}.
\item
  Do not check GitHub Actions or remote CI status.
\item
  Local GitHub Actions workflow files may be modified when requested or
  needed, but only in the local repository.
\end{itemize}

\hypertarget{reporting-format}{%
\subsubsection{Reporting format}\label{reporting-format}}

Every completed task summary must explicitly include:

\begin{itemize}
\tightlist
\item
  assumptions;
\item
  changed files;
\item
  commands run;
\item
  test or validation outcomes;
\item
  blockers or limitations;
\item
  confirmation that no \texttt{git\ push} was run and remote GitHub
  Actions were not checked.
\end{itemize}

\hypertarget{project-attribution}{%
\subsubsection{Project attribution}\label{project-attribution}}

\begin{itemize}
\tightlist
\item
  Author: Chong Shik Park.
\item
  Affiliation: Department of Accelerator Science and Center for
  Accelerator Research, Korea University, Sejong, 30019 Republic of
  Korea.
\end{itemize}

\hypertarget{unclassified-markdown-sources}{%
\section{Unclassified Markdown
Sources}\label{unclassified-markdown-sources}}

Source Markdown: \texttt{docs/beam\_state\_physics\_api.md}

\hypertarget{beam-state-physics-api}{%
\subsection{Beam-State Physics API}\label{beam-state-physics-api}}

\texttt{synapticTrack.beam.state\_physics} is the canonical
implementation for the LEBT/MEBT 10D beam state:

\begin{Shaded}
\begin{Highlighting}[]
\NormalTok{x\_c, xp\_c, y\_c, yp\_c, \textless{}x\^{}2\textgreater{}, \textless{}xx\textquotesingle{}\textgreater{}, \textless{}x\textquotesingle{}\^{}2\textgreater{}, \textless{}y\^{}2\textgreater{}, \textless{}yy\textquotesingle{}\textgreater{}, \textless{}y\textquotesingle{}\^{}2\textgreater{}}
\end{Highlighting}
\end{Shaded}

The module centralizes component units and determinant tolerances. Use
\texttt{state10\_physics(state,\ mode="differentiable")} in PyTorch
losses: it applies one finite determinant floor to preserve gradients.
Use \texttt{state10\_physics(state,\ mode="strict")} for validation and
publication metrics: invalid covariance-derived quantities remain
\texttt{NaN} and \texttt{validity} reports finite-input, determinant,
emittance, beta, and overall validity flags.

\texttt{state10\_from\_cholesky\_parameters} provides the explicit
full-state positive-definite reconstruction mode;
\texttt{covariance\_from\_cholesky\_params} and related functions remain
available for individual planes. New code should import this module from
\texttt{synapticTrack.beam}. \texttt{synapticTrack.ml.beam\_physics},
\texttt{ml.metrics.moments\_to\_twiss},
\texttt{data.production\_audit.moments\_to\_twiss}, and
\texttt{simulation.dataset\_generator.derive\_twiss\_from\_state\_moments}
remain compatibility wrappers and should not be used in new code.

Publication and audit reports must not clamp invalid determinants. They
retain invalid values as \texttt{NaN}, count them explicitly, and omit
them from derived error aggregates while recording \texttt{n\_invalid}.

\hypertarget{metric-reports}{%
\subsubsection{Metric Reports}\label{metric-reports}}

\texttt{compute\_state\_physics\_metrics()} in
\texttt{synapticTrack.ml.metrics} is the canonical strict reporting API
for \texttt{state\_last\_ws} and \texttt{state\_lebt\_exit}. Its
per-group output includes component and normalized errors,
physical-class metrics, derived Courant-Snyder/emittance errors,
covariance-validity fractions, warnings, sample counts, and the recorded
tolerance configuration.

Physics-regularized training writes a compatibility
\texttt{physics\_metrics\_\textless{}split\textgreater{}.json} report in
addition to canonical metrics. It preserves historical flat fields, but
new consumers should read the canonical \texttt{groups} report or
\texttt{state\_physics\_metrics\_\textless{}split\textgreater{}.json}.
Mixed-unit state RMSE is secondary and must not be interpreted as a
physical acceptance criterion.

Source Markdown:
\texttt{docs/codex\_refactor\_tasks/01\_fixed\_parameter\_inference\_contract.md}

\hypertarget{codex-prompt-fixed-parameter-inference-contract}{%
\subsection{Codex Prompt: Fixed-Parameter Inference
Contract}\label{codex-prompt-fixed-parameter-inference-contract}}

Read \texttt{AGENTS.md}, the current LEBT no-offset schema, surrogate
API, and tests.

\hypertarget{task}{%
\subsubsection{Task}\label{task}}

Refactor full 24D LEBT surrogate input validation so it checks each
fixed parameter against its schema-defined fixed value rather than
assuming every fixed parameter is zero.

\texttt{beam\_current} is fixed at \texttt{0.05} in the current
no-offset production schema; the injection coordinates remain zero.

\hypertarget{requirements}{%
\subsubsection{Requirements}\label{requirements}}

\begin{enumerate}
\def\labelenumi{\arabic{enumi}.}
\tightlist
\item
  Replace the public
  \texttt{fixed\_zero\_tol}/\texttt{allow\_nonzero\_fixed} terminology
  with names that reflect validation against fixed expected values.
  Preserve a backward-compatible API where practical.
\item
  Keep accepting the sampled 19D input directly.
\item
  For full 24D inputs, compare \texttt{x0}, \texttt{xp0}, \texttt{y0},
  \texttt{yp0}, and \texttt{beam\_current} against
  \texttt{LEBT\_FIXED\_PARAMETER\_VALUES} using a configurable
  tolerance.
\item
  Reject an input whose fixed current differs from \texttt{0.05}, while
  accepting a valid full row with that current by default.
\item
  Ensure errors identify the parameter, expected value, observed maximum
  error, and tolerance.
\item
  Reuse the same fixed-value helper in production audit code where
  possible; do not maintain two independent fixed-value rules.
\item
  Update affected notebooks and docs that still say ``fixed zero'' when
  they mean fixed injection/current.
\end{enumerate}

\hypertarget{tests}{%
\subsubsection{Tests}\label{tests}}

Add or update tests for valid full 24D input with current \texttt{0.05},
rejection of incorrect current, rejection of nonzero orbit offsets,
direct 19D input, and backward-compatible public arguments.

Do not modify generated datasets or retrain models.

Source Markdown:
\texttt{docs/codex\_refactor\_tasks/02\_beam\_current\_unit\_contract.md}

\hypertarget{codex-prompt-beam-current-unit-contract}{%
\subsection{Codex Prompt: Beam-Current Unit
Contract}\label{codex-prompt-beam-current-unit-contract}}

Read \texttt{AGENTS.md}, TRACK input handling, Gaussian beam-state code,
LEBT/MEBT production configs, and relevant documentation.

\hypertarget{task-1}{%
\subsubsection{Task}\label{task-1}}

Create one explicit, tested beam-current unit contract between
synapticTrack and TRACK.

The code currently labels beam current as \texttt{uA}, while TRACK
templates contain \texttt{current\ =\ 0.05}, and injection writes values
directly to \texttt{track.dat}.

\hypertarget{requirements-1}{%
\subsubsection{Requirements}\label{requirements-1}}

\begin{enumerate}
\def\labelenumi{\arabic{enumi}.}
\tightlist
\item
  Verify TRACK's \texttt{current} unit from project documentation or the
  TRACK reference available locally. If it cannot be verified locally,
  do not guess: define a clearly marked configuration boundary and
  document the unresolved external assumption.
\item
  Define canonical synapticTrack and TRACK units in one reusable module.
\item
  Add explicit conversion functions with finite/non-negative validation.
\item
  Make \texttt{update\_track\_dat\_twiss()} use the conversion rather
  than copying a raw value.
\item
  Ensure \texttt{GaussianBeamState}, \texttt{Beam},
  \texttt{DEFAULT\_LEBT\_UNITS}, \texttt{TRACK\_BEAM\_METADATA}, LEBT
  production, and MEBT production agree on the contract.
\item
  Store both numeric value and unit in generated metadata. Preserve
  existing files and provide compatibility handling for metadata without
  units.
\item
  Update documentation and paper-facing text only after the unit is
  verified.
\end{enumerate}

\hypertarget{tests-1}{%
\subsubsection{Tests}\label{tests-1}}

Test conversion round trips, TRACK-file replacement, non-negative
rejection, LEBT fixed current, MEBT fixed current, and output-beam
metadata propagation.

Do not silently reinterpret existing production data.

Source Markdown:
\texttt{docs/codex\_refactor\_tasks/03\_mebt\_surrogate\_input\_schema.md}

\hypertarget{codex-prompt-mebt-surrogate-input-schema}{%
\subsection{Codex Prompt: MEBT Surrogate Input
Schema}\label{codex-prompt-mebt-surrogate-input-schema}}

Read \texttt{AGENTS.md}, MEBT Sobol production code, LEBT schema
abstractions, and the MEBT notebook workflow.

\hypertarget{task-2}{%
\subsubsection{Task}\label{task-2}}

Define reusable MEBT dataset and ML schemas that include injected beam
state in the surrogate input contract.

The current MEBT Sobol workflow samples 19 lattice controls and starts
from a fixed \texttt{scratch.\#07} distribution. That is useful for a
fixed-distribution diagnostic study, but not for the stated MEBT
surrogate requirement.

\hypertarget{requirements-2}{%
\subsubsection{Requirements}\label{requirements-2}}

\begin{enumerate}
\def\labelenumi{\arabic{enumi}.}
\tightlist
\item
  Keep the current fixed-distribution 19D dataset format valid and
  explicitly identify it as such.
\item
  Add a versioned MEBT schema for a general surrogate with:

  \begin{itemize}
  \tightlist
  \item
    lattice controls;
  \item
    injected transverse beam state, using the established 10D
    state-moment convention or a documented equivalent physical
    parameterization;
  \item
    fixed-current value and unit metadata;
  \item
    validity masks and failure reasons.
  \end{itemize}
\item
  Reuse \texttt{DatasetSchema}/ML schema patterns rather than defining a
  notebook-only format.
\item
  Define target groups and units, including diagnostics and any
  requested MEBT beam states.
\item
  Separate sampling bounds for controls and injected state, validating
  positive beta/emittance and positive-definite covariance.
\item
  Document how \texttt{scratch.\#07} becomes a nominal reference
  distribution versus a per-sample injected state.
\item
  Do not change the 44D LEBT output format or run a production campaign.
\end{enumerate}

\hypertarget{tests-2}{%
\subsubsection{Tests}\label{tests-2}}

Test schema dimensions/names, deterministic sample expansion, valid
state acceptance, invalid covariance rejection, metadata serialization,
and compatibility loading of the current fixed-distribution MEBT
batches.

Source Markdown:
\texttt{docs/codex\_refactor\_tasks/04\_physics\_training\_output\_scaling.md}

\hypertarget{codex-prompt-physics-regularized-output-scaling}{%
\subsection{Codex Prompt: Physics-Regularized Output
Scaling}\label{codex-prompt-physics-regularized-output-scaling}}

Read \texttt{AGENTS.md}, Task 4/Task 5 training modules, normalizers,
loss functions, and the physics-regularized training configuration.

\hypertarget{task-3}{%
\subsubsection{Task}\label{task-3}}

Make output scaling a real, reproducible part of physics-regularized
training.

The current physics model trains against physical-space targets
directly, even though output normalizers are fitted and
\texttt{output\_normalization} is configurable. This leaves mixed
physical scales to dominate data and derived losses.

\hypertarget{requirements-3}{%
\subsubsection{Requirements}\label{requirements-3}}

\begin{enumerate}
\def\labelenumi{\arabic{enumi}.}
\tightlist
\item
  State clearly which quantities are predicted in normalized coordinates
  and which physics losses must be evaluated after conversion to
  physical units.
\item
  Introduce a reusable differentiable output transform based on
  train-split statistics. It must work on CPU/GPU and preserve
  dtype/device.
\item
  Apply scaling consistently to diagnostics, state components,
  Courant-Snyder loss, covariance penalties, and evaluation output.
\item
  Make \texttt{output\_normalization} either operational or remove it
  with a documented migration path; it must not remain a no-op
  configuration field.
\item
  Preserve checkpoint loading and the public
  \texttt{predict\_physical()} API.
\item
  Record output scaling values, source split, and transform version in
  model artifacts.
\item
  Avoid changing model architecture or retraining existing production
  models.
\end{enumerate}

\hypertarget{tests-3}{%
\subsubsection{Tests}\label{tests-3}}

Test normalizer round trips, differentiable transforms, physical-space
metric consistency, component balancing for disparate scales, checkpoint
reload, and deterministic small training runs.

Source Markdown:
\texttt{docs/codex\_refactor\_tasks/05\_unify\_beam\_state\_physics\_api.md}

\hypertarget{codex-prompt-unify-beam-state-physics-api}{%
\subsection{Codex Prompt: Unify Beam-State Physics
API}\label{codex-prompt-unify-beam-state-physics-api}}

Read \texttt{AGENTS.md} and the current beam-physics, metrics,
production-audit, physics-training, inference, and beam modules.

\hypertarget{task-4}{%
\subsubsection{Task}\label{task-4}}

Consolidate the duplicated 10D state-moment, Courant-Snyder, emittance,
and covariance-validity calculations into one documented API supporting
NumPy and PyTorch.

\hypertarget{requirements-4}{%
\subsubsection{Requirements}\label{requirements-4}}

\begin{enumerate}
\def\labelenumi{\arabic{enumi}.}
\tightlist
\item
  Establish one canonical state order:
  \texttt{x\_c,\ xp\_c,\ y\_c,\ yp\_c,\ \textless{}x\^{}2\textgreater{},\ \textless{}xx\textquotesingle{}\textgreater{},\ \textless{}x\textquotesingle{}\^{}2\textgreater{},\ \textless{}y\^{}2\textgreater{},\ \textless{}yy\textquotesingle{}\textgreater{},\ \textless{}y\textquotesingle{}\^{}2\textgreater{}}.
\item
  Define explicit modes for:

  \begin{itemize}
  \tightlist
  \item
    differentiable training calculations, with safe finite gradients;
  \item
    strict validation/reporting, with invalid values represented as
    \texttt{NaN} and associated validity flags;
  \item
    positive-definite Cholesky reconstruction.
  \end{itemize}
\item
  Centralize determinant thresholds and unit definitions. Do not let
  each caller choose unrelated defaults.
\item
  Migrate \texttt{ml.metrics}, \texttt{data.production\_audit},
  \texttt{ml.physics\_training}, and surrogate inference to the shared
  implementation.
\item
  Preserve existing metric JSON keys where needed, but deprecate
  duplicate helpers and document replacements.
\item
  Ensure derived metrics count non-finite values explicitly and never
  produce silently misleading aggregate errors.
\end{enumerate}

\hypertarget{tests-4}{%
\subsubsection{Tests}\label{tests-4}}

Add cross-module consistency tests for valid/invalid covariances,
NumPy/Torch agreement, gradients, all 10D components, both LEBT state
groups, and JSON-safe reporting.

Do not change the physical definitions or silently clamp invalid values
in publication metrics.

Source Markdown:
\texttt{docs/codex\_refactor\_tasks/06\_legacy\_physics\_cleanup.md}

\hypertarget{codex-prompt-legacy-physics-and-runtime-cleanup}{%
\subsection{Codex Prompt: Legacy Physics and Runtime
Cleanup}\label{codex-prompt-legacy-physics-and-runtime-cleanup}}

Read \texttt{AGENTS.md}, legacy beam utilities, TRACK runtime helpers,
and their tests.

\hypertarget{task-5}{%
\subsubsection{Task}\label{task-5}}

Repair known legacy correctness issues and remove brittle assumptions
while keeping public behavior compatible.

\hypertarget{requirements-5}{%
\subsubsection{Requirements}\label{requirements-5}}

\begin{enumerate}
\def\labelenumi{\arabic{enumi}.}
\tightlist
\item
  Fix \texttt{beam.twiss.Twiss.compute\_twiss()} so its self-consistency
  warning does not reference undefined \texttt{label}.
\item
  Route legacy validity summaries through the canonical state-physics
  API from the beam-state refactor. Non-finite beta/emittance values
  must be counted as invalid, not omitted by comparisons with
  \texttt{NaN}.
\item
  Make MEBT lattice segment discovery order numeric suffixes, so
  \texttt{elements10.json} follows \texttt{elements9.json}.
\item
  Replace LEBT-specific TRACK runtime defaults such as
  \texttt{scratch.\#02} with a small runtime specification object or
  validated arguments. Keep existing defaults for callers.
\item
  Replace print-and-continue output-file handling with structured
  warnings or results where failures affect downstream correctness.
\item
  Add focused regression tests for each repaired path.
\end{enumerate}

Do not remove legacy public functions abruptly. Keep compatibility shims
and mark them deprecated where appropriate.

Source Markdown:
\texttt{docs/codex\_refactor\_tasks/07\_track\_execution\_orchestration.md}

\hypertarget{codex-prompt-unify-track-execution-orchestration}{%
\subsection{Codex Prompt: Unify TRACK Execution
Orchestration}\label{codex-prompt-unify-track-execution-orchestration}}

Read \texttt{AGENTS.md}, \texttt{utils/track\_runtime.py},
\texttt{track/run\_track.py},
\texttt{simulation/lebt\_track\_adapter.py},
\texttt{simulation/mebt\_sobol\_production.py}, and their tests.

\hypertarget{task-6}{%
\subsubsection{Task}\label{task-6}}

Create one reusable, path-safe TRACK segmented-execution service for
LEBT and MEBT. Remove global-current-working-directory dependence from
legacy execution while preserving existing public functions as
compatibility wrappers.

\hypertarget{motivation}{%
\subsubsection{Motivation}\label{motivation}}

\texttt{run\_calibrated\_track()} changes the process-wide working
directory, while the LEBT adapter and MEBT Sobol producer independently
prepare runtimes, write segments, hand off scratch distributions,
collect outputs, and clean work directories. This makes concurrent
execution and failure handling brittle.

\hypertarget{requirements-6}{%
\subsubsection{Requirements}\label{requirements-6}}

\begin{enumerate}
\def\labelenumi{\arabic{enumi}.}
\tightlist
\item
  Add a typed execution request/result API in a reusable runtime module.
  It must include source/work directories, \texttt{TrackRuntimeSpec},
  lattice segments, output directories, TRACK executable, verbosity, and
  the initial scratch index.
\item
  Execute every path explicitly. Do not call \texttt{os.chdir()} in
  library code.
\item
  Keep the existing scratch handoff contract: a segment loading
  \texttt{scratch.\#NN} must make \texttt{scratch.\#(NN+1)} available
  for the next segment.
\item
  Centralize preparation, segmented execution, final-output collection,
  and cleanup policy. MEBT must retain \texttt{scratch.\#07} as its
  configured initial distribution; LEBT retains its default
  \texttt{scratch.\#02}.
\item
  Make \texttt{run\_calibrated\_track()} and
  \texttt{run\_track\_steps()} compatibility wrappers around the new
  path-safe API where practical. Do not break notebook callers.
\item
  Have the LEBT adapter and MEBT producer use the shared service rather
  than maintaining independent orchestration sequences.
\item
  Return structured execution data including segment outputs, final
  directory, scratch handoff records, and cleanup outcome. Preserve
  existing numerical diagnostic returns through wrappers.
\end{enumerate}

\hypertarget{tests-5}{%
\subsubsection{Tests}\label{tests-5}}

Add tests for:

\begin{enumerate}
\def\labelenumi{\arabic{enumi}.}
\tightlist
\item
  current working directory is unchanged on success and exceptions;
\item
  LEBT default initial scratch and MEBT configured initial scratch;
\item
  multi-segment scratch handoff in numeric order;
\item
  absolute and relative output paths;
\item
  cleanup policies for successful and failed runs;
\item
  compatibility of existing \texttt{run\_calibrated\_track()} and
  adapter APIs.
\end{enumerate}

Do not invoke the real TRACK executable in tests. Do not retrain models
or change dataset formats.

Source Markdown:
\texttt{docs/codex\_refactor\_tasks/08\_track\_diagnostic\_output\_contract.md}

\hypertarget{codex-prompt-enforce-track-diagnostic-output-contracts}{%
\subsection{Codex Prompt: Enforce TRACK Diagnostic Output
Contracts}\label{codex-prompt-enforce-track-diagnostic-output-contracts}}

Read \texttt{AGENTS.md}, \texttt{utils/track\_outputs.py}, TRACK output
collection code, LEBT/MEBT production schemas, and their tests.

\hypertarget{task-7}{%
\subsubsection{Task}\label{task-7}}

Define one explicit, validated contract for segmented TRACK diagnostic
outputs and use it for LEBT and MEBT collection.

\hypertarget{motivation-1}{%
\subsubsection{Motivation}\label{motivation-1}}

Current helpers use \texttt{zip()} and positional indexing when copying
coordinate files and combining \texttt{z\_scanner.out}, centroids, and
RMS values. A missing or extra path can silently truncate collection or
produce a misaligned diagnostic vector.

\hypertarget{requirements-7}{%
\subsubsection{Requirements}\label{requirements-7}}

\begin{enumerate}
\def\labelenumi{\arabic{enumi}.}
\tightlist
\item
  Add typed diagnostic-location and output-spec objects containing the
  diagnostic name, segment output path, final coordinate filename, and
  unit.
\item
  Validate exact cardinality, unique final filenames, readable
  coordinate files, and exact z-position alignment before creating
  output vectors.
\item
  Replace positional \texttt{zip()} truncation in output-copy helpers
  with explicit length checks and informative errors.
\item
  Validate output arrays are finite and all RMS quantities are positive
  for samples marked valid. Preserve lost/failed samples with masks and
  reasons.
\item
  Return a structured diagnostic collection result with source paths, z
  positions, centroids, RMS values, units, and validation findings.
\item
  Retain \texttt{get\_beam\_data()},
  \texttt{diagnostics\_to\_output\_vector()}, and existing copy helpers
  as documented compatibility wrappers.
\item
  Reuse schema names and units instead of duplicating ordered strings in
  LEBT and MEBT producers.
\end{enumerate}

\hypertarget{tests-6}{%
\subsubsection{Tests}\label{tests-6}}

Add tests for:

\begin{enumerate}
\def\labelenumi{\arabic{enumi}.}
\tightlist
\item
  valid LEBT and MEBT diagnostic layouts;
\item
  unequal segment/output/final-file lengths;
\item
  duplicate final filenames;
\item
  missing coordinate output and missing/misaligned z-scanner entries;
\item
  nonfinite centroid and nonpositive RMS handling;
\item
  legacy helper output equivalence for valid inputs.
\end{enumerate}

Do not run real TRACK or modify generated datasets.

Source Markdown:
\texttt{docs/codex\_refactor\_tasks/09\_production\_batch\_atomicity\_and\_failure\_taxonomy.md}

\hypertarget{codex-prompt-atomic-production-batches-and-failure-taxonomy}{%
\subsection{Codex Prompt: Atomic Production Batches and Failure
Taxonomy}\label{codex-prompt-atomic-production-batches-and-failure-taxonomy}}

Read \texttt{AGENTS.md}, LEBT/MEBT Sobol production modules, batch
planning, HDF5 I/O, generation runtime helpers, and tests.

\hypertarget{task-8}{%
\subsubsection{Task}\label{task-8}}

Make staged LEBT and MEBT sample production resumable and auditable
under worker, TRACK, and filesystem failures.

\hypertarget{motivation-2}{%
\subsubsection{Motivation}\label{motivation-2}}

MEBT workers currently serialize broad exceptions as unstructured
strings and write batch HDF5 files directly. Interrupted writes,
malformed partial batches, and different failure representations across
LEBT/MEBT can corrupt resume and audit decisions.

\hypertarget{requirements-8}{%
\subsubsection{Requirements}\label{requirements-8}}

\begin{enumerate}
\def\labelenumi{\arabic{enumi}.}
\tightlist
\item
  Define a reusable typed sample outcome with global index, validity
  mask, failure category, human-readable reason, optional exception
  class, runtime, and work-directory disposition.
\item
  Establish a small stable failure taxonomy, including input validation,
  runtime preparation, TRACK process, required output, parsing, and
  internal errors. Preserve original messages as context.
\item
  Write batches to a same-filesystem temporary path, flush/close
  successfully, then atomically replace the final HDF5 path. Set
  \texttt{batch\_completed} only after durable completion.
\item
  Make resume/audit reject partial, unreadable, schema-incomplete, or
  noncontiguous batches with actionable findings; never treat them as
  complete.
\item
  Keep valid/lost/failed rows and reasons in the dataset rather than
  silently dropping them. Maintain deterministic Sobol index ordering in
  serial and multiprocess paths.
\item
  Preserve existing HDF5 fields and add new metadata/fields compatibly
  with explicit schema-version handling.
\item
  Share the outcome and batch-commit logic between LEBT and MEBT where
  their contracts overlap.
\end{enumerate}

\hypertarget{tests-7}{%
\subsubsection{Tests}\label{tests-7}}

Add tests for:

\begin{enumerate}
\def\labelenumi{\arabic{enumi}.}
\tightlist
\item
  atomic success commit and absence of temporary files afterward;
\item
  simulated write failure leaves no completed final batch;
\item
  corrupt or incomplete batch rejection during resume;
\item
  every failure category and JSON/HDF5-safe reason storage;
\item
  serial/multiprocess ordering equivalence using a fake backend;
\item
  backward-compatible reading of existing completed batches.
\end{enumerate}

Do not invoke TRACK, delete user production data, or retrain models.

Source Markdown:
\texttt{docs/codex\_refactor\_tasks/10\_versioned\_artifact\_and\_config\_loading.md}

\hypertarget{codex-prompt-versioned-artifact-and-configuration-loading}{%
\subsection{Codex Prompt: Versioned Artifact and Configuration
Loading}\label{codex-prompt-versioned-artifact-and-configuration-loading}}

Read \texttt{AGENTS.md}, \texttt{ml/io\_utils.py}, production
configuration discovery, checkpoint/normalizer artifact loading, and
tests.

\hypertarget{task-9}{%
\subsubsection{Task}\label{task-9}}

Replace ambiguous silent artifact/config parsing with explicit loading
results and numeric version-aware discovery, while preserving the
current convenience APIs.

\hypertarget{motivation-3}{%
\subsubsection{Motivation}\label{motivation-3}}

Several readers return \texttt{None} or an empty list for both a missing
file and a malformed file. Versioned configuration discovery also relies
on lexical path sorting, which misorders versions such as \texttt{v10}
and \texttt{v2}. These behaviors can cause training or production to run
with fallback defaults without a visible root cause.

\hypertarget{requirements-9}{%
\subsubsection{Requirements}\label{requirements-9}}

\begin{enumerate}
\def\labelenumi{\arabic{enumi}.}
\tightlist
\item
  Add a reusable artifact-load result or typed exception contract that
  distinguishes missing, invalid syntax, invalid schema, unsupported
  version, and successful load.
\item
  Add strict reader variants for JSON, YAML, CSV, normalizer manifests,
  and production configs. Existing tolerant helpers must remain
  available and be documented as compatibility behavior.
\item
  Parse versioned filenames numerically and deterministically. Reject
  ambiguous names and duplicate logical versions with actionable errors.
\item
  Validate minimum required metadata before an artifact is consumed:
  schema version, dimensions where applicable, units/current-contract
  metadata, and source/config identity.
\item
  Record the resolved artifact paths and validation result in training
  and production manifests, without modifying historical artifact
  contents.
\item
  Remove broad exception swallowing where an invalid artifact would
  otherwise be mistaken for a missing optional artifact.
\end{enumerate}

\hypertarget{tests-8}{%
\subsubsection{Tests}\label{tests-8}}

Add tests for:

\begin{enumerate}
\def\labelenumi{\arabic{enumi}.}
\tightlist
\item
  missing versus malformed JSON/YAML/CSV;
\item
  numeric ordering of \texttt{v1}, \texttt{v2}, and \texttt{v10}
  configurations;
\item
  duplicate and malformed versioned filenames;
\item
  strict metadata rejection and tolerant compatibility behavior;
\item
  checkpoint/normalizer manifest path resolution;
\item
  JSON-safe load diagnostics.
\end{enumerate}

Do not change trained model weights, regenerate datasets, or require
remote resources.

Source Markdown:
\texttt{docs/codex\_refactor\_tasks/11\_mebt\_batch\_schema\_and\_fixed\_current\_contract.md}

\hypertarget{codex-prompt-mebt-batch-schema-and-fixed-current-contract}{%
\subsection{Codex Prompt: MEBT Batch Schema and Fixed-Current
Contract}\label{codex-prompt-mebt-batch-schema-and-fixed-current-contract}}

Read \texttt{AGENTS.md}, the MEBT Sobol producer,
\texttt{DatasetSchema}, batch I/O, MEBT schema tests, and existing LEBT
no-offset batch tests.

\hypertarget{task-10}{%
\subsubsection{Task}\label{task-10}}

Make the fixed-distribution MEBT Sobol batch format conform to one
versioned, schema-driven dataset contract, without changing existing
generated datasets.

The current producer writes a custom layout, while the declared MEBT
schema and generic batch tools describe a richer, incompatible layout.
The producer also accepts \texttt{fixed\_beam\_current} even though its
schema declares a fixed value of \texttt{0.05}.

\hypertarget{requirements-10}{%
\subsubsection{Requirements}\label{requirements-10}}

\begin{enumerate}
\def\labelenumi{\arabic{enumi}.}
\tightlist
\item
  Define a generic batch representation that supports:

  \begin{itemize}
  \tightlist
  \item
    sampled and full inputs;
  \item
    fixed parameters;
  \item
    optional target groups;
  \item
    validity, completion, failure reason, and production outcomes;
  \item
    per-field names and units from \texttt{DatasetSchema}.
  \end{itemize}
\item
  Make the fixed-distribution MEBT producer write the generic
  representation for new batches. Its 19 sampled controls and fixed
  current must be stored as an auditable full 20D input row.
\item
  Preserve the existing MEBT batch reader behavior through a
  compatibility adapter for legacy custom-layout files. Do not rewrite
  or delete data.
\item
  Make MEBT current derive from the fixed-distribution schema. Reject a
  production config whose current differs from the schema value, unless
  a new explicitly versioned schema and dataset name are supplied.
\item
  Preserve the current TRACK conversion boundary and unit-verification
  metadata. Do not claim a verified TRACK physical unit.
\item
  Make consolidated MEBT artifacts retain completion, validity, failure,
  and current metadata needed for downstream audit.
\end{enumerate}

\hypertarget{tests-9}{%
\subsubsection{Tests}\label{tests-9}}

Add tests for:

\begin{enumerate}
\def\labelenumi{\arabic{enumi}.}
\tightlist
\item
  a newly generated dry-run MEBT batch satisfying its schema;
\item
  full 20D controls-plus-fixed-current reconstruction;
\item
  rejection of a config current inconsistent with the schema;
\item
  legacy MEBT batch compatibility reading;
\item
  consolidated artifact preservation of validity and current metadata;
\item
  no regression in LEBT batch read/write tests.
\end{enumerate}

\hypertarget{do-not-do}{%
\subsubsection{Do Not Do}\label{do-not-do}}

\begin{itemize}
\tightlist
\item
  Do not regenerate production datasets.
\item
  Do not silently reinterpret current units.
\item
  Do not remove support for existing MEBT batch files.
\end{itemize}

\hypertarget{completion-summary}{%
\subsubsection{Completion Summary}\label{completion-summary}}

Report changed files, old/new batch layouts, current validation
behavior, compatibility behavior, tests run, and migration limitations.

Source Markdown:
\texttt{docs/codex\_refactor\_tasks/12\_mebt\_lattice\_control\_manifest.md}

\hypertarget{codex-prompt-mebt-lattice-control-manifest-and-drift-detection}{%
\subsection{Codex Prompt: MEBT Lattice-Control Manifest and Drift
Detection}\label{codex-prompt-mebt-lattice-control-manifest-and-drift-detection}}

Read \texttt{AGENTS.md}, the MEBT Sobol producer, MEBT schema, RAON MEBT
lattice files, lattice parsing code, and current MEBT production tests.

\hypertarget{task-11}{%
\subsubsection{Task}\label{task-11}}

Make the MEBT 19D control-vector contract explicit and detect lattice
changes that would otherwise silently change feature meaning while
preserving vector length.

\hypertarget{requirements-11}{%
\subsubsection{Requirements}\label{requirements-11}}

\begin{enumerate}
\def\labelenumi{\arabic{enumi}.}
\tightlist
\item
  Define a versioned lattice-control manifest containing, at minimum:

  \begin{itemize}
  \tightlist
  \item
    control names in canonical order;
  \item
    element type and stable lattice identifier where available;
  \item
    nominal value, unit, and sampled bounds;
  \item
    source lattice path and content hash.
  \end{itemize}
\item
  Generate the runtime control vector from this manifest rather than
  relying only on traversal order and synthetic
  \texttt{q01}/\texttt{cav01} labels.
\item
  Validate that the parsed lattice exactly matches the manifest in
  count, order, element identity/type, and control names before sample
  generation.
\item
  Store the resolved manifest and lattice hash in batch and consolidated
  metadata. Record the same information in the dataset-level config and
  metadata files.
\item
  Keep generated notebook lattice JSON as an optional input convenience,
  not the production source of truth.
\item
  Provide a clear error that identifies the first mismatched control and
  the expected versus observed lattice identity.
\end{enumerate}

\hypertarget{tests-10}{%
\subsubsection{Tests}\label{tests-10}}

Add tests for:

\begin{enumerate}
\def\labelenumi{\arabic{enumi}.}
\tightlist
\item
  the committed MEBT lattice matching the canonical manifest;
\item
  a reordered or substituted element causing validation failure;
\item
  stable nominal vector and bounds generation;
\item
  manifest/hash persistence in dry-run metadata;
\item
  fallback from notebook JSON to committed \texttt{sclinac.dat}
  retaining the same manifest.
\end{enumerate}

\hypertarget{do-not-do-1}{%
\subsubsection{Do Not Do}\label{do-not-do-1}}

\begin{itemize}
\tightlist
\item
  Do not change physical control ranges.
\item
  Do not alter \texttt{scratch.\#07} handling or TRACK segment ordering.
\item
  Do not regenerate datasets.
\end{itemize}

\hypertarget{completion-summary-1}{%
\subsubsection{Completion Summary}\label{completion-summary-1}}

Report the canonical control order, metadata fields, mismatch
protections, changed files, and tests run.

Source Markdown:
\texttt{docs/codex\_refactor\_tasks/13\_unify\_state\_physics\_metric\_reports.md}

\hypertarget{codex-prompt-unify-beam-state-physics-metric-reports}{%
\subsection{Codex Prompt: Unify Beam-State Physics Metric
Reports}\label{codex-prompt-unify-beam-state-physics-metric-reports}}

Read \texttt{AGENTS.md}, \texttt{ml/metrics.py},
\texttt{ml/physics\_training.py}, training summaries, publication-metric
code, and Task 4/Task 5 tests.

\hypertarget{task-12}{%
\subsubsection{Task}\label{task-12}}

Make \texttt{compute\_state\_physics\_metrics()} the sole implementation
of derived Courant-Snyder/emittance errors and covariance-validity
metrics used by plain MLP and physics-regularized MLP evaluation.

\hypertarget{requirements-12}{%
\subsubsection{Requirements}\label{requirements-12}}

\begin{enumerate}
\def\labelenumi{\arabic{enumi}.}
\tightlist
\item
  Remove duplicated calculation logic from physics-regularized training
  in favor of the canonical metric API.
\item
  Ensure both model types emit the same detailed report fields for each
  state group: component metrics, normalized metrics, derived metrics,
  validity fractions, warnings, and sample counts.
\item
  Preserve existing
  \texttt{physics\_metrics\_\textless{}split\textgreater{}.json} files
  through a documented compatibility serializer. Do not make notebooks
  or prior reports fail.
\item
  Make publication table generation consume the canonical metrics first;
  retain legacy fallback only for historical artifacts.
\item
  Define one tolerance configuration path for strict covariance
  reporting and record it in the emitted metric JSON.
\item
  Avoid mixed-unit aggregate RMSE as a physical acceptance metric.
\end{enumerate}

\hypertarget{tests-11}{%
\subsubsection{Tests}\label{tests-11}}

Add or update tests for:

\begin{enumerate}
\def\labelenumi{\arabic{enumi}.}
\tightlist
\item
  plain and physics-regularized evaluation yielding equivalent canonical
  metrics for identical predictions;
\item
  compatibility JSON containing historical required keys;
\item
  invalid covariance and nonfinite derived quantity reporting;
\item
  JSON serialization of all report fields;
\item
  publication tables loading both canonical and legacy artifacts.
\end{enumerate}

\hypertarget{do-not-do-2}{%
\subsubsection{Do Not Do}\label{do-not-do-2}}

\begin{itemize}
\tightlist
\item
  Do not retrain models.
\item
  Do not change loss definitions or model architecture.
\item
  Do not change saved metric values manually.
\end{itemize}

\hypertarget{completion-summary-2}{%
\subsubsection{Completion Summary}\label{completion-summary-2}}

Report canonical fields, compatibility fields, changed files, tests run,
and any historical-report limitations.

Source Markdown:
\texttt{docs/codex\_refactor\_tasks/14\_package\_import\_boundary\_cleanup.md}

\hypertarget{codex-prompt-package-import-boundary-cleanup}{%
\subsection{Codex Prompt: Package Import Boundary
Cleanup}\label{codex-prompt-package-import-boundary-cleanup}}

Read \texttt{AGENTS.md}, all package \texttt{\_\_init\_\_.py} files, the
previous circular import regression, CLI entry points, and
import-related tests.

\hypertarget{task-13}{%
\subsubsection{Task}\label{task-13}}

Reduce eager package re-exports so importing \texttt{synapticTrack} or a
low-level submodule does not transitively import unrelated analysis,
data generation, ML training, visualization, or TRACK runtime code.

\hypertarget{requirements-13}{%
\subsubsection{Requirements}\label{requirements-13}}

\begin{enumerate}
\def\labelenumi{\arabic{enumi}.}
\tightlist
\item
  Map the import graph around \texttt{synapticTrack}, \texttt{utils},
  \texttt{data}, \texttt{simulation}, and \texttt{ml} package
  initializers.
\item
  Keep only lightweight stable public types/functions eagerly imported.
  Require direct submodule imports for heavy workflows, or use narrowly
  scoped lazy exports where compatibility requires them.
\item
  Prevent data modules from importing simulation package initialization
  and prevent simulation configuration from importing the full ML
  package.
\item
  Preserve documented public imports where practical; issue deprecation
  guidance for imports that must move.
\item
  Add a fast import smoke test with a time budget for:

  \begin{itemize}
  \tightlist
  \item
    \texttt{import\ synapticTrack};
  \item
    beam/current and schema modules;
  \item
    MEBT schema and producer modules;
  \item
    ML metrics module.
  \end{itemize}
\item
  Ensure imports do not initialize TRACK execution, read datasets, or
  require optional plotting/training dependencies.
\end{enumerate}

\hypertarget{tests-12}{%
\subsubsection{Tests}\label{tests-12}}

Add tests that exercise the listed imports in isolated subprocesses and
confirm no circular import error. Run focused package, data, simulation,
and ML tests.

\hypertarget{do-not-do-3}{%
\subsubsection{Do Not Do}\label{do-not-do-3}}

\begin{itemize}
\tightlist
\item
  Do not remove top-level imports without a compatibility plan.
\item
  Do not move application behavior into \texttt{\_\_init\_\_.py} files.
\item
  Do not change model or physics behavior.
\end{itemize}

\hypertarget{completion-summary-3}{%
\subsubsection{Completion Summary}\label{completion-summary-3}}

Report the import graph changes, compatibility decisions, import
timings, changed files, tests run, and remaining heavyweight imports.

Source Markdown:
\texttt{docs/codex\_refactor\_tasks/15\_training\_artifact\_path\_contract.md}

\hypertarget{codex-prompt-training-artifact-path-contract}{%
\subsection{Codex Prompt: Training Artifact Path
Contract}\label{codex-prompt-training-artifact-path-contract}}

Read \texttt{AGENTS.md}, \texttt{ml/run\_artifacts.py}, Task 4 and Task
5 trainers, staged training, model comparison, publication metrics, and
relevant tests.

\hypertarget{task-14}{%
\subsubsection{Task}\label{task-14}}

Centralize every per-run artifact path under one versioned artifact
contract.

\hypertarget{requirements-14}{%
\subsubsection{Requirements}\label{requirements-14}}

\begin{enumerate}
\def\labelenumi{\arabic{enumi}.}
\tightlist
\item
  Extend \texttt{RunArtifactPaths} to include standard metrics,
  canonical state-physics metrics, compatibility physics metrics, plots,
  summaries, normalizer manifest, checkpoints, config, and timing
  artifacts.
\item
  Replace manually assembled artifact filenames in training, staged
  training, comparison, and publication modules with this contract.
\item
  Add strict discovery/loading helpers that report missing,
  incompatible, or legacy artifacts explicitly rather than silently
  substituting empty data.
\item
  Preserve current filenames for existing training outputs unless a
  versioned migration adapter is provided.
\item
  Make emitted manifests include artifact paths and artifact-format
  version.
\item
  Keep plain MLP and physics-regularized MLP output trees comparable.
\end{enumerate}

\hypertarget{tests-13}{%
\subsubsection{Tests}\label{tests-13}}

Add tests for:

\begin{enumerate}
\def\labelenumi{\arabic{enumi}.}
\tightlist
\item
  paths for both model types and all splits;
\item
  strict missing-artifact diagnostics;
\item
  backward-compatible loading of current run directories;
\item
  staged-training and publication-table use of the centralized paths.
\end{enumerate}

\hypertarget{do-not-do-4}{%
\subsubsection{Do Not Do}\label{do-not-do-4}}

\begin{itemize}
\tightlist
\item
  Do not retrain models.
\item
  Do not rename or delete existing model artifacts in place.
\item
  Do not alter reported numerical metrics.
\end{itemize}

\hypertarget{completion-summary-4}{%
\subsubsection{Completion Summary}\label{completion-summary-4}}

Report the artifact manifest format, compatibility behavior, changed
files, tests run, and any artifacts intentionally left legacy-only.

Source Markdown:
\texttt{docs/codex\_refactor\_tasks/16\_versioned\_hdf5\_dataset\_compatibility.md}

\hypertarget{codex-prompt-versioned-hdf5-dataset-compatibility-layer}{%
\subsection{Codex Prompt: Versioned HDF5 Dataset Compatibility
Layer}\label{codex-prompt-versioned-hdf5-dataset-compatibility-layer}}

Read \texttt{AGENTS.md}, dataset generation modules,
\texttt{data/batch\_io.py}, dataset schemas, input/target view helpers,
and tests for LEBT and MEBT data loading.

\hypertarget{task-15}{%
\subsubsection{Task}\label{task-15}}

Create a versioned dataset-reader compatibility layer for the legacy
LEBT generator layout, no-offset Sobol batches, and MEBT Sobol batches.

\hypertarget{requirements-15}{%
\subsubsection{Requirements}\label{requirements-15}}

\begin{enumerate}
\def\labelenumi{\arabic{enumi}.}
\tightlist
\item
  Define explicit format identifiers and supported versions for each
  existing HDF5 layout.
\item
  Implement read-only adapters that expose a normalized dataset
  interface: schema identity, inputs, full/fixed parameters when
  available, target groups, units, validity, completion, failure
  categories, and provenance.
\item
  Route dataset loading and audit code through format detection rather
  than file-path-specific assumptions.
\item
  Surface unavailable fields as explicit capability limitations, never
  as fabricated values.
\item
  Keep writers unchanged in this task except for adding unambiguous
  format metadata where it is missing.
\item
  Document a future opt-in migration command, but do not migrate or
  rewrite existing datasets automatically.
\end{enumerate}

\hypertarget{tests-14}{%
\subsubsection{Tests}\label{tests-14}}

Add fixture-based tests for each format covering detection, normalized
reads, unsupported capability reporting, target/input view selection,
and invalid or unknown format errors.

\hypertarget{do-not-do-5}{%
\subsubsection{Do Not Do}\label{do-not-do-5}}

\begin{itemize}
\tightlist
\item
  Do not alter production datasets.
\item
  Do not make legacy files appear to contain state labels or fixed
  inputs they never stored.
\item
  Do not change sampling, TRACK execution, or model architecture.
\end{itemize}

\hypertarget{completion-summary-5}{%
\subsubsection{Completion Summary}\label{completion-summary-5}}

Report supported formats, normalized interface capabilities,
compatibility limitations, changed files, tests run, and the proposed
opt-in migration path.

Source Markdown: \texttt{docs/codex\_refactor\_tasks/README.md}

\hypertarget{refactor-tasks}{%
\subsection{Refactor Tasks}\label{refactor-tasks}}

These prompts capture the physics and code refactors identified in the
2026 repository review. Tasks 01--16 have been executed in numeric
order; the files remain as implementation records and regression-scope
specifications. Future work should preserve the contracts established by
Tasks 11--16 for MEBT batch schema/current provenance, lattice control
manifests, canonical state metrics, package import boundaries, training
artifacts, and HDF5 compatibility.

\begin{longtable}[]{@{}
  >{\raggedright\arraybackslash}p{(\columnwidth - 2\tabcolsep) * \real{0.5000}}
  >{\raggedright\arraybackslash}p{(\columnwidth - 2\tabcolsep) * \real{0.5000}}@{}}
\toprule\noalign{}
\begin{minipage}[b]{\linewidth}\raggedright
Task
\end{minipage} & \begin{minipage}[b]{\linewidth}\raggedright
Scope
\end{minipage} \\
\midrule\noalign{}
\endhead
\bottomrule\noalign{}
\endlastfoot
\texttt{01\_fixed\_parameter\_inference\_contract.md} & Make full 24D
LEBT inference validate fixed values, including fixed current. \\
\texttt{02\_beam\_current\_unit\_contract.md} & Establish and enforce
one explicit current-unit boundary for TRACK. \\
\texttt{03\_mebt\_surrogate\_input\_schema.md} & Define a reusable MEBT
dataset/surrogate schema with injected beam state. \\
\texttt{04\_physics\_training\_output\_scaling.md} & Make output scaling
effective for physics-regularized training. \\
\texttt{05\_unify\_beam\_state\_physics\_api.md} & Consolidate state
moments, Courant-Snyder, and covariance-validity logic. \\
\texttt{06\_legacy\_physics\_cleanup.md} & Repair legacy Twiss/validity
paths and remove avoidable runtime assumptions. \\
\texttt{07\_track\_execution\_orchestration.md} & Unify path-safe
LEBT/MEBT TRACK segment execution and runtime cleanup. \\
\texttt{08\_track\_diagnostic\_output\_contract.md} & Validate
diagnostic output alignment, cardinality, units, and physical
validity. \\
\texttt{09\_production\_batch\_atomicity\_and\_failure\_taxonomy.md} &
Make staged production batches atomic, resumable, and
failure-auditable. \\
\texttt{10\_versioned\_artifact\_and\_config\_loading.md} & Make
versioned artifact/config discovery strict, numeric, and diagnosable. \\
\texttt{11\_mebt\_batch\_schema\_and\_fixed\_current\_contract.md} &
Align MEBT Sobol batches with the schema contract and make fixed-current
provenance enforceable. \\
\texttt{12\_mebt\_lattice\_control\_manifest.md} & Version the MEBT
control order and detect lattice/control semantic drift before
generation. \\
\texttt{13\_unify\_state\_physics\_metric\_reports.md} & Make canonical
component, Courant-Snyder, and covariance metrics common to both MLP
workflows. \\
\texttt{14\_package\_import\_boundary\_cleanup.md} & Reduce eager
package imports and prevent circular, heavyweight import paths. \\
\texttt{15\_training\_artifact\_path\_contract.md} & Centralize training
artifact paths, manifests, discovery, and compatibility loading. \\
\texttt{16\_versioned\_hdf5\_dataset\_compatibility.md} & Add explicit
adapters for legacy, LEBT no-offset, and MEBT HDF5 dataset formats. \\
\end{longtable}

Do not retrain production models as part of these refactors unless a
task explicitly requests it. Preserve dataset compatibility or provide a
documented migration path.

Source Markdown: \texttt{docs/hdf5\_dataset\_compatibility.md}

\hypertarget{hdf5-dataset-compatibility}{%
\subsection{HDF5 Dataset
Compatibility}\label{hdf5-dataset-compatibility}}

\texttt{synapticTrack.data.dataset\_compat} provides read-only detection
and normalized views for these supported layouts:

\begin{itemize}
\tightlist
\item
  \texttt{legacy\_lebt\_dataset\_v2}
\item
  \texttt{lebt\_no\_offset\_sobol\_batch\_v1}
\item
  \texttt{lebt\_no\_offset\_sobol\_consolidated\_v1}
\item
  \texttt{mebt\_sobol\_schema\_batch\_v2}
\item
  \texttt{mebt\_sobol\_schema\_consolidated\_v2}
\item
  \texttt{mebt\_sobol\_legacy\_batch\_v1}
\end{itemize}

The normalized view reports capabilities for named inputs, full and
fixed parameters, beam-state targets, completion flags, and structured
failure categories. A missing capability is reported explicitly;
compatibility readers do not invent absent fixed parameters, labels, or
state targets.

\hypertarget{future-migration}{%
\subsubsection{Future Migration}\label{future-migration}}

No dataset is migrated automatically. A future opt-in CLI is intended to
use a command shaped as:

\begin{Shaded}
\begin{Highlighting}[]
\ExtensionTok{synapticTrack}\NormalTok{ dataset migrate }\AttributeTok{{-}{-}input}\NormalTok{ legacy.h5 }\AttributeTok{{-}{-}output}\NormalTok{ migrated.h5 }\AttributeTok{{-}{-}target{-}format} \OperatorTok{\textless{}}\NormalTok{format{-}id}\OperatorTok{\textgreater{}}
\end{Highlighting}
\end{Shaded}

The command must write a new file, preserve source provenance and source
hash, and require the user to select the target format. Until that
command exists, use \texttt{read\_hdf5\_dataset\_compat()} for legacy
and current datasets.

Source Markdown: \texttt{docs/mebt\_surrogate\_input\_schema.md}

\hypertarget{mebt-surrogate-input-schema}{%
\subsection{MEBT Surrogate Input
Schema}\label{mebt-surrogate-input-schema}}

\hypertarget{dataset-versions}{%
\subsubsection{Dataset Versions}\label{dataset-versions}}

\texttt{mebt\_production\_scratch07\_sobol\_v1} remains the valid
fixed-distribution MEBT diagnostic dataset. It has 19 sampled MEBT
lattice controls and a fixed beam-current metadata value. Every run
starts from the reference TRACK particle distribution
\texttt{RAON/MEBT\_lattice/scratch.\#07}. It is appropriate for studying
control-to-diagnostic response around that one injected distribution,
but it is not the general MEBT surrogate contract.

\texttt{mebt\_general\_surrogate\_v1} defines the future general
surrogate input:

\begin{Shaded}
\begin{Highlighting}[]
\NormalTok{19D MEBT controls + 10D injected Gaussian state {-}\textgreater{} 30D full input}
\NormalTok{                                             + fixed beam\_current}
\end{Highlighting}
\end{Shaded}

The sampled 29D input contains the existing 11 quadrupole gradients,
four cavity fields/phases, and the following injected-state
parameterization:

\begin{Shaded}
\begin{Highlighting}[]
\NormalTok{x0 [mm], xp0 [mrad], y0 [mm], yp0 [mrad],}
\NormalTok{alpha\_x, beta\_x [mm/mrad], emit\_x [mm*mrad],}
\NormalTok{alpha\_y, beta\_y [mm/mrad], emit\_y [mm*mrad]}
\end{Highlighting}
\end{Shaded}

This is an explicit physical equivalent of the established 10D
state-moment convention. It maps to
\texttt{x\_c,\ xp\_c,\ y\_c,\ yp\_c,\ \textless{}x\^{}2\textgreater{},\ \textless{}xx\textquotesingle{}\textgreater{},\ \textless{}x\textquotesingle{}\^{}2\textgreater{},\ \textless{}y\^{}2\textgreater{},\ \textless{}yy\textquotesingle{}\textgreater{},\ \textless{}y\textquotesingle{}\^{}2\textgreater{}}
through the Courant-Snyder covariance. Positive beta and emittance are
validated, and the resulting horizontal and vertical covariance
determinants must be positive.

\hypertarget{fixed-distribution-control-manifest}{%
\subsubsection{Fixed-Distribution Control
Manifest}\label{fixed-distribution-control-manifest}}

The fixed-distribution Sobol workflow uses a versioned
\texttt{mebt\_lattice\_control\_manifest\_v1} derived from
\texttt{RAON/MEBT\_lattice/sclinac.dat}. The manifest fixes the 19D
control order, control type, segment and element location, nominal
value, bounds, and source SHA-256 hash. Generation validates the parsed
source lattice, or optional JSON segment convenience files, against this
manifest before sampling. A mismatch reports the first divergent control
rather than silently applying controls to a changed lattice.

The fixed current is \texttt{0.05} in the configured synapticTrack/TRACK
conversion contract. It is fixed, not zero, and is recorded with its
unit boundary and verification status in batch and consolidated
metadata.

\hypertarget{targets-and-metadata}{%
\subsubsection{Targets and Metadata}\label{targets-and-metadata}}

The general schema reserves these target groups:

\begin{itemize}
\tightlist
\item
  \texttt{diagnostics}: 20D MEBT centroid and rms measurements at
  MWS1--4 and MEBT exit, in mm;
\item
  \texttt{state\_mws4}: 10D centroid/second-moment state at MWS4;
\item
  \texttt{state\_mebt\_exit}: 10D centroid/second-moment state at MEBT
  exit.
\end{itemize}

All general batches must record validity masks and failure reasons.
Current metadata uses canonical synapticTrack current values in
\texttt{uA}, plus the TRACK current conversion boundary and its
verification status. The local repository does not establish TRACK's
physical current unit, so that boundary remains explicitly unverified.

\hypertarget{sampling-and-injection}{%
\subsubsection{Sampling and Injection}\label{sampling-and-injection}}

\texttt{MEBTGeneralSamplingBounds} keeps lattice-control bounds separate
from injected-state bounds. Unit-cube samples are expanded
deterministically, and every expanded injected state is validated before
it can be used. A future production generator must rewrite the
per-sample input distribution and \texttt{track.dat} current before
segmented TRACK execution. \texttt{scratch.\#07} remains the nominal
reference distribution from which those sampling bounds and injection
validation should be calibrated; it is not reused unchanged for all
samples in the general dataset.

No existing fixed-distribution MEBT batch is rewritten. The read-only
HDF5 compatibility layer recognizes both current schema batches and
legacy MEBT fixed-distribution batches. Legacy files expose only fields
physically stored on disk; they do not acquire reconstructed
fixed-current/full-input columns, control labels, state targets, or
completion flags.

Source Markdown: \texttt{docs/package\_import\_boundaries.md}

\hypertarget{package-import-boundaries}{%
\subsection{Package Import Boundaries}\label{package-import-boundaries}}

Package initializers keep lightweight stable types eagerly available and
retain historical workflow exports through lazy compatibility
resolution. This avoids loading analysis, TRACK execution, data
generation, visualization, and ML training stacks merely because a
low-level module is imported.

For new code, import workflow APIs from their owning submodule, for
example:

\begin{Shaded}
\begin{Highlighting}[]
\ImportTok{from}\NormalTok{ synapticTrack.simulation.no\_offset\_sobol\_production }\ImportTok{import}\NormalTok{ run\_generation}
\ImportTok{from}\NormalTok{ synapticTrack.ml.metrics }\ImportTok{import}\NormalTok{ compute\_state\_physics\_metrics}
\ImportTok{from}\NormalTok{ synapticTrack.data.dataset\_compat }\ImportTok{import}\NormalTok{ read\_hdf5\_dataset\_compat}
\end{Highlighting}
\end{Shaded}

Top-level and package-level historical imports remain supported where
possible, but they may load the relevant workflow on first attribute
access. The isolated import smoke tests enforce a five-second budget and
check that low-level imports do not initialize TRACK runtime, complete
ML training modules, or unrelated simulation packages.

Source Markdown: \texttt{docs/physics\_regularized\_output\_scaling.md}

\hypertarget{physics-regularized-output-scaling}{%
\subsection{Physics-Regularized Output
Scaling}\label{physics-regularized-output-scaling}}

Physics-regularized training keeps the model's structured output in
physical coordinates. This preserves the positive-definite state head
and makes existing checkpoints and
\texttt{LEBTSurrogateModel.predict\_physical()} backward compatible.

The train-split output normalizer is nevertheless operational. The data
loss is computed from normalized prediction and target residuals, using
one reusable per-component \texttt{DifferentiableOutputTransform}. This
balances diagnostics and state components with disparate numerical
scales. The transform is fitted only on the training split and recorded
in the run normalizer manifest and checkpoint metadata.

Courant-Snyder target loss, covariance positive-definiteness penalty,
diagnostic rms-positivity penalty, reported metrics, and prediction APIs
use physical units. Therefore their loss weights retain physical
interpretation; output normalization does not silently rescale the
beam-physics constraints.

The transform version is \texttt{train\_split\_output\_transform\_v1}.
Existing production checkpoints without this manifest field remain
physical-output checkpoints and continue to load unchanged. No existing
model is reinterpreted or retrained.

Source Markdown: \texttt{docs/training\_artifact\_contract.md}

\hypertarget{training-artifact-contract}{%
\subsection{Training Artifact
Contract}\label{training-artifact-contract}}

Each plain MLP and physics-regularized MLP run uses the
\texttt{training\_run\_artifacts\_v2} contract from
\texttt{synapticTrack.ml.run\_artifacts}. The contract preserves
established filenames and adds \texttt{artifact\_manifest.json} for
explicit discovery.

\hypertarget{standard-paths}{%
\subsubsection{Standard Paths}\label{standard-paths}}

The manifest records checkpoints, configuration, normalizer manifest,
timing, training history, plots, and per-split artifacts for
\texttt{train}, \texttt{val}, and \texttt{test}:

\begin{itemize}
\tightlist
\item
  \texttt{metrics\_\textless{}split\textgreater{}.json} for standard
  model metrics;
\item
  \texttt{state\_physics\_metrics\_\textless{}split\textgreater{}.json}
  for canonical state-physics reports;
\item
  \texttt{physics\_metrics\_\textless{}split\textgreater{}.json} for
  historical compatibility reports;
\item
  \texttt{derived\_physics\_metrics\_\textless{}split\textgreater{}.json}
  and
  \texttt{covariance\_validity\_metrics\_\textless{}split\textgreater{}.json}
  for focused table consumers.
\end{itemize}

\texttt{discover\_run\_artifacts()} distinguishes a valid current
manifest, a supported legacy run directory with current filenames but no
manifest, and invalid or unsupported manifests.
\texttt{load\_run\_json\_artifact()} returns structured status for
missing or malformed files rather than treating missing metrics as valid
empty data.

New training code must use \texttt{run\_artifacts(run\_dir)}. Existing
model directories are not renamed or migrated in place.

\hypertarget{references}{%
\section*{References}\label{references}}
\addcontentsline{toc}{section}{References}

\hypertarget{refs}{}
\begin{CSLReferences}{1}{0}
\leavevmode\vadjust pre{\hypertarget{ref-courant1958theory}{}}%
Courant, E. D., and H. S. Snyder. 1958. {``Theory of the
Alternating-Gradient Synchrotron.''} \emph{Annals of Physics} 3 (1):
1--48. \url{https://doi.org/10.1016/0003-4916(58)90012-5}.

\leavevmode\vadjust pre{\hypertarget{ref-harris2020array}{}}%
Harris, Charles R., K. Jarrod Millman, Stéfan J. van der Walt, Ralf
Gommers, Pauli Virtanen, David Cournapeau, Eric Wieser, et al. 2020.
{``Array Programming with NumPy.''} \emph{Nature} 585: 357--62.
\url{https://doi.org/10.1038/s41586-020-2649-2}.

\leavevmode\vadjust pre{\hypertarget{ref-ostroumov2006track}{}}%
Ostroumov, P. N., and V. N. Aseev. 2006. {``TRACK: The New Beam Dynamics
Code.''} In \emph{Proceedings of the 2006 International Linear
Accelerator Conference}.

\leavevmode\vadjust pre{\hypertarget{ref-paszke2019pytorch}{}}%
Paszke, Adam, Sam Gross, Francisco Massa, Adam Lerer, James Bradbury,
Gregory Chanan, Trevor Killeen, et al. 2019. {``PyTorch: An Imperative
Style, High-Performance Deep Learning Library.''} In \emph{Advances in
Neural Information Processing Systems}. Vol. 32.

\leavevmode\vadjust pre{\hypertarget{ref-pedregosa2011scikit}{}}%
Pedregosa, Fabian, Gaël Varoquaux, Alexandre Gramfort, Vincent Michel,
Bertrand Thirion, Olivier Grisel, Mathieu Blondel, et al. 2011.
{``Scikit-Learn: Machine Learning in Python.''} \emph{Journal of Machine
Learning Research} 12: 2825--30.

\leavevmode\vadjust pre{\hypertarget{ref-raffin2021stable}{}}%
Raffin, Antonin, Ashley Hill, Adam Gleave, Anssi Kanervisto, Maximilian
Ernestus, and Noah Dormann. 2021. {``Stable-Baselines3: Reliable
Reinforcement Learning Implementations.''} In \emph{Journal of Machine
Learning Research}, 22:1--8. 268.

\leavevmode\vadjust pre{\hypertarget{ref-sobol1967distribution}{}}%
Sobol', I. M. 1967. {``On the Distribution of Points in a Cube and the
Approximate Evaluation of Integrals.''} \emph{USSR Computational
Mathematics and Mathematical Physics} 7 (4): 86--112.
\url{https://doi.org/10.1016/0041-5553(67)90144-9}.

\leavevmode\vadjust pre{\hypertarget{ref-pandas2020}{}}%
team, The pandas development. 2020. {``Pandas-Dev/Pandas: Pandas.''} In
\emph{Zenodo}. \url{https://doi.org/10.5281/zenodo.3509134}.

\leavevmode\vadjust pre{\hypertarget{ref-towers2024gymnasium}{}}%
Towers, Mark, Ariel Kwiatkowski, Jordan Terry, John U. Balis, Gianluca
De Cola, Tristan Deleu, Manuel Goulão, et al. 2024. {``Gymnasium: A
Standard Interface for Reinforcement Learning Environments.''}
\emph{arXiv Preprint arXiv:2407.17032}.

\leavevmode\vadjust pre{\hypertarget{ref-virtanen2020scipy}{}}%
Virtanen, Pauli, Ralf Gommers, Travis E. Oliphant, Matt Haberland, Tyler
Reddy, David Cournapeau, Evgeni Burovski, et al. 2020. {``SciPy 1.0:
Fundamental Algorithms for Scientific Computing in Python.''}
\emph{Nature Methods} 17: 261--72.
\url{https://doi.org/10.1038/s41592-019-0686-2}.

\leavevmode\vadjust pre{\hypertarget{ref-wiedemann2015particle}{}}%
Wiedemann, Helmut. 2015. \emph{Particle Accelerator Physics}. 4th ed.
Springer. \url{https://doi.org/10.1007/978-3-319-18317-6}.

\end{CSLReferences}

\end{document}
